
\medterm{Iatrogenic} Describing an illness, injury, or other harm caused unintentionally by medical examination, treatment, a medicine, a procedure, or a healthcare device. Examples include side effects, procedure-related bleeding, infection, nerve damage, or radiation injury. The term does not imply negligence.

\medterm{Iatrogenic Disease} A disease or injury caused unintentionally by a medical examination, treatment, medication, procedure, or other healthcare intervention.

\medterm{Iatrogenic Injury} Harm caused unintentionally by a medical examination, treatment, procedure, medicine, or healthcare device. It may be temporary or permanent and can include infection, bleeding, medication effects, nerve damage, or injury during surgery.

\medterm{Ibandronic Acid} A bisphosphonate medicine used in some settings to reduce bone loss and complications from cancer affecting bone. It can rarely damage the jaw or cause unusual thigh fractures.
\medterm{IBD} Inflammatory bowel disease, including Crohn's disease and ulcerative colitis, causing long-term digestive-tract inflammation, diarrhea, abdominal pain, bleeding, weight loss, or fatigue.

\synonyms
Inflammatory bowel disease



\medterm{Ibogaine} Ibogaine is a psychoactive plant alkaloid studied for addiction treatment but associated with dangerous cardiac arrhythmias, seizures, and death.

\medterm{Ibritumomab Tiuxetan} An antibody linked to a radioactive substance and used in selected B-cell lymphomas. The antibody attaches to CD20-positive cells and delivers radiation; low blood counts, infection, and delayed marrow problems are important risks.

\medterm{IBS} Alternate index label; see \textit{Irritable bowel syndrome}.

\medterm{IBS-D Irritable Bowel Syndrome With Diarrhea} IBS with diarrhea causes recurrent abdominal pain associated with loose stools without structural disease explaining the symptoms. Management may include diet changes, soluble fiber, gut-directed medicines, behavioral therapy, and evaluation for alarm features or alternative diagnoses.

\synonyms
Diarrhea IBS (IBS-D Irritable Bowel Syndrome With Diarrhea)

\medterm{Ibuprofen} A medicine that reduces pain, fever, and inflammation. It can irritate or bleed the stomach, harm the kidneys, cause fluid retention, or increase heart risks in some people.

\medterm{Ibuprofen Dosing For Children} Children’s ibuprofen dosing is based mainly on weight, age, product concentration, and medical advice; the label dose should not be exceeded and it should be avoided in dehydration or certain kidney conditions.

\medterm{Ibuprofen Overdose} Toxic exposure to more ibuprofen than recommended, which can cause nausea, vomiting, abdominal pain, gastrointestinal bleeding, kidney injury, metabolic acidosis, drowsiness, seizures, or coma.

\medterm{ICD} An implantable cardioverter-defibrillator, a device that monitors heart rhythm and can deliver pacing or an electrical shock to stop dangerous rhythms and prevent sudden cardiac death.

\synonyms
Implantable cardioverter defibrillator

\medterm{ICD Shock Therapy} Treatment delivered by an implanted defibrillator to stop dangerous heart rhythms. The device may use painless pacing or an electrical shock, depending on the rhythm.

\medterm{ICD-O Neoplasm Morphology Classification} A coding system used by cancer registries and pathology services to record what type of tumor is present and whether it is benign or cancerous. It does not replace the diagnosis, location, or stage of the cancer.
\medterm{Ice Pack} A device that applies cold to a body area to help relieve pain or swelling. It should be wrapped to protect the skin and used for limited periods.

\medterm{Ice Pick Headache} A sudden, severe, stabbing head pain that usually lasts seconds. Attacks may recur in different head areas without warning; new, persistent, or unusually severe pain needs medical assessment to exclude serious causes.

\synonyms
Primary stabbing headache

\medterm{ICF Framework} A World Health Organization framework for describing how a health condition affects the body, daily activities, participation in life, and the person's environment.

\medterm{Ichthyosis} A group of conditions that cause very dry, rough, scaly skin. Some forms are inherited and may be present from birth; others are linked to allergies or another illness.

\medterm{Ichthyosis Vulgaris} A common inherited skin condition causing dry, fine scaling, especially on the arms and legs. It often improves with age and can be managed with regular moisturizers and prescribed scale-reducing treatments.

\medterm{Icteric} Describing yellowing of the skin, whites of the eyes, or other tissues caused by too much bilirubin in the blood.

\medterm{Icteric Index} The icteric index is a laboratory estimate of serum yellow discoloration related mainly to bilirubin and can signal jaundice or interfere with testing.

\medterm{Icterus} Yellowing of the skin or whites of the eyes caused by excess bilirubin. It may result from liver disease, blocked bile flow, or rapid breakdown of red blood cells and needs evaluation.

\medterm{Ictus} A sudden neurological event, such as a stroke or seizure.

\medterm{ICU Psychosis} A colloquial term for delirium occurring in an intensive-care setting, with acute changes in attention, awareness, and
cognition. Causes include severe illness, infection, medicines, sleep disruption, pain, and metabolic abnormalities.
It requires prompt evaluation and treatment of reversible causes, with supportive measures to orient and calm the patient.

\synonyms
Psychosis ICU (ICU Psychosis)

\medterm{Idaiovirus} Likely a misspelling of Idaeovirus, a genus of plant-infecting RNA viruses that includes viruses associated with raspberry and privet; it is not a recognized human pathogen.

\medterm{Idarubicin} Idarubicin is an anthracycline chemotherapy drug that intercalates DNA and inhibits topoisomerase II, used mainly in acute myeloid leukaemia; myelosuppression and cumulative cardiotoxicity are important risks.

\medterm{IDC} Alternate index label; see \textit{Invasive ductal carcinoma}.

\medterm{Idea Of Reference} A belief that ordinary events, comments, or media messages have a special personal meaning, often seen in psychotic or mood disorders when held with strong conviction.

\medterm{Ideas Of Reference} The belief that ordinary events, comments, media messages, or other people’s actions have a special personal meaning. If the belief is held with strong conviction despite evidence, it may occur in a psychotic or mood disorder.

\medterm{Ideation} The formation or content of thoughts, especially thoughts related to a particular behavior or symptom. Clinical examples include suicidal ideation and homicidal ideation, which require assessment of intent, plan, means, and immediate safety.

\medterm{Identical Twins} Two individuals formed when one fertilized egg separates into two embryos. They usually have nearly identical DNA, but differences can develop through mutations, epigenetic changes, environment, and chance.

\medterm{Idiopathic} Describing a condition for which no specific cause has been identified after appropriate evaluation. It does not mean that the condition is imaginary, unimportant, or untreatable; it means that its cause is currently unknown.
\medterm{Idiopathic Facial Paralysis (Facial Nerve Problems)} Alternate index label for Facial Nerve Problems.

\medterm{Idiopathic Hypercalciuria} Excess calcium in the urine without an identified cause. It can increase the risk of kidney stones and, less often, bone loss.

\medterm{Idiopathic Hypersomnia} A sleep disorder that causes extreme daytime sleepiness, long or unrefreshing sleep, and difficulty waking up, even when a person has enough time to sleep. Other causes must be ruled out.

\medterm{Idiopathic Interstitial Pneumonias} A group of interstitial lung diseases with no identified cause, characterized by varying degrees of lung inflammation and fibrosis; diagnosis requires clinical, radiologic, and sometimes pathologic assessment.

\medterm{Idiopathic Intracranial Hypertension} Increased pressure around the brain without a tumor or other clear cause. It can cause headaches, a pulsing sound in the ears, and vision problems and may damage sight if untreated.

\medterm{Idiopathic Myelofibrosis} Alternate index label; see \textit{Myelofibrosis}.

\medterm{Idiopathic Pulmonary Fibrosis} A long-term lung disease of unknown cause in which scar tissue gradually builds up in the lungs. It causes increasing shortness of breath and a dry cough and requires specialist care.

\medterm{Idiopathic Thrombocytopenic Purpura} An older name for immune thrombocytopenia, a disorder in which the immune system destroys platelets and may cause easy bruising, tiny skin spots, nosebleeds, heavy menstrual bleeding, or other bleeding.
\medterm{Idiopathic Toe Walking} Alternate index label; see \textit{Toe walking in children}.

\medterm{Idiosyncrasy} An unusual individual reaction to a medicine or exposure that is not explained by the usual pharmacologic effect.

\medterm{Idiosyncratic Reaction} An uncommon, unexpected, and often genetically determined adverse drug reaction that is unrelated to the drug's known pharmacology and usually dose-independent (a Type B reaction).
\medterm{Idiotypic} Relating to the unique antigen-binding features of an antibody or immune-cell receptor. These features can be recognized by other antibodies and are useful in studying immune responses and some blood cancers.

\medterm{Idioventricular Rhythm} A slow heart rhythm that starts in the lower chambers of the heart instead of the normal pacemaker. It may occur when the usual electrical signal is blocked or when the heart is severely ill.

\medterm{Idoxuridine} Idoxuridine is a thymidine analogue antiviral formerly used mainly topically for herpes simplex keratitis; it inhibits viral DNA synthesis but can irritate or damage the corneal epithelium.

\medterm{Ifosfamide} Ifosfamide is an alkylating chemotherapy drug that forms DNA cross-links and is used for selected sarcomas and other cancers; haemorrhagic cystitis, encephalopathy, infertility, and myelosuppression require monitoring and mesna protection.

\medterm{iFR} Alternate name; see \textit{Instantaneous Wave-Free Ratio}.

\medterm{Ig} Abbreviation for immunoglobulin, an antibody protein produced by B cells and plasma cells; major classes include IgG, IgA, IgM, IgE, and IgD.

\medterm{IgA} A type of antibody found mainly in body fluids and linings, including saliva, tears, breast milk, and the digestive and breathing passages. It helps protect these surfaces from infection.

\medterm{IgA Nephropathy} A kidney disease in which IgA antibodies build up in the kidney's filtering units. It can cause blood or protein in the urine, high blood pressure, and sometimes kidney failure.

\medterm{IgA Nephropathy (Berger Disease)} A kidney disease caused by deposits of immunoglobulin A in the kidney filters, potentially leading to blood in urine, protein loss, high blood pressure, or kidney failure.

\synonyms
Berger's Disease

\medterm{IgA Vasculitis} An illness in which inflammation damages small blood vessels, often causing a raised purple rash, joint pain, stomach pain, and kidney problems. It is most common in children and may follow an infection.

\synonyms
Anaphylactoid Purpura

\medterm{IgA Vasculitis (Henoch-Schonlein Purpura)} A disease in which the immune system inflames small blood vessels, usually after an infection. It can cause purple spots on the legs, joint pain, abdominal pain, and kidney problems, especially in children.

\synonyms
Henoch-Schonlein Purpura

\medterm{IgD} A less common type of antibody found mainly on the surface of certain immune cells. It helps those cells recognize signals, but its full role in human disease is not yet clear.

\medterm{IgE} A type of antibody involved in allergies and defense against some parasites. When it reacts to an allergen, it can trigger symptoms such as sneezing, hives, asthma, or anaphylaxis.

\medterm{IgG} The most common type of antibody in the blood. It helps neutralize germs and toxins, supports immune protection after infection or vaccination, and can cross the placenta to protect a developing baby.

\medterm{IgG And Complement Deposition At The Dermis-Epidermis Junction} An immunofluorescence finding that may occur in the lupus band test, in which immune
complexes containing immunoglobulin and complement deposit along the dermoepidermal
junction. The pattern and clinical context are needed for interpretation.

\medterm{IgG4-Related Disease} An uncommon immune disorder that causes inflammation and scar-like growths in one or more organs. It may affect the pancreas, salivary glands, eyes, kidneys, lungs, or blood vessels; diagnosis usually requires imaging, blood tests, and sometimes a biopsy.

\medterm{IgG4-Related Disease Arteritis} Inflammation of a large or medium-sized artery as part of IgG4-related disease. It may thicken or weaken the artery and can cause narrowing, an aneurysm, pain, or reduced blood flow.

\medterm{IgG4-Related Ophthalmic Disease} An immune-mediated fibroinflammatory disorder that can affect the lacrimal glands, orbit, eyelids, or
extraocular muscles. It may cause swelling, mass lesions, double vision, or visual
symptoms and can occur with IgG4-related disease in other organs.

\medterm{Igh Gene} A gene region that helps make the heavy part of an antibody. Immune cells rearrange these gene segments so the body can produce antibodies that recognize many different targets.
\medterm{IgM} A large type of antibody that is often made early during an infection. It helps clump germs and activate other parts of the immune system; a blood test may suggest recent exposure but must be interpreted with other results.

\medterm{IIH} Alternate index label; see \textit{Pseudotumor cerebri (idiopathic intracranial hypertension)}.

\medterm{Ikappab Kinase} A group of enzymes that helps turn on genes involved in inflammation and immune responses.

\medterm{Ileal Atresia} A condition present at birth in which part of the ileum, the last section of the small intestine, is closed or missing. It causes bowel blockage in a newborn, with a swollen abdomen, green vomiting, and failure to pass the first stool.

\medterm{Ileal Conduit} A surgically created passage that diverts urine through a segment of ileum to an opening in the abdominal wall.

\medterm{Ileal Fistula} An abnormal connection involving the ileum and another organ, bowel segment, or skin surface, which may cause intestinal contents to leak and lead to infection, malnutrition, or fluid loss.

\medterm{Ileal Hemorrhage} Bleeding from the ileum, the distal portion of the small intestine, caused by ulceration, inflammation, vascular lesions, tumors, trauma, or other disease.

\medterm{Ileal Neoplasm} An ileal neoplasm is an abnormal growth in the ileum, the last part of the small intestine; it may be benign or cancerous.

\medterm{Ileal Obstruction} Partial or complete blockage of the ileum, the final portion of the small intestine. Causes include
adhesions, hernia, inflammatory bowel disease, tumors, and twisting of the bowel. Symptoms may
include cramping pain, vomiting, abdominal distension, and inability to pass stool or gas.

\medterm{Ileal Pouch-Anal Anastomosis} An operation that removes the colon and rectum and creates a pouch from the end of the small intestine. The pouch is connected to the anus so bowel movements can continue through the usual route.

\medterm{Ileal Stenosis} Narrowing of the ileum that may impair passage of intestinal contents and result from Crohn disease, inflammation, scarring, ischemia, tumor, or surgery.

\medterm{Ileal Ulcer} An ileal ulcer is a break in the lining of the ileum caused by inflammation, infection, medicines, ischemia, trauma, or another intestinal disorder.

\medterm{Ileitis} Inflammation of the ileum, the distal segment of the small intestine. Causes include
Crohn disease, infection, ischemia, medication-related injury, and other inflammatory
disorders; symptoms may include right-lower-quadrant pain, diarrhea, fever, or bleeding.

\medterm{Ileitis (Crohn'S Disease)} Alternate index label for Crohn's Disease.

\medterm{Ileitis, Crohn} Crohn disease involving the ileum, the final part of the small intestine, causing chronic or relapsing inflammation that may produce abdominal pain, diarrhea, weight loss, strictures, or fistulas.

\medterm{Ileocaecal} Ileocaecal means relating to the junction between the ileum and the caecum, where the small intestine joins the large intestine.

\medterm{Ileocecal Valve} A sphincter-like fold at the junction of the terminal ileum and caecum. It regulates
passage of intestinal contents into the colon and helps limit reflux of colonic contents
into the small intestine.

\medterm{Ileocolitis, Crohn} Crohn disease affecting both the ileum and colon, causing patchy chronic inflammation that may lead to abdominal pain, diarrhea, weight loss, strictures, fistulas, or perianal disease.

\medterm{Ileostomy} An operation that brings the end of the small intestine through an opening in the abdomen so stool drains into an external pouch. Output is often liquid, so adequate fluids, salt replacement, and skin care are important.







\medterm{Ileum} The final and longest part of the small intestine. It absorbs vitamin B12, bile salts, and remaining nutrients before food passes into the large intestine.

\medterm{Ileus} Failure of coordinated intestinal propulsion without a mechanical blockage, usually from
postoperative bowel inhibition, inflammation, metabolic disturbance, or medication
effect. It causes distension, nausea, vomiting, and reduced passage of flatus or stool.

\medterm{Iliac} Relating to the ilium or the iliac region, including the iliac bones and iliac blood vessels.

\medterm{Iliac Artery} One of the large arteries that carries blood to the pelvis and lower limb. The common iliac arteries divide into internal branches for the pelvic organs and external branches that continue into the legs.

\medterm{Iliac Crest} The iliac crest is the curved upper border of the ilium and a landmark and attachment site for abdominal, back, and pelvic muscles.

\medterm{Iliac Vein} An iliac vein drains blood from the pelvis and lower limb toward the inferior vena cava; thrombosis can cause swelling and pulmonary embolism.

\medterm{Common Iliac Vein} A vein formed by the union of the internal and external iliac veins. The right and left common iliac veins unite to form the inferior vena cava.

\medterm{External Iliac Vein} A continuation of the femoral vein above the inguinal ligament that joins the internal iliac vein to form the common iliac vein.

\medterm{Internal Iliac Vein} A vein draining the pelvic organs, pelvic wall, and gluteal region that joins the external iliac vein to form the common iliac vein.

\medterm{Iliopsoas Muscles} The psoas major and iliacus muscles working together to bend the hip. They help lift the thigh, stabilize the spine, and maintain posture, and may cause groin or back pain when strained.

\synonyms
Iliopsoas

\medterm{Iliotibial Band} A strong band of tissue running along the outer side of the thigh from the hip to the knee. It helps stabilize the hip and knee during movement.

\medterm{Iliotibial Band Syndrome} Pain on the outer side of the knee caused by irritation of the iliotibial band, especially during repeated activities such as running or cycling.

\synonyms
ITBS (Iliotibial Band Syndrome)

\medterm{Ilium} Largest hip bone forming the superior portion. Body contributes to acetabulum; wing has
a concave iliac fossa. Iliac crest runs ASIS to PSIS, a palpable landmark for lumbar
puncture level. Sacroiliac joint articulates with sacrum. Gluteal muscles originate from
its external surface; iliacus fills its fossa.

\medterm{Illness Anxiety Disorder} Preoccupation with having or acquiring a serious illness, with minimal or no somatic
symptoms present. The individual has a high level of anxiety about health, performs
excessive health-related behaviors, and experiences maladaptive avoidance. Symptoms
persist for at least six months. The individual is not delusional and insight may vary.

\medterm{Illusion} A misperception or misinterpretation of a real external stimulus, such as mistaking a shadow for a person. Illusions may occur with visual impairment, delirium, seizures, intoxication, or severe psychiatric illness and differ from hallucinations, which occur without an external stimulus.

\medterm{Iloprost} A prostacyclin analogue that causes vasodilation and inhibits platelet aggregation; inhaled or infused formulations are used for selected forms of pulmonary hypertension and severe peripheral ischemia.

\medterm{Image View} An image view is a particular angle or projection used when producing a medical image, such as a radiograph.

\medterm{Image-Guided Ablation} Destruction of abnormal tissue using an energy source such as heat, cold, or chemicals while imaging guides placement and
monitors the procedure.

\medterm{Image-Guided Biopsy} Removal of a tissue sample while ultrasound, CT, MRI, mammography, or fluoroscopy guides the needle or instrument to the target and helps avoid critical structures.

\medterm{Image-Guided Drainage} Placement of a catheter or needle under imaging guidance to evacuate fluid, pus, blood,
or bile from a pathologic collection. Indications include abscess, biloma, urinoma,
pleural collection, and symptomatic cysts.

\medterm{Image-Guided Therapy} A minimally invasive procedure performed with imaging to guide instruments or treatment to a target tissue.

\medterm{IGRT} Image-guided radiation therapy uses medical imaging, such as CT, MRI, ultrasound, or X-ray, to position the patient and target radiation accurately before or during treatment. It helps deliver radiation to a tumor while reducing exposure to nearby healthy tissue.

\medterm{Imaging} The production of pictures of internal body structures using methods such as radiography, ultrasound,
computed tomography, magnetic resonance imaging, or nuclear medicine.

\medterm{Imaging And Radiology} Imaging and radiology use methods such as X-ray, ultrasound, CT, and MRI to create pictures that help diagnose or guide treatment of disease.

\medterm{Imaging Technique} An imaging technique is a method for producing pictures of body structures or processes, such as ultrasound, CT, MRI, or nuclear imaging.

\medterm{Imaging-Guided Biopsy} Alternate index label; see \textit{Image-guided biopsy}. Appropriate specimen handling, bleeding-risk assessment, informed consent, and post-procedure monitoring remain essential.

\medterm{Imatinib Mesylate} The mesylate salt of imatinib, a tyrosine-kinase inhibitor used to treat BCR-ABL-positive chronic myeloid leukemia, gastrointestinal stromal tumors with KIT alterations, and other selected cancers.

\medterm{Imerslund-GräSbeck Syndrome} A rare inherited disorder in which the small intestine cannot absorb vitamin B12 properly. It can cause low B12 levels, anemia, tiredness, numbness, balance problems, or other nerve damage despite adequate dietary intake.

\medterm{Imidazoles} A group of antifungal medicines, including clotrimazole, miconazole, and ketoconazole, used to treat fungal infections of the skin, mouth, or other body areas.

\medterm{Imide} An imide is a chemical compound containing nitrogen bonded to two carbonyl groups and may have pharmaceutical, metabolic, or industrial significance.

\medterm{Imines} Imines are compounds containing a carbon-nitrogen double bond and occur in chemical synthesis, metabolism, and drug chemistry.

\medterm{Imino Acid} A compound related to amino acids whose nitrogen is part of a ring or imine structure. Proline is commonly grouped with imino acids in biochemistry.
\medterm{Iminoacetic Acid Scan} A nuclear medicine scan that uses a small amount of radioactive medicine to show how bile moves through the liver, gallbladder, and bile ducts. It is also called a HIDA scan.

\medterm{Imipenem} A strong antibiotic, usually given with cilastatin, used for serious bacterial infections. It can cause allergic reactions, stomach problems, seizures, or antibiotic resistance.

\medterm{Imipramine} A medicine used for depression and some bladder problems. It can cause dry mouth, constipation, dizziness, abnormal heart rhythms, and dangerous effects in overdose.

\medterm{Imipramine Overdose} Toxic ingestion of imipramine, a tricyclic antidepressant, which can cause anticholinergic effects, seizures, hypotension, ventricular arrhythmias, coma, and life-threatening cardiac toxicity.

\medterm{Imiquimod} Imiquimod is a topical immune-response modifier that activates Toll-like receptor 7 and induces cytokines; it treats selected actinic keratoses, superficial basal-cell carcinomas, and external genital warts, often causing local inflammation.

\medterm{Immature} Immature means not fully developed or differentiated, such as an immature cell, organ, tissue, or developmental stage.

\medterm{Immature Ovarian Teratoma} A malignant germ-cell tumor containing immature embryonal tissues, usually diagnosed in adolescents or young adults and graded by the amount of immature neuroepithelium.

\medterm{Immaturity} Development that is less physically, neurologically, emotionally, or socially advanced than expected for a person's age or stage of life.

\medterm{Immediate Hypersensitivity} A rapid allergic reaction, usually involving IgE antibodies and mast cells. It can cause hives, itching, swelling, wheezing, vomiting, or anaphylaxis within minutes of exposure.

\medterm{Immersion Foot (Trench Foot)} A painful cold injury caused by prolonged exposure of the feet to cold, wet conditions without freezing, leading to numbness, swelling, and skin changes.

\synonyms
Trench Foot

\medterm{Immobilization} Restriction of movement of a body part or the whole patient, using rest, splints, casts, braces, traction, or other methods. It may protect an injury but can cause stiffness, pressure injury, muscle loss, venous thromboembolism, and other complications if prolonged.

\medterm{Immobilize} To restrict movement of a body part, often with a splint or cast, to support healing.

\synonyms
Immobilization

\medterm{Immune} Having an immune response that recognizes and helps protect against a particular antigen, infection, or toxin; immunity may be innate or acquired.

\medterm{Immune And Infectious Diseases} A medical field encompassing disorders of the immune system and illnesses caused by
pathogenic microorganisms. Immune system disorders include autoimmune diseases (e.g.,
lupus, rheumatoid arthritis), immunodeficiencies, and hypersensitivity reactions.

\medterm{Immune Checkpoint} A control signal that helps prevent immune cells from becoming too active or damaging healthy tissue. Some cancer medicines block these signals to help immune cells attack tumors, but this can cause inflammation in normal organs.

\medterm{Immune Complex} A cluster formed when an antibody attaches to a foreign or abnormal substance. If these clusters are not cleared, they can collect in tissues and cause inflammation, blood-vessel damage, or kidney disease.

\medterm{Immune Evasion} The ability of a pathogen, tumor, or other abnormal cell to avoid recognition or destruction by the immune system.

\medterm{Immune Hemolytic Anemia} Anemia caused by immune-mediated destruction of red blood cells, either by
autoantibodies or antibodies directed against foreign red-cell antigens. It may cause
fatigue, jaundice, dark urine, splenomegaly, and laboratory evidence of hemolysis.

\medterm{Immunohemolysis} Destruction of red blood cells by an immune reaction involving antibodies, complement, or both. It may occur in autoimmune hemolytic anemia, incompatible transfusion reactions, or other antibody-mediated disorders.

\medterm{Immune Response} The body’s reaction to an infection, foreign substance, or abnormal cell. It includes rapid defenses such as inflammation and slower, targeted responses involving antibodies and immune cells that remember previous exposures.

\medterm{Immune Status} Immune status describes the functional state of a person's immune defenses, including immunocompetence, prior infection or vaccination, immunosuppression, and susceptibility to infection.

\medterm{Immune System} The body's network of cells, tissues, organs, and proteins that recognizes and responds to infections, abnormal cells, and harmful substances. It includes rapid innate defenses and learned adaptive defenses; overactivity can cause allergy or autoimmunity, while weakness increases infection risk.

\medterm{Immune Thrombocytopenia} An autoimmune disorder in which the immune system destroys platelets or reduces their production. Low platelets can cause easy bruising, tiny red or purple skin spots, nosebleeds, heavy periods, or other bleeding; causes may be primary or related to another illness or medicine.

\medterm{Immune Thrombocytopenic Purpura} An older name for immune thrombocytopenia, in which the immune system destroys platelets. It may cause bruising, pinpoint skin spots, nosebleeds, heavy menstrual bleeding, or no symptoms.

\medterm{Immune Thrombocytopenic Purpura (Idiopathic Thrombocytopenic Purpura)} Alternate index label for Idiopathic Thrombocytopenic Purpura.

\medterm{Immune Tolerance} The ability of the immune system to avoid attacking the body’s own tissues while still responding to infections. Failure of tolerance can contribute to autoimmune disease; excessive tolerance can allow some cancers or infections to persist.

\medterm{Immune-Mediated Necrotizing Myopathy} A subacute severe proximal myopathy with very high CK (often >10,000), minimal
inflammation on biopsy, necrotic fibers. Two major autoantibodies: anti-HMGCR
(3-hydroxy-3-methylglutaryl-CoA reductase, statin-associated or spontaneous) and
anti-SRP (signal recognition particle).

\medterm{Immune-Related Adverse Event} An inflammatory or autoimmune toxicity caused by immune checkpoint blockade or another
immune-activating cancer treatment. It may affect skin, bowel, liver, endocrine organs,
lungs, heart, kidneys, or nervous system and often requires treatment-specific
immunosuppression.

\medterm{Immunity} The ability of the body to recognize and resist infectious agents or other foreign substances through innate or adaptive defenses.

\medterm{Immunization} The process of inducing protection against an infection, usually through vaccination. Recommended
vaccines and schedules depend on age, previous doses, pregnancy, immune status, medical
conditions, travel, exposure risk, and current public-health guidance.

\medterm{Immunizations For People With Diabetes} Vaccination is especially important in diabetes because infections may be more severe; recommendations depend on age, health status, vaccines received, and local guidance.

\medterm{Immunoassay} An immunoassay is a laboratory test that uses antibodies to detect or measure a substance such as a hormone, drug, antigen, or antibody.

\medterm{Immunoblast Count} An immunoblast count measures activated lymphocytes, usually in a tissue or blood sample, and may help evaluate immune responses or lymphoid disease.

\medterm{Immunoblotting} Immunoblotting, or Western blotting, separates proteins by electrophoresis and detects a target with specific antibodies; band pattern and intensity are interpreted with controls and clinical context.

\medterm{Immunochemistry} The study or laboratory measurement of biological molecules using specific antigen-antibody reactions, including immunoassays and immunostaining.

\medterm{Immunocompetence} Immunocompetence is the capacity of the immune system to recognize antigens and mount an effective, appropriately regulated response; deficiency increases infection risk and may impair vaccine responses.

\medterm{Immunocompetent} Having an immune system capable of mounting an appropriate response to common infectious or foreign antigens.

\medterm{Immunocompromised} Having weakened or impaired immune defenses because of disease, medication, cancer treatment, transplantation, or another cause, increasing susceptibility to some infections and complications.

\medterm{Immunocompromised Host} A person whose immune defenses are impaired by congenital disease, infection,
malignancy, immunosuppressive treatment, organ failure, or another condition. Such
patients have increased susceptibility to opportunistic infection, severe disease, and
atypical presentations.

\medterm{Immunoconjugate} A therapeutic or diagnostic molecule in which an antibody or other targeting protein is chemically linked to a drug, toxin, enzyme, or imaging agent.

\medterm{Immunocytochemistry} A laboratory method using labeled antibodies to identify specific proteins or antigens within cells or tissue sections.

\medterm{Immunodeficiency} Impaired development or function of one or more components of the immune system,
resulting in increased susceptibility to infection, malignancy, or autoimmunity. It may
be primary and genetic or secondary to disease, medication, malnutrition, or another
acquired cause.

\medterm{Immunodeficiency And Infections} Immunocompromised individuals are at increased risk for opportunistic infections.
Primary immunodeficiencies are genetic disorders affecting immune function. Secondary
immunodeficiencies include HIV/AIDS, organ transplantation, chemotherapy, and
corticosteroid use.

\medterm{Immunodeficiency Diseases} Disorders in which one or more components of the immune system are absent, deficient, or
dysfunctional. They may be inherited or acquired and can cause recurrent, severe,
unusual, or persistent infections, autoimmune disease, and increased malignancy risk.

\medterm{Immunodeficiency Disorders} Conditions in which part of the immune system is missing, weak, or does not work properly. They may be inherited or acquired and can cause repeated, severe, unusual, or persistent infections and sometimes autoimmune disease or cancer.

\medterm{Immunoelectrophoresis} A laboratory method that separates proteins by electrophoresis and identifies them with specific antibodies, often to study immunoglobulins.

\medterm{Immunoelectrophoresis - Urine} A laboratory technique that separates and identifies urinary proteins with specific antibodies, used to investigate abnormal immunoglobulin excretion and monoclonal light chains.

\medterm{Immunoelectrophoresis – Blood} A laboratory technique that separates blood proteins by electrophoresis and identifies them with specific antibodies, used historically to evaluate abnormal immunoglobulins and monoclonal proteins.

\medterm{Immunofixation - Urine} A urine test that uses electrophoresis and antibodies to identify and type monoclonal immunoglobulins or light chains, including Bence Jones proteins.

\medterm{Immunofixation Blood Test} A blood test that identifies and types monoclonal immunoglobulins, helping evaluate monoclonal gammopathy, multiple myeloma, and related plasma-cell disorders.

\medterm{Immunofluorescence} A technique using fluorescently labeled antibodies to detect specific antigens in
tissues, cells, or microorganisms; visualized by fluorescence microscopy.

\medterm{Immunofluorescence Technique} An immunofluorescence technique uses fluorescently labelled antibodies to locate specific proteins, organisms, or immune deposits in a sample.

\medterm{Immunogenetics} The study of how genes influence immune-system development, function, and disease. It includes inherited susceptibility to autoimmune disease, differences in immune responses, and genetic matching for transplantation.

\medterm{Immunogenicity} The ability of an antigen, vaccine, or other substance to provoke an immune response, including antibody or cellular responses.

\medterm{Immunoglobulin} An antibody protein made by immune cells that recognizes specific foreign substances and helps block or remove infections and other harmful targets.

\medterm{Immunoglobulin (Antibody)} A protein made by immune cells that recognizes a particular foreign substance, germ, or abnormal cell. Antibodies can block harmful targets, mark them for removal, or activate other immune defenses.

\synonyms
Antibody

\medterm{Immunoglobulin A} An antibody found mainly in mucus, saliva, tears, breast milk, and intestinal fluid. It protects surfaces exposed to the outside world and helps prevent germs from attaching and entering tissues.

\medterm{Immunoglobulin D} A type of antibody found in very small amounts in the blood and on the surface of some B cells. Its exact role is not fully understood, but it may help regulate immune responses in the airways and other tissues.

\medterm{Immunoglobulin E} An antibody involved in many immediate allergic reactions and defense against some parasites. It activates mast cells, which release chemicals causing itching, swelling, mucus, wheezing, or other allergy symptoms.

\synonyms
IgE

\medterm{Immunoglobulin G} The most common antibody in blood and tissue fluid, helping protect against many infections. It can remain after infection or vaccination, cross the placenta during pregnancy, and sometimes cause harmful immune reactions.

\medterm{Immunoglobulin Idiotype} The distinctive part of an antibody that determines what it recognizes. Different immune-cell families make antibodies with different idiotypes; the concept is mainly used in laboratory research and some cancer treatments.

\medterm{Immunoglobulin Isotype} The class of an antibody, determined by its constant region. The main isotypes are IgG, IgA, IgM, IgD, and IgE; each has different roles and locations in immune defense.

\medterm{Immunoglobulin M} The first major antibody produced during an initial immune response.

\medterm{Immunoglobulin Replacement Therapy} Treatment that supplies pooled antibodies through a vein or under the skin when a person's immune system makes too few effective antibodies. It can reduce recurrent infections in selected inherited or acquired antibody deficiencies; headache, fever, and infusion reactions may occur.

\medterm{Immunohistochemistry} A laboratory method that uses specially marked antibodies to find particular proteins in a tissue sample. It helps identify tumor type, infection, or immune disease and can guide treatment choices.

\medterm{Immunologic Adjuvant} A substance added to some vaccines to strengthen or guide the immune response. It may improve protection but can also increase temporary pain, swelling, fever, or other vaccine reactions.
\medterm{Immunologic Factor} A substance, cell, gene, or process that affects immune-system activity or control. Examples include antibodies, immune cells, signaling proteins, inherited traits, medicines, and conditions that alter defense against disease.

\medterm{Immunologic Memory} The capacity of the adaptive immune system to respond more rapidly and strongly after re-exposure to an antigen because of long-lived memory lymphocytes.

\medterm{Immunologic Skin Test} An immunologic skin test places a small amount of an antigen or allergen in or on the skin to assess immune reactivity, such as allergy or tuberculosis exposure.

\medterm{Immunologic Surveillance} The immune system’s ongoing detection and removal of abnormal, infected, or damaged cells. This protection helps control infections and may recognize early cancerous changes, although it is not always successful.

\medterm{Immunologic Technique} A laboratory or clinical method that uses antibodies, antigens, or immune reactions to detect or measure a substance.

\medterm{Immunologic Test} An immunologic test examines immune cells, antibodies, antigens, or immune-related proteins to help evaluate infection, allergy, autoimmune disease, or immune deficiency.

\medterm{Immunological Memory} The ability of the adaptive immune system to mount a faster, larger, and more effective
response upon re-exposure to an antigen, conferred by long-lived memory B cells and
memory T cells generated during the primary response.

\medterm{Immunological Synapse} The organized contact point where an immune cell meets a target or another immune cell. Signals gathered at this contact help decide whether the immune cell should become active.

\medterm{Immunological Tolerance} The reduced or absent immune response to a specific antigen, including the normal lack of attack on the body’s own tissues.

\medterm{Immunologically Privileged Site} A body area where immune responses are partly limited, helping protect delicate tissues from inflammation while still allowing important defense against infection.

\medterm{Immunology} The medical and biological study of the immune system, including its normal defenses and disorders such as allergy, autoimmunity, immunodeficiency, infection, and transplant rejection.
\medterm{Immunomodulation} Modification of immune activity to enhance, suppress, or redirect an immune response. Immunomodulatory medicines are used in cancer, autoimmune disease, transplantation, and infection but may increase susceptibility to infection or other adverse effects.

\medterm{Immunomodulators In Dermatology} Topical immunomodulators treat inflammatory skin diseases without corticosteroid side
effects. Pimecrolimus 1 percent cream and tacrolimus 0.03 and 0.1 percent ointments are
calcineurin inhibitors used for atopic dermatitis, facial eczema, and intertriginous
areas. They are steroid-sparing agents.

\medterm{Immunophenotyping} Identification and characterization of cells according to their surface, cytoplasmic, or
nuclear antigens, usually by flow cytometry or immunohistochemistry. It is used to
classify leukemias, lymphomas, immune-cell subsets, and some immunodeficiency disorders.

\medterm{Immunoprecipitation} A laboratory technique that uses an antibody to bind and isolate a specific antigen or protein from a mixture. It is used in research and some validated diagnostic or analytical methods.

\medterm{Immunoprotein} A protein involved in immune defense or immune regulation, such as an antibody, complement protein, cytokine, or antigen-presenting molecule. Immunoproteins identify threats, coordinate immune cells, or help remove damaged or infected material.

\medterm{Immunoradiometric Assay} A laboratory test using radioactive labels on antibodies to measure a specific substance, such as a hormone, protein, or other marker.

\medterm{Immunoregulation} Control of the strength, duration, and target of immune responses, including activation, suppression, and tolerance. Abnormal immunoregulation contributes to autoimmunity, allergy, immunodeficiency, and chronic inflammation.

\medterm{Immunosenescence} Age-related changes in innate and adaptive immunity that alter infection susceptibility,
vaccine responses, inflammation, and cancer surveillance. It is influenced by chronic
disease, nutrition, medications, and functional status.

\medterm{Immunosuppression} Reduced immune-system activity caused by disease, medicines, cancer treatment, or another condition. It can prevent rejection of a transplant or control autoimmune disease but increases susceptibility to infections and some cancers.
\medterm{Immunosuppression In Transplant} Medications to prevent rejection. Induction with thymoglobulin or basiliximab.
Maintenance with tacrolimus, mycophenolate, and prednisone. Monitor drug levels and
renal function. Calcineurin inhibitor nephrotoxicity is major concern requiring dose
optimization.

\medterm{Immunosuppression Screening} Pre-treatment evaluation for latent or active infection, vaccination status, malignancy
risk, and organ dysfunction before immunosuppressive therapy. Testing commonly includes
tuberculosis, hepatitis B and C, HIV, blood counts, liver and kidney function, and
pregnancy when relevant.

\medterm{Immunosuppressive} Describing a medicine, treatment, or condition that reduces immune-system activity. Immunosuppression can prevent transplant rejection or control autoimmune disease, but it may increase the risk of infection, some cancers, and reduced vaccine responses.

\medterm{Immunotherapy} Cancer treatment that helps the immune system recognize or attack cancer cells, or supplies immune cells or antibodies that do so. Main approaches include checkpoint medicines, engineered immune cells, antibodies, and immune-stimulating treatments. Benefits and side effects depend on the cancer and treatment.

\medterm{Immunotherapy And Biologic Treatments} Advanced therapies that harness or modify the immune system to treat diseases,
particularly cancer, autoimmune disorders, and inflammatory conditions. Biologics are
large-molecule drugs derived from living organisms that target specific immune pathways,
such as TNF inhibitors, monoclonal antibodies, and checkpoint inhibitors.

\medterm{Immunotherapy For Cancer} Cancer treatment that stimulates, restores, redirects, or supplies immune activity to recognize and attack malignant cells, including immune-checkpoint inhibitors, cellular therapies, and antibodies.

\medterm{Immunotherapy In Dermatology} Cancer immunotherapy uses medicines that alter immune activity against skin cancers.
Checkpoint inhibitors target pathways such as PD-1, PD-L1, or CTLA-4. Their use depends
on the tumor type, stage, prior treatment, and the patient’s health.


\medterm{Immunotoxicity} Harm to immune-system function caused by chemicals, medicines, infections, or other exposures. It may weaken defenses, trigger excessive sensitivity, cause autoimmunity, or change responses to vaccines and infections.

\medterm{Immunotoxin} A targeted toxin made by linking a toxic protein or drug to an antibody or binding molecule that directs it toward selected cells.

\medterm{Immunoturbidimetry} A laboratory method that measures cloudiness produced when antibodies react with a target antigen, allowing estimation of the target’s concentration.

\medterm{Imp Dehydrogenase} IMP dehydrogenase converts inosine monophosphate to xanthosine monophosphate in guanine-nucleotide synthesis and is a target of some immunosuppressive drugs.

\medterm{Impact} The forceful contact or collision of one object or body part with another, or the effect a disease, treatment, or event has on health; context determines the meaning.

\medterm{Impacted} Firmly wedged or unable to move from its normal position, as with an impacted tooth or stool. The result may cause pain, blockage, inflammation, infection, or damage to nearby tissues.
\medterm{Impacted Fracture} A fracture in which one broken bone fragment is driven firmly into another, often producing relative stability but sometimes causing shortening or deformity.

\medterm{Impacted Tooth} Tooth unable to erupt into functional position. Most common: mandibular third molars,
maxillary canines. May cause pathology: cyst, root resorption, caries. Treatment:
monitoring or extraction.

\medterm{Impacted Wisdom Teeth} Wisdom teeth that cannot emerge normally because they are blocked by bone, gum, or another tooth.

\synonyms
Wisdom Teeth, Impacted

\medterm{Impaction} Blockage caused by material that is firmly lodged, such as hardened stool, a foreign object, or an impacted tooth. Symptoms depend on the site and may include pain, swelling, vomiting, or inability to pass contents.
\medterm{Impaired Fasting Glucose} A fasting blood-glucose level above normal but below the range used to diagnose diabetes. It indicates increased future diabetes risk and may improve with weight management, activity, and treatment of risk factors.
\medterm{Impaired Glucose Tolerance} Blood glucose above the normal range after a glucose challenge but below the diagnostic threshold for diabetes.

\medterm{Impedance} Impedance is the opposition to the flow of an alternating electrical current; in medicine it can help assess tissue properties or body fluids.

\medterm{Impedance Hearing Testing} A group of tests, including tympanometry, that measure the movement of the eardrum and middle-ear function.

\medterm{Impedance Plethysmography} A noninvasive method that detects changes in limb electrical impedance caused by changes in blood volume. It has been used to assess venous blood flow and suspected deep-vein thrombosis, although duplex ultrasonography is generally preferred for current diagnostic evaluation.

\medterm{Impella Device} Percutaneous microaxial left ventricular assist device that actively pumps blood from
the LV into the aorta, reducing LV workload and increasing cardiac output. Types:
Impella 2.5 (up to 2.5 L/min), Impella CP (up to 3.7 L/min), Impella 5.0 (up to 5 L/min,
surgically placed).

\medterm{Imperforate Anus} A congenital anomaly where the anus is completely absent or fails to open at the
perineum. Classified as low (distal to the levator ani complex) or high (proximal to the
levator ani). Often associated with other anomalies including vertebral, anal, cardiac,
tracheoesophageal, renal, and limb defects.

\medterm{Imperforate Anus Repair} Surgery to create or reposition an anal opening and connect the rectum so stool can pass normally in a baby born with an absent or blocked anus. Some babies need temporary bowel diversion; follow-up assesses healing, constipation, and bowel control.

\medterm{Imperforate Hymen} A congenital obstruction of the vaginal opening by an intact hymen without a normal
central opening. After puberty it may cause absent menstruation, cyclic pelvic pain,
urinary retention, or a pelvic mass from retained menstrual blood.

\medterm{Imperforate Vagina} Imperforate vagina is a rare congenital absence or obstruction of the vaginal canal, often associated with genital outflow obstruction, cyclical pelvic pain, haematocolpos, and urinary or reproductive complications.

\medterm{Impetigo} Impetigo is a superficial contagious bacterial skin infection producing fragile blisters or honey-colored crusts, commonly caused by Staphylococcus aureus or Streptococcus pyogenes.
affecting children.

\synonyms
School Sores (Impetigo)


\medterm{Impingement Syndrome} Pain and restricted movement caused when a tendon or other soft tissue is compressed between bones, commonly in the shoulder during elevation of the arm.

\medterm{Implant} A medical device, tissue, or material placed inside the body to replace, support, monitor, or deliver treatment.

\medterm{Implant (Dental)} A biocompatible fixture placed in the jawbone to support a dental prosthesis. It may be used for one missing tooth or to support a partial or complete denture; complications include infection, bone loss, and mechanical failure.

\synonyms
Dental

\medterm{Implant Failure} Loss of an implant’s intended structural, biologic, or functional performance. Causes
include infection, fracture, loosening, wear, malposition, thrombosis, material
degradation, and patient or tissue factors.

\medterm{Implant Sterilization} Validated processing that removes or inactivates all viable microorganisms, including
bacterial spores, from an implant. Methods include moist heat, ethylene oxide,
radiation, and low-temperature plasma according to material and device requirements.

\medterm{Implant-Associated Infection} Infection involving an implanted medical device or the tissue surrounding it. Microbial
adherence and biofilm formation can impair antimicrobial penetration and often require
device-directed therapy, source control, or removal.

\medterm{Implant-Supported Prosthesis} A fixed or removable dental prosthesis that is supported or retained by one or more dental implants.

\medterm{Implantable Cardioverter Defibrillator} An implanted device that detects dangerous ventricular rhythms and treats them with antitachycardia pacing or an electrical shock.


\medterm{Implantable Cardioverter-Defibrillator} An implanted device that monitors heart rhythm and delivers electrical therapy
(antitachycardia pacing or defibrillation) for life-threatening ventricular arrhythmias.

\medterm{Implantable Defibrillator} A surgically implanted device that detects dangerous ventricular arrhythmias and delivers pacing or an electrical shock to restore an effective rhythm.

\medterm{Implantable Loop Recorder} Subcutaneous device providing continuous ECG monitoring for up to 3 years. Used for
evaluation of unexplained syncope, recurrent palpitations, cryptogenic stroke, and
occult atrial fibrillation detection. Small size (coin-sized), inserted under local
anesthesia in the chest.

\medterm{Implantable Port} A subcutaneously placed vascular-access device consisting of a reservoir connected to a
catheter, used for repeated administration of medication, fluids, blood products, or
parenteral nutrition. It requires sterile access and may be complicated by infection,
thrombosis, occlusion, or mechanical failure.

\medterm{Implantation} The embedding of the blastocyst into the endometrium, occurring approximately 6-10 days
after ovulation. The blastocyst hatches from the zona pellucida, and the trophoblast
invades the decidua. Implantation triggers hCG production, which maintains the corpus
luteum. Failure of implantation is a common cause of early pregnancy loss.

\medterm{Implantation Procedure} Placement of a device, tissue, medicine-delivery system, or other material inside the body for treatment or support. Examples include joint, dental, heart, and breast implants; risks include bleeding, infection, movement, failure, or rejection of donor tissue.

\medterm{Implants Endometrial (Endometriosis)} Alternate index label for Endometriosis.

\medterm{Importance Weight} An importance weight is a statistical weight used to adjust observations so a sample better represents a target population or distribution.

\medterm{Impotence (Erectile Dysfunction [ED])} The ongoing inability to get or keep an erection firm enough for sexual activity. It may result from blood-vessel disease, nerve problems, medicines, hormone changes, or emotional factors.
\synonyms
Erectile dysfunction

\medterm{Impotence (Erectile Dysfunction (ED, Impotence))} Alternate index label for Erectile Dysfunction (ED, Impotence).

\medterm{Impression} Summary of key findings and recommendations in radiology report. The concise conclusion of a medical or radiology report, summarizing the most important findings, likely meaning, and recommendations for next steps.
\medterm{Imprinting} A pattern of gene expression determined by whether a gene copy was inherited from the mother or father.

\medterm{Impulse} An impulse is a brief signal or force; in physiology it commonly refers to an electrical signal travelling along a nerve or through the heart.

\medterm{Impulse Oscillometry} A pulmonary function testing technique that measures respiratory impedance by applying
small pressure oscillations during tidal breathing. It assesses large airway resistance
(R5), small airway resistance (R5-R20), and reactance.

\medterm{Impulsive Behavior} Impulsive behavior is action with limited forethought or inhibition despite possible harm and may occur with psychiatric, neurologic, substance-related, or developmental conditions.

\medterm{Impulsivity} A tendency to act quickly without adequate consideration of consequences or difficulty delaying a response; it can occur in ADHD, mania, substance use, and other conditions.

\medterm{In} In clinical measurements, an abbreviation for inch or inches; the same letters may have other meanings in different contexts.

\medterm{In Situ} In its original place. In pathology, carcinoma in situ consists of abnormal or cancer-like cells that remain within the surface layer and have not invaded deeper tissue. It has not spread through blood or lymph vessels, but some lesions can become invasive if untreated.

\synonyms
In place

\medterm{In Situ Carcinoma} Abnormal malignant cells confined to the tissue layer in which they arose and not yet invading through the basement membrane.

\medterm{In Situ Hybridization} A laboratory method that uses labeled nucleic-acid probes to detect specific DNA or RNA sequences within cells or tissue. It can identify gene copy changes, translocations, viral material, or tissue-specific gene expression.

\medterm{In Vitro} In vitro means performed outside a living organism, such as in a test tube, culture dish, or laboratory system.

\medterm{In Vitro Fertilisation} An assisted reproductive technology in which eggs are retrieved from the ovaries,
fertilized with sperm in a laboratory, and resulting embryos are transferred to the
uterus.

\medterm{In Vitro Fertilization} In vitro fertilization fertilizes retrieved eggs in a laboratory and transfers resulting embryos to the uterus, forming part of assisted reproductive treatment.

\medterm{Intracytoplasmic Sperm Injection} A fertility treatment in which a single sperm is injected directly into an egg in a laboratory to help achieve fertilisation. It is also called ICSI.

\medterm{In Vivo} Taking place or being studied inside a living person, animal, or other organism. For example, an in vivo test examines how a medicine works in the body rather than only in a laboratory dish.

\medterm{In-Vivo} In-vivo means occurring or being studied within a living organism.

\medterm{Inactivated Vaccines} Vaccines made from microorganisms that have been killed or from parts of them. They cannot reproduce or cause the infection they are designed to prevent, but they may need several doses or booster doses to maintain protection.

\medterm{Inactivator} A substance or process that makes a biological agent inactive. Inactivators may disable viruses, toxins, enzymes, or medicines and are used in immunity, laboratory testing, sterilization, and treatment.

\medterm{Inalimarev} An investigational antiviral medicine studied against selected viral infections. Its effectiveness, approved uses, dosing, interactions, and safety remain dependent on clinical evidence and regulatory authorization.

\medterm{Inamrinone} Inamrinone is an intravenous medicine that increases heart contraction and widens blood vessels; it was used short term in severe heart failure but can cause low blood pressure and abnormal rhythms.

\medterm{Inanition} Severe physical depletion caused by inadequate food or nutrient intake, often accompanied by weight loss, weakness, dehydration, and metabolic disturbance. Causes may include starvation, malabsorption, chronic disease, or inability to eat.

\medterm{Inborn Error Of Metabolism} An inherited disorder caused by a defect in a metabolic pathway, often resulting in
accumulation of harmful substances or deficiency of essential products.

\medterm{Inborn Errors Of Metabolism} Inherited disorders caused by defects in enzymes, transporters, or other metabolic proteins, leading to abnormal buildup or deficiency of substances.

\medterm{Inbreed} To reproduce between genetically related individuals; inbreeding can increase the chance that recessive inherited disorders will occur.

\medterm{Inbreeding} Reproduction between people who share a recent ancestor, increasing the chance that a child inherits two copies of the same recessive disease-causing variant. Genetic counseling can help assess family risk.

\medterm{Inbreeding, Coefficient Of} The probability that the two copies of a gene at a locus in an individual are identical by descent from a common ancestor; it is used in population and medical genetics.

\medterm{Incarcerated Hernia} A hernia whose protruded tissue cannot be returned to the abdominal cavity. It may progress to strangulation and requires urgent clinical assessment.

\medterm{Incentive Spirometry} A breathing device that provides visual feedback while a person inhales slowly and deeply. It may be used with coughing, mobilization, and other measures to encourage lung expansion after illness or surgery; it is a lung-expansion exercise and is not the same as diagnostic spirometry.

\medterm{Incest} Sexual activity or reproduction between close biological relatives. It can involve abuse or coercion and increases the chance of inherited disorders in children; affected people need safe, confidential medical and safeguarding support.

\medterm{Incidence} Incidence is the number of new cases of a disease or event arising in a population during a specified period.

\medterm{Incidental Finding} An unexpected observation discovered during an examination performed for another reason. Its importance depends on the
finding and the clinical context.

\medterm{Incidentaloma} An unexpected mass or abnormality found during imaging or another test performed for a different reason. It may be harmless or important, so follow-up depends on its size, appearance, location, growth, and the person's risks.

\medterm{Incised Wound} Clean-cut wound produced by a sharp-edged instrument (scalpel, knife, glass) that is
longer on the skin surface than it is deep, featuring clean edges without tissue
bridging or marginal abrasion.

\medterm{Incision} A deliberate cut made by a healthcare professional through the skin or another tissue, usually during surgery or to allow drainage. The wound is closed or left to heal according to its size, location, contamination, and purpose.

\medterm{Incisional Biopsy} Surgical removal of only part of a lesion for diagnosis, generally used when the lesion is too large or located where complete removal would be inappropriate as an initial procedure.

\medterm{Incisional Hernia} A hernia through a weakness in a previous surgical incision. It may appear as a bulge that
becomes more prominent with standing or straining. Risk factors include wound infection, obesity, smoking, poor
nutrition, chronic cough, and repeated surgery.

\medterm{Incisor} A front tooth shaped mainly for cutting and biting food. Adults usually have eight incisors, and each has a crown above the gum and a root anchored in the jawbone.

\medterm{Incisors} The flat front teeth used mainly for biting and cutting food. Adults usually have eight incisors, and their health can be affected by decay, injury, gum disease, wear, or poor alignment.

\medterm{Inclusion Body} A circumscribed mass of foreign, metabolic, or abnormal material within a cell's cytoplasm or nucleus. Inclusion bodies may contain aggregated or misfolded proteins, viral components, glycogen, lipids, or other substances, and their appearance and location can help identify cellular injury or disease.

\medterm{Inclusion Body Myositis} A slowly progressive inflammatory and degenerative muscle disease, usually causing
asymmetric weakness of the finger flexors and quadriceps. Muscle biopsy may show endomysial inflammation and rimmed
vacuoles. Response to immunotherapy is often limited.

\medterm{Inclusion Conjunctivitis} Conjunctivitis caused by \textit{Chlamydia trachomatis}, usually presenting with persistent mucopurulent discharge, follicular inflammation, and sometimes genital infection. Neonatal disease requires prompt evaluation and systemic treatment to prevent complications and transmission.

\medterm{Inclusion-Cell (I-Cell) Disease} A rare inherited lysosomal-storage disorder caused by impaired tagging and trafficking of lysosomal enzymes,
leading to their secretion outside cells. It causes severe developmental delay, skeletal abnormalities, coarse facial
features, and progressive multisystem disease. Management is supportive and coordinated by specialists.

\synonyms
I-Cell

\medterm{Incompetence} Failure of a structure to close or function adequately; examples include valvular incompetence of the heart and cervical incompetence during pregnancy.

\medterm{Incompetent Cervix} Early painless opening of the cervix during pregnancy, which may lead to pregnancy loss or premature birth.

\medterm{Incomplete Miscarriage} Partial loss of a pregnancy in which some pregnancy tissue remains inside the uterus, potentially causing bleeding, pain, infection, or continued pregnancy symptoms.

\medterm{Incomplete Ossification} Failure of bone or cartilage to mineralize or develop fully. It may be part of normal development or result from genetic, metabolic, nutritional, or skeletal disorders and is interpreted with age, imaging, and laboratory findings.

\medterm{Incontinence} Involuntary loss of control over urination, bowel movements, or both. It may result from weakened muscles, nerve disease, infection, constipation, childbirth, prostate problems, medicines, mobility limits, or other conditions.

\medterm{Incontinence (Fecal)} Inability to control bowel movements, causing accidental leakage of stool or mucus. It may result from diarrhea, constipation, weakened anal muscles, nerve disease, childbirth injury, or medicines. Assessment looks for the cause; treatment may include bowel regulation, exercises, retraining, medicines, or surgery.

\synonyms
Fecal

\medterm{Incontinence (Urinary)} Alternate index label; see Urinary incontinence. Involuntary urine leakage. Types: stress (cough/sneeze), urge (overactive bladder), overflow (retention), functional (immobility/cognition), mixed. Assessment: history, bladder diary, urinalysis, post-void residual.
\synonyms
Urinary


\medterm{Incontinence, Bowel} Alternate index label; see \textit{Fecal incontinence}.

\medterm{Incontinence, Fecal} Alternate index label; see \textit{Fecal incontinence}.

\medterm{Incontinence, Urinary} Alternate index label; see \textit{Urinary incontinence}.

\medterm{Incontinentia Pigmenti} An X-linked dominant disorder affecting skin, teeth, eyes, and the nervous system,
caused by pathogenic variants in IKBKG. Skin lesions evolve through vesicular,
verrucous, hyperpigmented, and atrophic stages and may be associated with retinal
disease or seizures.

\medterm{Incoordinatio} Lack of coordination between muscles or body movements; the term is often used in older medical language.

\medterm{Incoordination} Impaired ability to perform smooth, accurate, coordinated movements, often called ataxia. Causes include cerebellar disease, sensory loss, vestibular disorders, intoxication, medicines, and other neurologic conditions.

\medterm{Incoordination (Ataxia)} Impaired coordination of voluntary movement, often caused by disease affecting the cerebellum or other parts of the nervous system.

\synonyms
Ataxia

\medterm{Increased Head Circumference} Increased head circumference means head size above the expected range for age and sex and may reflect familial size, hydrocephalus, swelling, or a genetic condition.

\medterm{Increased Intracranial Pressure} Abnormally elevated pressure within the skull caused by brain swelling, hemorrhage, mass
lesion, impaired cerebrospinal-fluid circulation, or venous obstruction. It can cause
headache, vomiting, visual disturbance, reduced consciousness, and brain herniation.

\medterm{Increased Libido} Increased libido is heightened sexual desire that may be normal or associated with hormones, medicines, mood elevation, neurologic disease, or substance use.

\medterm{Increased Osteoid} Increased osteoid is excess unmineralized bone matrix, seen when bone formation outpaces mineralization as in osteomalacia or some metabolic disorders.

\medterm{Increased Vascular Permeability} A fundamental feature of acute inflammation that allows plasma proteins, fluid, and
leukocytes to exit the vascular space and enter the extravascular tissue.

\medterm{Increment} An increase or stepwise change in an amount, measurement, dose, or treatment intensity. The size and timing of an increment should be chosen carefully to achieve benefit while limiting harm.

\medterm{Incretins} Incretins are gut hormones, mainly GLP-1 and GIP, released after meals that enhance glucose-dependent insulin secretion; GLP-1 also suppresses glucagon, slows gastric emptying, and reduces appetite.

\medterm{Incubation} Incubation is the time between exposure to an infectious agent and the appearance of symptoms.

\medterm{Incubation Period} The time between exposure to an infectious agent and the beginning of symptoms. It varies with the organism, dose, route of exposure, and person’s immune response.
\medterm{Incubator} A controlled-environment chamber culturing microorganisms, cells, or tissues at defined temperature, humidity, and atmospheric composition.
\medterm{Incus} The incus, commonly known as the anvil, is the middle ossicle of the three bones in the
middle ear, articulating with the malleus at the incudomalleolar joint and with the
stapes at the incudostapedial joint.

\medterm{Incyclinide} An investigational tetracycline-related medicine studied for effects on enzymes involved in tissue breakdown and inflammation. Its potential cancer or anti-inflammatory uses remain dependent on clinical research and regulatory approval.

\medterm{Indapamide} Indapamide is a thiazide-like diuretic used mainly for hypertension; it promotes renal sodium excretion and may cause hyponatraemia, hypokalaemia, dehydration, hyperuricaemia, or photosensitivity.

\medterm{Indel Mutation} A genetic change in which one or more DNA building blocks are inserted or deleted. If the number is not a multiple of three, it can disrupt the protein reading frame.
\medterm{Independence} The ability to perform activities and make decisions without assistance, although a person may still use devices, adaptations, or occasional support.

\medterm{Indeterminate} Indeterminate means not definitively classified or interpreted because available findings are insufficient, conflicting, or nonspecific.
\medterm{India Immunization Schedule} The Universal Immunization Schedule (UIP) of the Government of India, providing free vaccines to all children, adolescents, and pregnant women through the public health system.
\medterm{Indican} Indican is a compound formed from tryptophan metabolism and excreted in urine; abnormal levels may occur with intestinal bacterial overgrowth or metabolic disease.

\medterm{Indicate} To show, suggest, or point to a finding or conclusion; an indication is a clinical reason for using a test, treatment, or procedure.

\medterm{Indication} A clinical condition, symptom, or question that provides the reason for performing a test, procedure, or treatment.

\medterm{Indigestion} Common term for dyspepsia, discomfort or burning in the upper abdomen that may include early fullness, bloating, belching, or nausea. Persistent symptoms, bleeding, weight loss, vomiting, anemia, or difficulty swallowing require medical evaluation.

\medterm{Indigo Carmine} Indigo carmine is a blue dye used in selected medical procedures to visualize urinary or other fluid flow and in laboratory staining.

\medterm{Indinavir} An older HIV protease-inhibitor medicine that blocks production of mature virus particles. It is rarely used now because of resistance, drug interactions, kidney stones, and other adverse effects.

\medterm{Indirect Antiglobulin Test} A test that detects free, clinically significant antibodies in a patient’s serum or
plasma by incubating it with reagent erythrocytes. It is used for antibody screening,
compatibility testing, and prenatal red-cell antibody evaluation.

\medterm{Indirect Calorimetry} Measurement of oxygen consumption and carbon-dioxide production to estimate resting
energy expenditure. It is the preferred method for determining energy needs when
prediction equations are unreliable, particularly in critically ill or metabolically
complex patients.

\medterm{Indirect Veneer} A thin covering made outside the mouth and bonded to the front of a tooth to change its shape, color, or appearance. It may be made from porcelain or composite material. Tooth preparation, strength, hygiene, and long-term maintenance affect its suitability.


\medterm{Indium} Indium is a metal used in industry and medical radioisotopes; exposure effects depend on the compound, dose, and route.

\medterm{Indium-111} A radioactive form of indium used as a tracer in selected nuclear-medicine tests. It may label white blood cells to help locate inflammation or infection, or be attached to another substance to help show certain tumors. It emits detectable radiation that is recorded by a special camera.

\medterm{Indium-111 Scan} A nuclear medicine scan that uses the radioactive tracer indium-111 to make selected tissues visible. It may be attached to white blood cells to look for inflammation or infection, or to a substance that helps find some tumors.

\medterm{Indium-Labelled WBC Scan} A nuclear medicine study in which a patient’s white blood cells are labeled with indium-111 and reinjected to localize selected sites of inflammation or infection.

\medterm{Indocyanine Green} Indocyanine green is a fluorescent dye injected intravenously to assess liver function, blood flow, and retinal or choroidal circulation.

\medterm{Indole} A chemical structure containing a nitrogen atom, found in tryptophan and many natural substances and medicines. Indole-containing compounds can affect the nervous system, blood vessels, digestion, inflammation, or other body processes.

\medterm{Indole-3-Carbinol} A compound formed from glucosinolates in cruciferous vegetables and studied for effects on estrogen metabolism and cancer biology.

\medterm{Indolent} Slow to develop, progress, or cause symptoms; it often describes a relatively slow-growing disease or tumor.

\medterm{Indoles} A group of compounds containing an indole ring, including substances involved in biological signaling, metabolism, and medicines.

\medterm{Indometacin} Indometacin is a nonsteroidal anti-inflammatory medicine used for pain and inflammation; it can cause stomach bleeding, kidney problems, and cardiovascular complications.

\medterm{Indomethacin} Indomethacin is a potent NSAID used for inflammatory pain and selected acute conditions such as gout; gastrointestinal ulceration or bleeding, renal injury, cardiovascular events, and headache are important adverse effects.

\medterm{Indomethacin Overdose} Toxic exposure to excessive indomethacin, which can cause nausea, vomiting, stomach bleeding, kidney injury, seizures, or other serious effects.

\medterm{Indoor Allergens} Substances inside buildings that can trigger allergic or asthma symptoms, including dust mites, mold, cockroach particles,
and animal dander. Exposure may cause rhinitis, conjunctivitis, wheezing, or eczema. Reduction of exposure and
appropriate allergy or asthma treatment depend on the identified trigger.


\medterm{Indoximod} Indoximod is an investigational immunotherapy that modulates tryptophan metabolism and has been studied in combination cancer treatment.

\medterm{Induced Abortion} Deliberate termination of a pregnancy using medicines or a procedure. Safe care depends on gestational age, patient preference, accurate information, legal context, trained providers, and assessment for complications such as heavy bleeding or infection.

\medterm{Induced Hyperthermia} Induced hyperthermia is deliberate elevation of body or tissue temperature for a therapeutic purpose, with careful monitoring to prevent heat injury.

\medterm{Induced Labor} Medical initiation of uterine contractions before labor begins naturally to achieve delivery when continuing pregnancy may pose greater risk.

\medterm{Induced Sputum Analysis} A method of assessing airway inflammation by examining sputum produced after inhaled
saline stimulation. The types and proportions of inflammatory cells may support assessment of asthma, infection, COPD,
or bronchiectasis.

\medterm{Inducing Labor} Starting childbirth artificially with medicines, breaking the amniotic sac, or other methods when continuing pregnancy may pose greater risks than delivery.

\medterm{Induction} Initiation of a physiologic or therapeutic process, such as induction of labor, anesthesia, remission, or immune tolerance. The method, indication, benefits, and risks depend on the clinical context.

\medterm{Induction Of Labor} Starting uterine contractions medically before labor begins on its own. It may be recommended when continuing pregnancy is riskier, using medicines or procedures that soften the cervix and encourage contractions.

\medterm{Induction Therapy} The first, often intensive, part of treatment given to reduce or remove a disease. In acute leukemia, its main aim is remission, meaning that tests find no clear evidence of leukemia; further treatment is usually needed to reduce the chance of return.

\medterm{Indurated} Abnormally firm or hardened when palpated, usually because of inflammation, edema, fibrosis, infiltration, or a tumor within the skin or underlying tissue.
\medterm{Induration} Abnormal hardening of tissue, often detected by palpation and caused by inflammation, infection, edema, fibrosis, or a mass.

\medterm{Industrial Bronchitis} Industrial bronchitis is chronic airway irritation and inflammation related to workplace dusts, fumes, or chemicals, causing cough, phlegm, and breathing limitation.

\medterm{Industrial Disease} A disease caused or worsened by workplace exposures or conditions, such as chemicals, dust, noise, radiation, infection, or repetitive strain.

\medterm{Industrial Health} The branch of occupational health concerned with preventing and managing illness and injury related to industrial work, workplace exposures, and working conditions.

\medterm{Indwelling Catheter} A catheter left in place for ongoing drainage, infusion, monitoring, or access; urinary indwelling catheters require appropriate indications and care because they increase infection and trauma risk.

\medterm{Indwelling Catheter Care} A patient-care topic concerning indwelling catheter care. It addresses the relevant purpose, risks, management, and follow-up.

\medterm{Indwelling Pleural Catheter} A flexible tube placed through the chest wall into the space around the lung to drain fluid that keeps returning. It can allow drainage at home and relieve breathlessness, especially when cancer or a lung that cannot fully expand causes the fluid. Possible problems include infection, blockage, pain, or air leakage.

\synonyms
IPC; Tunneled Pleural Catheter; PleurX Catheter

\medterm{Inertia} Resistance to a change in motion; in obstetrics, uterine inertia means weak or ineffective contractions that slow labor.

\medterm{Infant} An infant is a child in the first year of life, a period of rapid growth and changing feeding, immune, neurologic, and developmental needs.

\medterm{Infant - Newborn Development} The early growth of a newborn’s body, brain, senses, movement, feeding, communication, and social responses during the first weeks and months.

\medterm{Infant Botulism} A neuromuscular illness caused by intestinal colonization with Clostridium botulinum and
production of botulinum neurotoxin in infants. It causes constipation, poor feeding,
weak cry, hypotonia, facial weakness, and descending paralysis.

\medterm{Infant Feeding} Provision of breast milk, appropriately prepared formula, or both to meet an infant's
energy, protein, fluid, and micronutrient requirements. Feeding recommendations depend
on age, growth, medical conditions, developmental readiness, and safe preparation
practices.

\medterm{Infant Formula} Infant formula is manufactured nutrition designed to provide adequate nutrients when breast milk is unavailable, insufficient, or not chosen; preparation and storage must be safe.


\medterm{Infant Formulas} Commercial preparations designed to provide nutrition for infants when breast milk is unavailable, insufficient, or not chosen.

\medterm{Infant Gynecomastia} Infant gynecomastia is temporary breast enlargement in a newborn or young infant caused by maternal or neonatal hormones and usually resolves spontaneously.

\medterm{Infant Harlequin Syndrome} A transient, benign cutaneous vascular phenomenon in neonates characterized by
asymmetric flushing of one half of the body with sharp midline demarcation, caused by
transient asymmetric sympathetic tone. The affected side is red and warm while the
contralateral side is pale.

\medterm{Infant Incubators} Enclosed devices providing controlled warmth, humidity, oxygen, and monitoring for premature, small, or ill newborns who need extra support after birth.

\medterm{Infant Jaundice} Yellowing of an infant’s skin or eyes caused by excess bilirubin. It is often temporary, but jaundice in the first day, rapid worsening, poor feeding, unusual sleepiness, or pale stools requires urgent assessment.
\medterm{Infant Leukemia} Infant leukemia is leukemia diagnosed during infancy, in which abnormal blood-forming cells multiply and interfere with normal blood-cell production.

\medterm{Infant Mortality} The death of an infant before the first birthday; the infant mortality rate is usually expressed per 1,000 live births.

\medterm{Infant Nutrition} Nutritional support for infants through breast milk, appropriately prepared formula, and developmentally suitable complementary foods. Requirements change with age, growth, illness, prematurity, and feeding ability; unsafe preparation can cause infection or malnutrition.

\medterm{Infant Of A Substance-Using Mother} A newborn exposed to alcohol or other substances during pregnancy who may need assessment for withdrawal, toxicity, feeding problems, or developmental effects.

\medterm{Infant Of Diabetic Mother} A newborn born to a mother with diabetes who may be at increased risk of large size, low blood sugar, breathing problems, jaundice, or congenital abnormalities.

\medterm{Infant Reflex} An infant reflex is an automatic response present during early development, such as rooting, sucking, grasping, or startle, whose timing and symmetry help assess neurologic function.

\medterm{Infant Reflexes} Automatic newborn responses such as rooting, sucking, Moro, grasp, and stepping that help assess neurologic development.

\medterm{Infant Reflux} Backflow of stomach contents into an infant’s esophagus, often causing uncomplicated spitting up. It usually improves with age, but poor growth, breathing problems, blood, forceful vomiting, or distress requires medical assessment.

\medterm{Infant Sleep Position} Placing an infant on the back for every sleep is recommended to reduce the risk of sleep-
related death. Infants should sleep on a firm, separate surface without pillows, loose bedding, bumpers, or stuffed
toys.

\medterm{Infant Test Procedure Preparation} Infant test procedure preparation includes age-appropriate instructions, feeding or fasting guidance, comfort measures, and safety steps before a test.

\medterm{Infant, Post-Term} An infant born at or after 42 weeks of gestation, beyond the usual term-pregnancy period; risks depend on placental function, fetal condition, and circumstances of delivery.

\medterm{Infanticide} The killing of a child, in law typically defined as the killing of a newborn or infant
(usually within the first year, or specifically within 24 hours of birth in some
jurisdictions) by its mother while mentally disturbed by the effects of childbirth or
lactation.

\medterm{Infantile Colic} A behavioral syndrome in healthy infants characterized by episodes of excessive,
inconsolable crying for >3 hours per day, >3 days per week, for >3 weeks (Wessel
criteria). Onset typically at 2-3 weeks, peaks at 6-8 weeks, and resolves by 3-4 months
of age.

\medterm{Infantile Fibrosarcoma} Infantile fibrosarcoma is a rare soft-tissue cancer that usually develops in infants or young children and requires specialist treatment.

\medterm{Infantile Hemangioma} Alternate index label; see \textit{Hemangioma}.

\medterm{Infantile Hepatic Hemangiomas} Benign growths made of extra blood vessels in an infant’s liver. Small lesions may cause no symptoms, while numerous or large lesions can enlarge the liver, strain the heart, lower thyroid hormone, or cause bleeding problems. Ultrasound and specialist follow-up guide monitoring and treatment.

\medterm{Infantile Paralysis (Polio)} Paralytic poliomyelitis caused by poliovirus infection involving motor neurons, which can produce acute flaccid paralysis and respiratory failure; vaccination prevents most cases.

\medterm{Infantile Spasm} Infantile spasms are brief epileptic seizures in infancy, often in clusters with flexion or extension movements, and require urgent diagnosis and treatment to protect development.

\medterm{Infantile Spasms} A severe seizure disorder of infancy characterized by brief, sudden flexion or extension spasms, often occurring in clusters and associated with developmental impairment.


\medterm{Infarct} An area of dead tissue caused by inadequate blood supply.

\medterm{Infarction} Tissue death caused by inadequate blood flow and oxygen, usually because a vessel is blocked. Examples include heart attack, brain infarction, and intestinal infarction.
infarction.

\medterm{Infection} Entry and multiplication of bacteria, viruses, fungi, parasites, or other microorganisms in the body. Infection may cause tissue damage and illness, although some infections produce no symptoms; the cause and site determine treatment.

\medterm{Infection (Geriatric)} Infectious diseases in elderly often present atypically: fever may be absent, confusion
may be sole symptom. Common infections: UTI, pneumonia, skin/soft tissue, bloodstream.
Prevention: vaccination, hand hygiene, aspiration precautions. Treatment: earlier
presentation, broader differential, drug interactions, renal dosing.

\synonyms
Geriatric


\medterm{Infection Control} Practices that prevent microorganisms from spreading between people, objects, and healthcare settings. They include hand cleaning, appropriate protective equipment, safe injection and device care, cleaning, sterilization, vaccination, and isolation when needed.

\medterm{Infection Cryptococcus (Cryptococcosis)} Alternate index label for Cryptococcosis.

\medterm{Infection Cyclospora (Cyclospora Infection (Cyclosporiasis))} Alternate index label for Cyclospora Infection (Cyclosporiasis).

\medterm{Infection Enterovirus (Enterovirus (Non-Polio Enterovirus Infection))} Alternate index label for Enterovirus (Non-Polio Enterovirus Infection).

\medterm{Infection HPV (HPV Infection Human Papillomavirus)} Alternate index label for HPV Infection Human Papillomavirus.

\medterm{Infection Human Papillomavirus (HPV Infection Human Papillomavirus)} Alternate index label for HPV Infection Human Papillomavirus.

\medterm{Infection Imaging} Imaging used to locate an infection when examination and routine tests do not show its source. Depending on the suspected problem, it may use ultrasound, CT, MRI, or a nuclear-medicine tracer; scan findings must be interpreted with symptoms, cultures, and other tests.

\medterm{Infection Listeria (Listeria)} Alternate index label for Listeria.

\medterm{Infection Non-Polio Enterovirus (Enterovirus (Non-Polio Enterovirus Infection))} Alternate index label for Enterovirus (Non-Polio Enterovirus Infection).

\medterm{Infection Pinworms (Pinworm Infection)} Alternate index label for Pinworm Infection.

\medterm{Infection Prevention} Measures that reduce the risk of infection, including hand hygiene, personal protective equipment, environmental cleaning, sterilization, disinfection, vaccination, and appropriate isolation precautions.


\medterm{Infection Testicle (Testicular Disorders)} Alternate index label for Testicular Disorders.

\medterm{Infection, Hospital-Acquired} An infection acquired during or as a result of healthcare that was not present or incubating at the time of admission; it may be caused by bacteria, viruses, fungi, or other organisms.

\medterm{Infection, Urinary Tract} Infection of part of the urinary system, commonly the bladder or urethra and sometimes the kidneys, usually caused by bacteria and producing symptoms such as burning urination, urgency, or flank pain.

\medterm{Infections In Pregnancy} Infections during pregnancy can affect the pregnant person, placenta, fetus, or newborn. Some, including listeriosis, group B streptococcal infection, urinary infection, and sexually transmitted infections, can cause miscarriage, premature birth, fetal illness, or newborn sepsis. Testing and treatment depend on the organism and stage of pregnancy.

\medterm{Infectious Arteritis} Inflammation and damage of an artery caused by infection, which can weaken the vessel, reduce blood flow, or produce aneurysm.

\medterm{Infectious Arthritis} Inflammation of a joint caused by infection, usually bacterial and sometimes viral or fungal.
Acute pain, swelling, warmth, fever, and restricted movement require urgent evaluation because
joint damage can develop quickly.

\medterm{Infectious Conjunctivitis} Inflammation of the conjunctiva caused by a virus, bacterium, or less commonly another infectious organism, often causing redness and discharge.

\medterm{Infectious Disease} Illness caused by a pathogen such as a bacterium, virus, fungus, or parasite. Infection may spread between people, from animals or the environment, or from organisms normally living in the body.

\medterm{Infectious Dose} The number of viable organisms required to establish infection in a host under specified
conditions. The dose depends on the pathogen, route of exposure, host defenses, and
organism viability.

\medterm{Infectious Encephalitis} Brain inflammation caused by an infectious organism, most often a virus. It may cause fever, headache, confusion, seizures, behavior changes, weakness, or reduced consciousness and requires urgent treatment.

\medterm{Infectious Encephalomyelitis} Infection-related inflammation of the brain and spinal cord that may cause fever, headache, confusion, weakness, seizures, or other neurologic signs.

\medterm{Infectious Enterocolitis} Infection-related inflammation of the small intestine and colon, commonly causing diarrhea, abdominal pain, fever, or vomiting.

\medterm{Infectious Esophagitis} Inflammation of the esophagus caused by an infectious organism, most often Candida,
herpes simplex virus, or cytomegalovirus in immunocompromised patients. It commonly
causes painful swallowing, difficulty swallowing, and retrosternal pain.

\medterm{Infectious Hepatitis} Alternate index label; see \textit{Viral hepatitis}.

\medterm{Infectious Mediastinitis} Infection and inflammation of the tissues in the middle of the chest, often after esophageal injury, surgery, or spread from a nearby infection.

\medterm{Infectious Meningitis} Infection and inflammation of the membranes surrounding the brain and spinal cord. Bacterial
meningitis is a medical emergency; headache, fever, neck stiffness, vomiting, confusion, or a
new rash requires immediate medical attention.

\medterm{Infectious Mononucleosis} A viral illness, usually caused by Epstein-Barr virus, characterized by fatigue, sore throat, and swollen lymph nodes.

\synonyms
Mono (Infectious Mononucleosis)

\medterm{Infectious Myringitis} Inflammation and infection of the tympanic membrane, often associated with viral or
bacterial respiratory infection. It can cause severe ear pain, fever, hearing reduction,
and small fluid-filled blisters on the eardrum.

\medterm{Infective Cervicitis} Infective cervicitis is inflammation of the cervix caused by organisms such as chlamydia, gonorrhea, trichomonas, or herpes, producing discharge, bleeding, or pelvic discomfort.

\medterm{Infective Cystitis} Infective cystitis is bacterial or other microbial inflammation of the bladder, causing urinary frequency, urgency, burning, and sometimes blood or fever.

\medterm{Infective Endocarditis} Infection of the endocardial surface of the heart, most commonly affecting the heart
valves. It is diagnosed using the Duke criteria, which include major criteria (positive
blood cultures, echocardiographic evidence of vegetations) and minor criteria
(predisposing heart condition, fever, vascular phenomena, immunologic phenomena).

\medterm{Infective Phlebitis} Infective phlebitis is infection and inflammation of a vein, often around an intravenous catheter, causing pain, redness, swelling, and possible bloodstream infection.

\medterm{Infective Vaginitis} Infective vaginitis is infection or inflammation of the vagina caused by yeast, bacteria, parasites, or viruses and may cause discharge, itching, odor, or pain.

\medterm{Inferior} Below or nearer the feet compared with another body part. For example, the diaphragm is inferior to the heart, and the inferior vena cava carries blood from the lower part of the body to the heart.

\medterm{Inferior Alveolar Nerve Block} A regional anesthetic injection near the mandibular foramen to anesthetize the inferior
alveolar nerve before it enters the mandibular canal. It usually anesthetizes mandibular
teeth on that side; the lingual nerve is often anesthetized incidentally.

\medterm{Inferior Colliculus} The inferior colliculus is a midbrain auditory relay that integrates sound information and helps orient reflexively to auditory stimuli.

\medterm{Inferior Hemifield Loss} Loss of vision in the lower half of the visual field in one or both eyes, which
localises to lesions of the superior retina or superior optic radiations passing through
the parietal lobe. Causes include superior retinal detachment, branch retinal artery
occlusion, and parietal lobe stroke.

\medterm{Inferior Mesenteric Artery} Third major abdominal aortic branch at L3, the hindgut artery. Supplies distal
transverse colon through rectum. Branches: left colic, sigmoid, superior rectal. The
marginal artery of Drummond provides SMA-IMA collateral circulation. Occlusion may cause
ischemic colitis at the splenic flexure watershed (Griffiths' point).

\medterm{Inferior Myocardial Infarction} A heart attack affecting the lower surface of the heart muscle. It usually results from blocked blood flow in a coronary artery and requires urgent treatment to limit permanent damage.

\medterm{Inferior Vena Cava} Body's largest vein, formed by common iliac vein union at L5, returning lower body
deoxygenated blood to the right atrium. Ascends right of the aorta, pierces the
diaphragm at T8, enters the pericardium. Receives hepatic, renal, gonadal, and lumbar
veins. IVC thrombosis/compression causes bilateral lower extremity edema and DVT.

\medterm{Inferolateral} Located below and toward the side of a reference structure.

\medterm{Infertility} Infertility is failure to achieve pregnancy after a defined period of regular unprotected intercourse, reflecting factors involving ovulation, sperm, reproductive anatomy, age, or unexplained causes.
to conceive or carry a pregnancy. Causes may involve ovulation, reproductive anatomy, sperm, age-related fertility,
medical conditions, or unexplained factors.

\synonyms
Subfertility


\medterm{Infertility, Female} Alternate index label; see \textit{Female infertility}.

\medterm{Infertility, Male} Alternate index label; see \textit{Male infertility}.

\medterm{Infestation} The presence or invasion of parasitic organisms, such as lice, mites, or worms, on or in the body.

\medterm{Infiltrate} An area of increased tissue density seen on an image, often in the lung. It is nonspecific and may represent infection, fluid, bleeding, inflammation, collapse, tumor, or another process.

\medterm{Infiltrating Cancer} Cancer that has grown beyond its original location into nearby tissue. Such growth may damage normal structures and can allow cancer cells to enter lymph vessels or blood vessels and spread elsewhere.

\synonyms
Invasive cancer

\medterm{Infiltrating Ductal Carcinoma} Alternate index label; see \textit{Invasive ductal carcinoma}.

\medterm{Infiltration} Entry or accumulation of cells, fluid, or another substance within tissue. In clinical care, intravenous infiltration means leakage of nonvesicant fluid into surrounding tissue; extravasation of an irritant or vesicant can cause serious injury.

\medterm{Inflamed Gallbladder} Alternate index label; see \textit{Cholecystitis}.

\medterm{Inflamed Pancreas} Alternate index label; see \textit{Pancreatitis}.

\medterm{Inflamed Pericardium} Alternate index label; see \textit{Pericarditis}.

\medterm{Inflammaging} Persistent, low-grade, age-associated inflammatory activity linked to cellular
senescence, immune dysregulation, chronic infection, tissue injury, and comorbidity. It
contributes to frailty, atherosclerosis, sarcopenia, and neurodegenerative disease.

\medterm{Inflammasomes} Groups of proteins inside immune cells that detect danger signals and activate inflammation. They help release inflammatory messenger proteins and can cause an infected or damaged cell to die. Excessive activation may contribute to autoimmune or inflammatory disease.

\medterm{Inflammation} The body’s protective response to injury, infection, or irritation. It may cause redness, warmth, swelling, pain, and reduced function. Short-term inflammation helps healing, while persistent inflammation can damage tissues and contribute to chronic disease.

\medterm{Inflammation (Dietary)} Chronic low-grade inflammation linked to diet. Western diet pro-inflammatory (saturated
fat, sugar, refined grains). Mediterranean diet anti-inflammatory.

\synonyms
Dietary

\medterm{Inflammation Imaging} Imaging used to find areas of active inflammation in the body. A scan using a sugar-like tracer may show infection or inflammation, but cancer can look similar, so scan findings must be interpreted with symptoms, examination, and other tests.

\medterm{Inflammation Of The Heart Valve (Endocarditis)} Alternate index label for Endocarditis.

\medterm{Inflammation Sclera (Scleritis)} Alternate index label for Scleritis.

\medterm{Inflammation Vaginal (Vaginitis Overview)} Alternate index label for Vaginitis Overview.

\medterm{Inflammatory Back Pain} A pattern of chronic back pain characterized by insidious onset, improvement with
exercise, lack of improvement with rest, nocturnal pain, and prolonged morning
stiffness. It suggests axial spondyloarthritis but is not diagnostic without supportive
clinical, laboratory, and imaging findings.

\medterm{Inflammatory Bowel Disease} A group of chronic conditions that cause inflammation in the digestive tract, mainly Crohn disease and ulcerative colitis. Symptoms may come and go and can include abdominal pain, diarrhea, rectal bleeding, and weight loss. The exact cause is unknown, but genes, the immune system, the environment, and gut bacteria may all play a role.

\medterm{Inflammatory Bowel Disease Diet} Crohn disease, ulcerative colitis. Nutritional support during flares. Exclusive enteral
nutrition for Crohn induction. Individual triggers vary.

\medterm{Inflammatory Bowel Disease: Intestinal Problems} Inflammatory bowel disease, including Crohn disease and ulcerative colitis, causes chronic immune-mediated intestinal inflammation with diarrhea, pain, bleeding, weight loss, and extraintestinal effects. Treatment aims for sustained remission and prevention of structural and nutritional complications.

\synonyms
IBD (Inflammatory Bowel Disease: Intestinal Problems)

\medterm{Inflammatory Breast Cancer} A rare, aggressive breast cancer causing redness, swelling, warmth, and thickened skin, sometimes without a distinct lump. It requires prompt imaging, biopsy, staging, and coordinated treatment.

\synonyms
Breast Cancer, Inflammatory
Cancer, Inflammatory Breast

\medterm{Inflammatory Myopathies} A group of disorders in which inflammation damages muscle, usually causing gradually worsening weakness around the hips, thighs, shoulders, or upper arms. Some forms also affect skin, lungs, swallowing, or other organs. Diagnosis and treatment require specialist assessment.

\medterm{Inflammatory Response} The coordinated vascular and immune reaction to infection, tissue injury, or other danger signals, producing mediator release, leukocyte recruitment, and changes in blood flow and permeability.

\medterm{Infliximab} A biologic medicine that blocks tumor necrosis factor, an inflammatory signal. It treats several inflammatory diseases but can increase the risk of serious infections, tuberculosis reactivation, infusion reactions, and heart-failure worsening.

\medterm{Influenza} A contagious respiratory illness caused by influenza A or B viruses, characterized by
sudden onset of fever, myalgia, headache, malaise, nonproductive cough, sore throat, and
nasal congestion. Seasonal epidemics occur annually. Complications: pneumonia (primary
viral or secondary bacterial), exacerbation of underlying conditions.

\medterm{Influenza (Flu)} A contagious respiratory infection caused by influenza viruses, commonly producing fever, cough, body aches, fatigue, and respiratory symptoms.



\medterm{Influenza A Virus} An influenza virus with segmented RNA that infects people and animals and can cause seasonal epidemics or, after major genetic changes, pandemics.

\medterm{Influenza Test} A test for influenza A or B using a respiratory specimen. Rapid antigen tests provide quick results
but may miss infection. Molecular tests are generally more sensitive and may identify or
subtype influenza viruses; results must be interpreted with symptoms, timing, and local
laboratory methods.

\medterm{Influenza Vaccination In Pregnancy} Annual inactivated influenza vaccine recommended for all pregnant women. Pregnancy
increases the risk of severe influenza (pneumonia, ICU admission, death) due to changes
in immune function and cardiopulmonary physiology.

\medterm{Influenza Virus} A segmented, negative-sense RNA virus belonging to the Orthomyxoviridae family.
Influenza causes seasonal epidemics and occasional pandemics of respiratory illness. The
virus has two surface glycoproteins: hemagglutinin (HA), which mediates attachment, and
neuraminidase (NA), which facilitates release.

\medterm{Influenza, Avian} Alternate index label; see \textit{Bird flu (avian influenza)}.

\medterm{Influenza, H1N1} Alternate index label; see \textit{H1N1 flu (swine flu)}.

\medterm{Influenza, Swine Flu} Alternate index label; see \textit{H1N1 flu (swine flu)}.

\medterm{Influenzavirus A} The virus genus containing influenza A viruses, which are classified into subtypes by their hemagglutinin and neuraminidase proteins.

\medterm{Influenzavirus B} The virus genus containing influenza B viruses, which mainly infect humans and cause seasonal respiratory illness.

\medterm{Informatics} Informatics applies information science, data, computing, and human factors to organize and improve healthcare, biology, research, or another field.

\medterm{Informative} Providing useful information or evidence that improves understanding, decision-making, diagnosis, treatment planning, research, or evaluation of a health problem.

\medterm{Informed Consent} A patient's voluntary agreement to a proposed intervention after receiving understandable information about its purpose, benefits, risks, alternatives, and the option of declining, with decision-making capacity or appropriate authorization.

\medterm{Informed Consent (Geriatric)} Process of obtaining voluntary authorization from elderly patient with decision-making
capacity. Capacity assessment: understanding, appreciation, reasoning, communication of
choice. Impaired capacity may require surrogate decision-maker.

\synonyms
Geriatric

\medterm{Informed Consent - Adults} Informed consent requires an adult patient to receive understandable information about a proposed intervention, benefits, risks, alternatives, and refusal before deciding voluntarily.

\medterm{Infrared Rays} Infrared rays are electromagnetic radiation with wavelengths longer than visible red light; excessive exposure can heat tissue and cause burns.

\medterm{Infrared Spectrometry} A method for identifying or measuring substances by examining how they absorb infrared light. Different chemical bonds absorb different wavelengths, creating a pattern that can help analyze medicines, tissues, samples, or pollutants.

\medterm{Infrarenal Aorta} The infrarenal aorta is the portion of the abdominal aorta below the renal arteries and above its division into the iliac arteries.

\medterm{Infratentorial Neoplasm} An infratentorial neoplasm is an abnormal growth in the posterior fossa, the part of the skull containing the cerebellum and brainstem.

\medterm{Infrequent} Infrequent means occurring less often than usual or expected, such as infrequent bowel movements, symptoms, or events.

\medterm{Infundibulopelvic Ligament} A fold of tissue running from each ovary to the side wall of the pelvis. It carries the ovary’s blood vessels, lymph vessels, and nerves; surgeons must protect the nearby ureter during pelvic operations.

\synonyms
Suspensory Ligament of the Ovary; IP Ligament

\medterm{Infuse} To introduce a fluid, medicine, or other substance gradually into the body, commonly through an intravenous line; it can also mean to soak a substance in a liquid.

\medterm{Infusion} Slow delivery of a fluid or medicine into a vein, under the skin, or another body site. It allows controlled dosing over time and may require monitoring for allergic reactions, infection, or fluid overload.

\synonyms
Intravenous infusion

\medterm{Infusion Pump} A device that delivers a measured amount of medicine or fluid at a controlled rate, usually into a vein, under the skin, or around the spinal cord. It can improve dosing accuracy but requires correct programming and monitoring for blockage, leakage, infection, or excessive delivery.

\medterm{Ingestion} The act of taking a substance into the digestive tract by swallowing. Ingestion of a toxic substance can cause poisoning.

\medterm{Ingredient} An ingredient is a substance included in a medicine, food, cosmetic, or other preparation, whether active or inactive.

\medterm{Ingrown Hair} A hair that curls back into or grows sideways beneath the skin, causing an itchy or painful bump. It may become inflamed or infected, especially after shaving, waxing, or close trimming.
\medterm{Ingrown Nail} An ingrown nail occurs when the nail edge grows into surrounding skin, causing pain, redness, swelling, and sometimes infection or granulation tissue.

\medterm{Ingrown Toenail} Ingrown toenail (onychocryptosis) occurs when the nail edge grows into the surrounding
skin, causing pain, redness, and swelling. It is common in adolescents and adults,
exacerbated by improper nail trimming, tight shoes, and trauma. Treatment includes
partial nail avulsion with matrixectomy (phenol application) for recurrent cases.

\synonyms
Onychocryptosis (Ingrown Toenail)


\medterm{Inguinal} Relating to the groin, the region where the lower abdomen meets the upper thigh and where certain hernias, lymph nodes, and injuries occur.

\medterm{Inguinal Canal} An oblique passage (4 cm) through the anterior abdominal wall, from the deep (internal)
inguinal ring to the superficial (external) inguinal ring, transmitting the spermatic
cord in males and round ligament in females, plus ilioinguinal and genital branch of
genitofemoral nerves.

\medterm{Inguinal Hernia} Protrusion of abdominal contents through the inguinal canal. Indirect (lateral to
inferior epigastric vessels, through deep inguinal ring): congenital, patent processus
vaginalis; most common overall (60-70\%), male predominance.

\medterm{Inguinal Hernia Repair} Surgery that returns tissue to the abdomen and closes or strengthens a weak area in the groin. It may be performed through a larger incision or several small ones, with or without mesh; recurrence, infection, and persistent pain are possible.


\medterm{Inguinal Orchiectomy} Inguinal orchiectomy is removal of a testicle through an incision in the groin, commonly used when treating suspected testicular cancer.

\medterm{Inguinal Region} Area at the thigh-trunk junction traversed by the inguinal ligament (ASIS to pubic
tubercle). Clinically significant for inguinal hernias: indirect (through deep inguinal
ring, most common) and direct (through Hesselbach's triangle).

\medterm{Inhalant} A substance breathed into the lungs as a gas, vapor, aerosol, or smoke; inhalants may be medicines, workplace substances, or abused chemicals.

\medterm{Inhalation} Drawing air, medicine, particles, or vapors into the lungs. Inhalation is a route for respiratory medicines and anesthetics, while accidental inhalation of food, vomit, smoke, or toxic substances can cause airway obstruction or lung injury.

\medterm{Inhalation Injury} Thermal or chemical injury to the airway or lungs caused by smoke, heated gases, or
toxic combustion products. It can produce upper-airway edema, bronchospasm, impaired
mucociliary clearance, and carbon-monoxide or cyanide toxicity and may worsen after the
initial burn.

\medterm{Inhalational Anesthesia} Anesthetic maintenance or induction using volatile liquids or gases delivered through
the respiratory tract. Depth is influenced by inspired concentration, alveolar
ventilation, uptake, distribution, and the agent's blood-gas solubility.

\medterm{Inhaled Corticosteroids} Anti-inflammatory medications delivered directly to the airways, reducing airway
inflammation, mucus production, and bronchial hyperresponsiveness. Examples include
fluticasone, budesonide, beclomethasone, and mometasone. They are the cornerstone of
asthma maintenance therapy. Side effects include oral candidiasis and dysphonia.

\medterm{Inhaled Nitric Oxide} A selective pulmonary vasodilator used in acute settings (acute respiratory failure,
right heart failure, post-cardiac surgery). It reduces pulmonary vascular resistance and
improves oxygenation by dilating well-ventilated pulmonary vessels (improving V/Q
matching). It has a short half-life and requires continuous delivery.

\medterm{Inhaler} A device that delivers medicine into the airways through the mouth. Pressurised inhalers release a measured spray, while dry-powder inhalers release powder during a strong breath in; correct technique is important for the medicine to work.

\medterm{Inherited Disease} A disease caused wholly or partly by genetic changes passed from parents or arising in a reproductive cell or early embryo.

\medterm{Inhibin} A glycoprotein hormone produced mainly by the gonads that suppresses pituitary follicle-stimulating hormone secretion and contributes to regulation of reproductive function.

\medterm{Inhibition} Inhibition is a reduction or suppression of an activity, such as nerve signalling, enzyme action, muscle contraction, or cell growth.

\medterm{Inhibitory Neurochemical} A chemical messenger that decreases the activity of nerve cells or muscles. Gamma-aminobutyric acid and glycine are important inhibitory neurotransmitters in the nervous system.

\synonyms
Inhibitory neurotransmitter

\medterm{Inhibitory Synapse} A connection between nerve cells that reduces the chance that the receiving cell will generate an electrical signal.

\medterm{Iniencephaly} A rare, usually lethal neural-tube defect characterized by occipital bone defect, severe
retroflexion of the head, and marked shortening or distortion of the cervical and
thoracic spine. It is often identified prenatally and may occur with other congenital
anomalies.

\medterm{Iniparib} An investigational anticancer drug studied for effects on DNA repair and tumor-cell survival; it is not a standard treatment for most cancers.

\medterm{Initial Insomnia} Initial insomnia is difficulty falling asleep at the beginning of the sleep period, often related to stress, habits, medicines, mood, or another sleep disorder.

\medterm{Injectable Fillers} Materials placed under or within the skin to restore volume, improve contours, or soften wrinkles. Effects vary by product and may include swelling, bruising, infection, lumps, or rare blood-vessel blockage.

\synonyms
Dermal fillers

\medterm{Injection} Introduction of a medication, fluid, or other substance into the body with a needle or other injection device.

\medterm{Injury} Damage to the body's tissues or function caused by trauma, force, heat, chemicals, radiation, or another harmful exposure.

\medterm{Injury - Kidney And Ureter} Damage to the kidney or ureter from trauma, obstruction, surgery, infection, or other causes; severe injury can impair urine flow or kidney function.

\medterm{Injury Classification (Abrasion, Laceration, Contusion, Incised Wounds)} Forensic categorization of mechanical wounds. Abrasion: superficial friction injury
affecting only the epidermis (e.g., scrape, graze, pattern abrasion).

\synonyms
Abrasion, Laceration, Contusion, Incised Wounds

\medterm{Injury Dating} Estimation of the age of an injury using clinical appearance, histopathology, imaging,
and the history of symptoms. Exact dating is often uncertain, particularly for bruises
and soft-tissue injuries, and findings should be reported with appropriate limitations.

\medterm{Injury Site} The body location where physical harm occurred or where its effects are being examined, documented, treated, and monitored during recovery.

\medterm{Ink Poisoning} Illness after swallowing or significant exposure to ink; most small pen-ink exposures cause little toxicity, but symptoms depend on the product and amount.

\medterm{Ink Remover Poisoning} Poisoning caused by chemicals used to remove ink, which may irritate the skin, eyes, lungs, or digestive tract and may cause systemic toxicity.

\medterm{Innate Immune Evasion} Mechanisms by which pathogens evade the innate immune system. Strategies include
capsules that prevent phagocytosis, biofilms that protect against immune cells, and
intracellular survival within macrophages. Some bacteria produce enzymes that degrade
complement components. Others modify their surface molecules to avoid recognition.

\medterm{Innate Immunity} First line of defense against pathogens, present from birth. Rapid (minutes to hours),
germline-encoded, lacks immunological memory.

\medterm{Innate Lymphoid Cells} ILCs: lymphoid lineage, no rearranged antigen receptors, mirror Th subsets. ILC1 (T-bet,
IFN-gamma): intracellular pathogens, IBD. ILC2 (GATA3, IL-5/13): helminths, allergy,
asthma, tissue repair. ILC3 (ROR-gammat, IL-17/22): extracellular bacteria/fungi,
lymphoid organogenesis.

\medterm{Inner Ear} The deepest part of the ear, containing the cochlea for hearing and balance organs for movement sensing. Disease there can cause hearing loss, ringing, dizziness, or balance problems.

\medterm{Inner Ear Infection} Inner-ear inflammation can affect hearing and balance, producing vertigo, nausea, tinnitus, or hearing loss. Causes include viral infection, bacterial disease, autoimmune inflammation, or other disorders; sudden hearing loss, severe imbalance, or neurologic signs requires urgent assessment.

\medterm{Innervation} The nerve supply to an organ, muscle, or area of the body, including the functional control provided by those nerves.

\medterm{Innominate} Innominate means unnamed; in anatomy it may refer to structures such as the innominate artery or innominate bone, depending on context.

\medterm{INOCA} Ischemia with non-obstructive coronary arteries: objective myocardial ischemia or ischemic symptoms in a patient whose coronary arteries do not have an obstructive epicardial stenosis. Mechanisms include coronary microvascular dysfunction and vasospastic angina; evaluation should identify the mechanism so treatment can be tailored.

\medterm{Inoculation} The introduction of microorganisms or clinical material into a culture medium, host, or
sterile site to initiate growth, infection, or immune response.

\medterm{Inoperable} Not safely treatable by surgery, often because a disease is too extensive, located near vital structures, or the person is too medically fragile. Other treatments may still be possible.
\medterm{Inorganic} Not based mainly on carbon-containing molecular structures. Inorganic substances include water, salts, minerals, and metals, many of which are essential to body function while others can be toxic.

\medterm{Inorganic Chemical} A chemical compound that is generally not based on carbon chains, such as water, salts, minerals, or metals. In the body, inorganic substances support fluid balance, nerve signals, bone structure, and many chemical reactions.

\medterm{Inosculation} The joining and establishment of a connection between blood vessels or vascularized tissues, important in wound healing, grafts, flaps, and transplantation.

\medterm{Inosine} A nucleoside made of hypoxanthine linked to ribose that participates in purine metabolism and cellular energy pathways.

\medterm{Inosine Monophosphate} A nucleotide containing inosine and one phosphate group; it is an intermediate in purine metabolism and a precursor of other nucleotides.

\medterm{Inosine Pranobex} A medicine combination marketed in some countries as an immunomodulator or antiviral; evidence and approved uses vary by jurisdiction.

\medterm{Inositol} A sugar-like compound involved in cell membranes and signaling; some forms are studied or used for selected metabolic, reproductive, or mental-health conditions.

\medterm{Inositol Phosphates} Phosphorylated forms of inositol that act in cell signaling, calcium regulation, and other intracellular processes.
\medterm{Inotropes In Anesthesia} Medicines used during anesthesia or critical illness to increase the force of heart contraction and support blood flow when the heart is failing. They are given by controlled infusion and require monitoring because they can cause abnormal heart rhythms, low or high blood pressure, and increased oxygen demand.

\medterm{Inotropic} Affecting the force of cardiac muscle contraction; a positive inotropic drug increases contractility, while a negative inotropic effect decreases it.

\medterm{Inpatient} An inpatient is a person admitted to a healthcare facility for monitoring, diagnosis, treatment, surgery, or care requiring an overnight or formal stay.

\medterm{Inpatient Rehabilitation} Intensive inpatient rehabilitation (3+ hours therapy/day) for acute conditions: stroke,
spinal cord injury, traumatic brain injury, major trauma. See acute rehabilitation.

\medterm{Inquilinus} A genus of gram-negative bacteria associated mainly with water and industrial environments; human infection is rare.

\medterm{Inquilinus Limosus} A species of Inquilinus bacteria rarely isolated from human clinical specimens and occasionally associated with opportunistic infection.

\medterm{INR} International normalized ratio: a standardized measure of the prothrombin time,
calculated as $(\mathrm{patient\ PT}/\mathrm{mean\ normal\ PT})^{\mathrm{ISI}}$, to ensure consistent monitoring of
vitamin K antagonist (warfarin) therapy across laboratories. Target INR: 2-3 for most
indications (AF, VTE); 2.5-3.5 for mechanical heart valves.



\medterm{Insect Bite} A cutaneous reaction resulting from the injection of salivary contents from an insect
during feeding, causing a localized or systemic hypersensitivity response. Common
presentations include a pruritic papule or wheal with surrounding erythema. Severe
reactions may include anaphylaxis, particularly with Hymenoptera stings.

\medterm{Insect Bites} Skin injuries caused by an insect feeding or injecting saliva or venom, producing local inflammation and sometimes allergic, infectious, or systemic reactions.

\medterm{Insect Bites And Stings} Bites and stings from insects such as mosquitoes, fleas, bedbugs, bees, wasps, hornets, yellow jackets, and ants can cause local pain, itching, swelling, allergic reactions, infection, or toxin effects; breathing difficulty or widespread symptoms is an emergency.

\medterm{Insect Sting Allergies} Immune reactions to venom from insects such as bees, wasps, hornets, and fire ants. Reactions may be local or
systemic; generalized hives, breathing difficulty, dizziness, vomiting, or throat swelling
may indicate anaphylaxis and require emergency treatment.

\synonyms
Allergy Insect (Insect Sting Allergies)

\medterm{Insect Virus} A virus that infects insects or other arthropods; some are studied for ecology or biological pest control rather than human disease.

\medterm{Insecta} The biological class containing insects. Some insects bite or sting, cause allergic reactions, contaminate food, or spread infections and parasites; medical importance depends on the species and exposure.

\medterm{Insecticide} A chemical or biological agent used to kill or control insects. Exposure can cause skin, eye, respiratory, neurologic, or cholinergic toxicity depending on the product; prevention and treatment require identifying the active ingredient.

\medterm{Insecticide Poisoning} Toxic effects caused by exposure to an insect-killing chemical; symptoms vary by agent and may include gastrointestinal, neurologic, respiratory, or cholinergic manifestations.

\medterm{Insemination} Introduction of sperm into the female reproductive tract, either through intercourse or an assisted-reproduction procedure such as intrauterine insemination. Indications, success rates, infection screening, and risks depend on the method and patient factors.

\medterm{Insertion} Movement of a muscle attachment toward its fixed attachment, or placement of one structure into another.

\medterm{Insidious} Developing slowly and quietly, often without obvious early symptoms. An insidious illness may be present for a long time before it is recognized, so the term describes its gradual pattern rather than its cause or severity.
\medterm{Insinuate} To introduce or insert gradually into a space, or to suggest something indirectly; in anatomy, a structure may insinuate between or around other structures.

\medterm{Insomnia} See \textit{Sleeplessness}.

\synonyms
Sleeplessness

\medterm{Insomnia Disorder} Difficulty initiating or maintaining sleep, or early-morning awakening with inability to
return to sleep, at least three nights per week for at least three months. The
individual experiences clinically significant distress or impairment in functioning.

\medterm{Inspection} Systematic visual examination of the body or a specimen for abnormalities such as color change, swelling, deformity, lesions, movement, or discharge.
the patient\'s overall appearance, posture, body habitus, gait, skin color, nutritional
status, respiratory effort, and specific anatomical regions. Requires adequate lighting
and proper patient exposure.

\medterm{Inspiration (Inhalation)} The part of breathing in which air enters the lungs. The diaphragm contracts and moves downward, while the chest muscles expand the chest, lowering pressure inside the lungs and drawing air through the airways into the air sacs.

\medterm{Inspiratory Capacity} Inspiratory capacity is the maximum volume inhaled after a normal exhalation and equals tidal volume plus inspiratory reserve volume.

\medterm{Inspiratory Crackles} Fine, discontinuous sounds heard during inspiration, caused by the sudden opening of
collapsed airways. They are a hallmark of interstitial lung disease (bibasilar,
Velcro-like crackles) and pulmonary fibrosis. They are also heard in pneumonia,
atelectasis, and heart failure.

\medterm{Inspiratory Time} The length of time spent breathing in during a natural breath or mechanical-ventilation cycle. Adjusting it can affect airflow, oxygen delivery, comfort, pressure, and the balance between inhaling and exhaling.

\medterm{Instantaneous Wave-Free Ratio} A measurement made during heart catheterization to judge whether a narrowed heart artery limits blood flow. A pressure wire compares pressure beyond and before the narrowing while the heart is relaxed. A result of 0.89 or less may support treatment when it fits the symptoms and other findings.

\synonyms
iFR; Wave-Free Ratio; Resting Coronary Physiology

\medterm{Instillation} Instillation is the gradual introduction of a liquid medicine or solution into a body cavity, organ, or tissue.

\medterm{Instrument} A tool or device used by healthcare workers for examination, measurement, diagnosis, surgery, treatment, patient care, or laboratory work.

\medterm{Instrument Sterilization} A validated process that destroys all forms of microbial life, including bacterial
spores, on or within a medical instrument. Common methods include steam under pressure,
ethylene oxide, hydrogen-peroxide plasma, and dry heat, selected according to device
material and design.

\medterm{Instrumental Activities Of Daily Living} More complex activities required for independent community living, including medication
management, shopping, cooking, finances, communication, transportation, and household
tasks.

\medterm{Instrumental ADL} A complex activity needed for independent community living, such as shopping, cooking, managing medicines or finances, and using transportation.



\medterm{Insufficient Cervix} Painless cervical shortening and dilation during the second trimester, also called
cervical insufficiency. It can cause recurrent mid-trimester pregnancy loss or
spontaneous preterm birth without regular uterine contractions.

\synonyms
cervical insufficiency.

\medterm{Insula} The insula is a cerebral cortical region involved in sensation, taste, emotion, interoception, pain, and integration of internal and external information.

\medterm{Insulin} A hormone made by the pancreas that helps cells use glucose and lowers blood sugar. It is used to treat diabetes, but too much can cause dangerously low blood sugar with sweating, confusion, seizures, or unconsciousness.

\medterm{Insulin And Syringes - Storage And Safety} A patient-care topic concerning insulin and syringes - storage and safety. It addresses the relevant purpose, risks, management, and follow-up.

\medterm{Insulin Antagonist} A hormone, medicine, or other factor that opposes insulin’s effects and can raise blood glucose.

\medterm{Insulin Aspart} A rapid-acting insulin analog (modified amino acid sequence) with onset 10-20 minutes,
peak 1-3 hours, duration 3-5 hours. Used to cover prandial glucose (injected at or just
after meals) in type 1 and type 2 diabetes; also used in continuous subcutaneous insulin
infusion (pump) and IV in acute settings.

\medterm{Insulin Autoantibodies} Autoantibodies directed against insulin or other pancreatic beta-cell antigens,
including GAD, IA-2, and ZnT8 antibodies. Their presence supports autoimmune type 1
diabetes but a negative panel does not exclude it.

\medterm{Insulin C-Peptide Test} A blood or urine test measuring C-peptide released in equal amounts with endogenous insulin; it helps distinguish native insulin production from injected insulin and evaluate hypoglycemia or diabetes classification.

\medterm{Insulin Coma} A state of severe hypoglycemia and unconsciousness caused by excessive insulin; it is a medical emergency and is not a current routine treatment.

\medterm{Insulin Detemir} A long-acting insulin analogue used to provide basal insulin replacement in diabetes mellitus; dosing is individualized and hypoglycemia is the principal serious adverse effect.

\medterm{Insulin Glargine} A long-acting basal insulin analog with a flat, peakless profile providing near-24-hour
coverage, administered once daily (or twice in some patients). Formed by modification of
human insulin producing slow, steady absorption from the injection depot. Used for type
1 and type 2 diabetes (basal component).

\medterm{Insulin Injection Technique} Insulin is injected into subcutaneous tissue using the prescribed device, dose, and site. Common sites include
the abdomen, thighs, buttocks, and upper arms; rotating sites helps reduce lipohypertrophy
and variable absorption. Technique should be individualized to the person and device.

\medterm{Insulin Isophane} Intermediate-acting human insulin, also called neutral protamine Hagedorn insulin, used for basal glucose control in diabetes mellitus.

\medterm{Insulin Lipohypertrophy} Thickened fatty tissue at repeated insulin-injection sites. Injecting into these areas can make
insulin absorption unpredictable.

\medterm{Insulin Lispro} A rapid-acting insulin analogue used around meals or for correction of high blood glucose in diabetes mellitus.

\medterm{Insulin Pen Devices} Pre-filled or refillable devices that deliver measured insulin through a small needle under the skin. They are portable, but the dose, injection technique, storage, and needle disposal must follow medical instructions.

\medterm{Insulin Preparations} Forms of insulin used to lower blood glucose by helping body cells take up glucose and reducing glucose release by the liver. They differ in how quickly they start and how long they work, so the type and dose must be individualized.

\medterm{Insulin Pump Therapy} Delivery of rapid-acting insulin through a small tube under the skin. The pump gives a continuous background dose and extra doses around meals or to correct high blood glucose. Users need training, monitoring, and a backup plan if the device fails.

\medterm{Insulin Pumps} Insulin pumps deliver rapid-acting insulin continuously and through programmed meal or correction doses, allowing flexible management of diabetes while requiring monitoring and device safety.

\medterm{Insulin Receptor} A transmembrane receptor tyrosine kinase that binds insulin and initiates signaling controlling glucose uptake, metabolism, growth, and anabolic processes.

\medterm{Insulin Resistance} Diminished biological response to physiological insulin concentrations resulting in
impaired glucose disposal in peripheral tissues and failure to suppress hepatic glucose
production. Pathophysiological hallmark of type two diabetes and metabolic syndrome.

\synonyms
IR (Insulin Resistance)

\medterm{Insulin Resistance Syndrome} Alternate index label; see \textit{Metabolic syndrome}.

\medterm{Insulin Resistance Syndromes} Rare inherited severe insulin resistance syndromes include Donohue syndrome also known
as leprechaunism, Rabson-Mendenhall syndrome, and type A insulin resistance syndrome
caused by mutations in the insulin receptor gene. Donohue syndrome is the most severe
with death in infancy.

\medterm{Insulin Secretion} Release of insulin from pancreatic beta cells in response primarily to increased plasma
glucose, with amplification by incretin hormones and nutrients. Insulin promotes glucose
uptake, glycogen synthesis, lipogenesis, and protein synthesis while suppressing hepatic
glucose production.

\medterm{Insulin Shock} Severe hypoglycemia, historically associated with excessive insulin and characterized by confusion, seizures, loss of consciousness, or other neurologic dysfunction.

\medterm{Insulin Signaling} The cellular process initiated when insulin binds its receptor and activates pathways that regulate glucose uptake, glycogen formation, protein synthesis, and fat metabolism.

\medterm{Insulin-Dependent Diabetes} Type 1 diabetes, in which the pancreas produces little or no insulin and replacement insulin is required.

\synonyms
Type 1 diabetes

\medterm{Insulin-Induced Lipohypertrophy} A common complication of subcutaneous insulin therapy characterized by localized
accumulation of adipose tissue at injection sites, caused by the lipogenic effect of
insulin on subcutaneous fat. It presents as firm, rubbery, non-tender subcutaneous
nodules that can cause unpredictable insulin absorption and glycemic variability.

\medterm{Insulin-Like Growth Factor} A hormone, mainly IGF-1, made largely by the liver in response to growth hormone. It supports growth of bone and other tissues and helps regulate metabolism. Abnormally high or low levels can occur with pituitary, liver, nutritional, or chronic illnesses.

\medterm{Insulinase} An enzyme that breaks down insulin and helps control how long insulin remains active. Its activity contributes to regulation of blood sugar and may influence insulin requirements or treatment research.

\medterm{Insulinoma} An insulinoma is a usually benign pancreatic beta-cell neuroendocrine tumour that secretes excessive insulin, causing recurrent fasting or exertional hypoglycaemia relieved by glucose.

\medterm{Integrase} Integrase is an enzyme that inserts viral DNA into a host genome, a process targeted by integrase inhibitors used in HIV treatment.

\medterm{Integrative Medicine} Patient-centered care that combines evidence-based conventional treatment with appropriate complementary approaches, emphasizing whole-person health, prevention, lifestyle, and shared decision-making.

\medterm{Integrative Rehabilitation} Combination of conventional and complementary therapies: acupuncture, massage, yoga,
meditation.

\medterm{Integrin} A cell-surface receptor that links cells to the extracellular matrix and helps control adhesion, movement, signaling, and immune-cell activity.

\medterm{Integrin Binding} The attachment of a molecule or cell to an integrin receptor, which can influence adhesion, migration, signaling, or clot formation.

\medterm{Integron} A genetic element containing an integrase and an attachment site that can capture and express gene cassettes; integrons commonly contribute to bacterial antimicrobial resistance.

\medterm{Integrons} Genetic elements that capture and express gene cassettes through site-specific
recombination. Integrons contain an integrase gene and a recombination site for gene
cassette insertion. They are important vehicles for the spread of antibiotic resistance
genes, particularly in gram-negative bacteria.

\medterm{Integumentary System} The body system comprising the skin and its appendages: hair, nails, sweat glands, and
sebaceous glands. It forms a protective barrier, limits fluid loss, contributes to
thermoregulation and sensation, supports immune defence, and participates in vitamin D
synthesis.

\medterm{Intein} An intein is a protein segment that excises itself from a precursor protein while joining the surrounding extein sequences.

\medterm{Intellectual Disability} Neurodevelopmental disorder with onset in the developmental period, characterised by
deficits in intellectual and adaptive functioning across conceptual, social, and
practical domains. DSM-5 grades severity by adaptive functioning: mild, moderate,
severe, profound.

\medterm{Intelligence} Intelligence is the capacity to learn, reason, solve problems, understand information, and adapt; tests measure selected abilities rather than a person’s entire worth or potential.

\medterm{Intelligence Test} An intelligence test uses standardized tasks to estimate cognitive abilities such as reasoning, memory, language, and problem-solving.

\medterm{Intense Pulsed Light} A treatment using brief pulses of broad-spectrum light for selected skin conditions and hair reduction. Light is absorbed by pigment or blood vessels, and treatment can cause temporary redness, burns, pigment changes, or eye injury.

\synonyms
IPL

\medterm{Intensity} The degree of effort or physiological demand during physical activity.

\medterm{Intensity-Modulated Radiation Therapy} A computer-optimized form of three-dimensional radiation therapy that shapes and
modulates the radiation dose to conform tightly to the target while sparing adjacent
normal tissue.

\medterm{Intensive Care} A patient-care topic concerning intensive care. It addresses the relevant purpose, risks, management, and follow-up.

\medterm{Intensive Care Rehabilitation} Early rehabilitation in ICU setting. Early mobilization, positioning, passive range of
motion. Improves outcomes, reduces complications.

\medterm{Intensive Care Unit} A specialist hospital unit providing continuous monitoring and life support for
critically ill patients. In obstetrics, women may require ICU admission for: severe
pre-eclampsia or eclampsia, massive postpartum hemorrhage, amniotic fluid embolism,
sepsis, cardiac disease in pregnancy, or complications of anesthesia.

\medterm{Intention Tremor} A tremor that becomes more pronounced during a purposeful movement toward a target. It commonly suggests cerebellar or cerebellar-pathway dysfunction, although medicines, toxins, and other neurologic disorders can contribute.

\medterm{Interatrial Septum} The wall between the heart's right and left upper chambers. It includes the fossa ovalis, a normal remnant of fetal circulation, and contains a thinner area that may be crossed during selected catheter procedures.

\medterm{Intercalated Disc} An intercalated disc is a specialized junction connecting cardiac muscle cells and transmitting electrical and mechanical force through the heart.

\medterm{Intercalating Agent} An intercalating agent inserts between DNA base pairs and can disrupt replication or transcription; some are laboratory reagents or anticancer drugs.

\medterm{Intercellular} Located between cells or occurring in the space between cells. Intercellular material and signals help cells attach, communicate, exchange substances, and organize tissues throughout the body.

\medterm{Intercostal} Intercostal means located between the ribs, especially the muscles, nerves, and blood vessels in those spaces.

\medterm{Intercostal Muscles} Muscles between the ribs that help expand and compress the chest during breathing and also assist movement of the trunk.

\medterm{Intercostal Retraction} Inward movement of the skin and soft tissues between the ribs during inspiration,
indicating increased work of breathing. It is common in infants and children with
respiratory distress and may accompany airway obstruction, bronchiolitis, pneumonia,
asthma, or other pulmonary disease.

\medterm{Intercostal Retractions} Inward pulling of the skin and soft tissues between the ribs during inspiration, indicating increased work of breathing and potentially significant respiratory distress.

\medterm{Intercurrent Disease} An illness or complication that develops during the course of another disease or treatment and may affect its management or outcome.

\medterm{Interdental Brush} Small brush for cleaning between teeth. Used where spaces exist.

\medterm{Interdigital Neuroma} Alternate index label; see \textit{Morton neuroma}.

\medterm{Interference} A factor that alters the accuracy or interpretation of a laboratory test, measurement, or drug effect.

\medterm{Interferometry} A measurement method that uses interference between waves to determine distance, motion, thickness, optical properties, or other physical quantities.

\medterm{Interferon} Interferons are cytokines that induce antiviral defenses, alter immune-cell activity, and regulate cell growth; type I, II, and III interferons have different receptors and clinical uses.

\medterm{Interferon Inducer} A substance that stimulates cells to produce interferons, proteins involved in antiviral and immune responses.
\medterm{Interferon Signaling} Type I IFN (IFN-alpha/beta): all nucleated cells produce; receptor IFNAR (ubiquitous);
JAK1/TYK2 -> STAT1/2/IRF9 (ISGF3) -> ISGs (hundreds: antiviral, antiproliferative,
immunomodulatory). Type II IFN (IFN-gamma): T/NK cells; receptor IFNGR; JAK1/JAK2 ->
STAT1 -> GAS elements. Type III IFN (IFN-lambda): mucosal epithelium; receptor IFNLR.

\medterm{Interferon-Alpha} A type I interferon cytokine with antiviral, antiproliferative, and immune-modulating effects; recombinant forms have been used for selected viral infections and cancers.

\medterm{Interferon-Beta} A type I interferon cytokine used mainly as an immunomodulatory treatment for relapsing forms of multiple sclerosis.

\medterm{Interferon-Gamma Release Assay} Blood test for Mycobacterium tuberculosis infection (latent or active). T-SPOT.TB(ELISPOT) or QuantiFERON-TB Gold Plus. Measures IFN-gamma released by T cells
after exposure to MTB-specific antigens (ESAT-6, CFP-10, TB7.7).

\synonyms
Interferon-Gamma Release Assay

\medterm{Interferonopathies} A group of rare disorders caused by abnormal activation of type I interferon pathways. They can produce inflammation,
neurologic disease, skin lesions, blood-vessel injury, or lung disease, often beginning
in childhood.

\medterm{Interferons} Natural proteins that cells release to warn nearby cells and help control the immune response. They are especially important in the body's defense against viruses and are also used as medicines for some infections, cancers, and immune conditions.

\medterm{Interleukin} A family of cytokines (interleukins) produced mainly by leukocytes that regulate immune
cell growth, differentiation, activation, and communication. Examples: IL-1 and IL-6
(inflammation, fever), IL-2 (T-cell growth), IL-4/IL-13 (Th2, IgE), IL-5 (eosinophils),
IL-10 (anti-inflammatory), IL-12 (Th1), IL-17 (Th17, autoimmunity), IL-23.

\medterm{Interleukin 1} A group of proinflammatory cytokines, especially IL-1 alpha and IL-1 beta, that promote fever, leukocyte activation, and inflammatory responses.

\medterm{Interleukin 13} A cytokine produced by type 2 helper T cells and other immune cells that contributes to allergic inflammation, mucus production, and tissue fibrosis.

\medterm{Interleukin 2} Interleukin-2 is a T-cell growth factor that promotes proliferation and survival of activated T cells and natural-killer cells; high-dose therapy can cause capillary-leak syndrome and systemic toxicity.

\medterm{Interleukin 7} A cytokine important for development, survival, and homeostasis of T lymphocytes and some B-cell precursors.
\medterm{Interleukin-6} A cytokine involved in immune regulation and inflammation that stimulates production of acute-phase proteins by the liver.

\medterm{Interleukins} Cytokines that carry signals between immune cells and regulate inflammation.

\synonyms
ILs

\medterm{Interlocutory} Interlocutory means provisional or occurring during an ongoing process, such as an interim clinical, legal, or research decision.

\medterm{INTERMACS Profile} A seven-level system for describing how sick a person is with advanced heart failure. Profile 1 means critical shock with very poor blood flow; profiles 2 and 3 describe worsening or dependence on medicines that strengthen the heart. Profiles 4 through 7 describe repeated hospital stays or increasing limits on activity. The profile helps doctors judge urgency and timing for a heart pump or transplant.

\synonyms
INTERMACS Staging; INTERMACS Scale

\medterm{Intermediate} Intermediate means between two stages, positions, or levels, such as an intermediate metabolite, cell, or developmental phase.

\medterm{Intermediate Cuneiform} Smallest cuneiform, centrally located between medial and lateral cuneiforms. Articulates
with navicular, medial/lateral cuneiforms, and second metatarsal. Forms the transverse
arch apex. Relatively fixed, contributing to midfoot stability. Injuries are typically
associated with high-energy midfoot trauma.

\medterm{Intermediate Filament} An intermediate filament is a strong cytoskeletal fiber that supports cell shape, mechanical stability, and tissue integrity.

\medterm{Intermediate Uveitis} Inflammation mainly affecting the gel inside the eye and nearby tissues at the edge of the retina. It may cause blurred vision, floaters, or little pain and can be linked to autoimmune disease, infection, or no identified cause. Eye examination and imaging guide treatment and monitoring for complications.

\medterm{Intermediate-Density Lipoprotein} A temporary particle formed when the body removes some fat from very-low-density lipoprotein. It can be changed into LDL, the type of cholesterol particle that can contribute to fatty buildup in artery walls.

\synonyms
IDL

\medterm{Intermittent Auscultation} Periodic assessment of the fetal heart rate with a fetoscope or Doppler device
during labor rather than continuous electronic monitoring. The frequency of assessment depends on the stage of labor,
risk factors, and local clinical protocol.

\medterm{Intermittent Claudication} Reproducible muscle discomfort, usually in the calf, thigh, or buttock, induced by
walking and relieved by rest, caused by exertional arterial insufficiency. It is
distinguished from rest pain and nonvascular musculoskeletal pain by history and
vascular testing.

\medterm{Intermittent Explosive Disorder} A disorder characterized by recurrent, impulsive, problematic outbursts of verbal or
physical aggression that are grossly disproportionate to the provocation. The outbursts
are not premeditated and cause marked distress or functional impairment. Age >6 years.

\synonyms
Road Rage

\medterm{Intermittent Fasting} An eating pattern that alternates planned periods of eating with periods of little or no food. Common schedules limit daily eating hours or use selected fasting days. It may help some adults lose weight, but it is not suitable for everyone, especially during pregnancy or with some medical conditions or eating disorders.

\medterm{Intermittent Fever} A fever that rises and falls with periods when temperature returns to normal or near normal.

\medterm{Intern} A physician in the early stage of postgraduate medical training, traditionally the first year after medical school; the exact training structure varies by country.

\medterm{Internal} Internal means located within the body, an organ, a structure, or a system rather than on its outer surface.

\medterm{Internal Bleeding} Internal bleeding is blood loss into a body cavity, tissue, or organ rather than through an obvious external wound. Trauma, ulcers, aneurysm, pregnancy complications, medicines, and tumors can cause it; faintness, shock, abdominal swelling, black stool, or vomiting blood is an emergency.

\medterm{Internal Capsule} A compact band of white-matter nerve fibres in the brain carrying important motor and sensory pathways between the cerebral cortex and lower centres.

\medterm{Internal Environment} The fluid and physiologic conditions inside the body that cells require, including temperature, pH, oxygen, electrolytes, and nutrients.

\medterm{Internal Fixation} Surgical stabilization of a fracture or unstable bone using implanted hardware such as plates, screws, rods, wires, or pins.

\medterm{Internal Fixator} An internal fixator is a surgically implanted plate, screw, rod, nail, or frame used to stabilize a fracture or correct skeletal alignment.

\medterm{Internal Hernia} Herniation of bowel through a mesenteric defect common after Roux-en-Y gastric bypass
creating a defect at the jejunojejunostomy. May cause intermittent or acute small bowel
obstruction. Presents with episodic abdominal pain, nausea, and vomiting.

\medterm{Internal Iliac Artery} Smaller common iliac terminal branch (hypogastric), supplying the pelvis, perineum,
gluteal region, and medial thigh. Posterior division: iliolumbar, lateral sacral,
superior gluteal. Anterior division: inferior gluteal, internal pudendal, obturator,
vesical, uterine, vaginal, middle rectal arteries.

\medterm{Internal Jugular Vein} Descends in the carotid sheath with the internal carotid artery and vagus nerve,
draining brain, face, and neck. Begins at the jugular foramen as a sigmoid sinus
continuation, joining the subclavian vein to form the brachiocephalic vein. Contains
backflow-preventing valves. Used for central venous catheterization.

\medterm{Internal Oblique Muscle} Middle flat abdominal layer deep to external oblique, from thoracolumbar fascia, iliac
crest, and lateral inguinal ligament to lower ribs, linea alba, and pubic crest.
Innervated by T7-T12 and L1. Ipsilateral rotation (opposite of external oblique).
Contributes to inguinal canal and abdominal wall integrity.

\medterm{Internal Podalic Version} A maneuver where the operator inserts a hand into the uterus, grasps one or both fetal
feet, and converts a transverse lie or cephalic presentation to a breech, then delivers
by total breech extraction. Indication: delivery of the second twin (transverse lie) or
severe fetal distress in the second stage. Rare in modern obstetrics.

\medterm{Internal Resorption} Destruction of tooth tissue from the inside of the tooth outward, caused by activity of specialized cells. It may produce a pink area or no symptoms and can weaken the tooth; dental imaging determines whether root-canal treatment can save it.

\medterm{International Normalized Ratio} The INR, a standardized calculation from the prothrombin time used to assess blood clotting and guide treatment with vitamin K antagonists such as warfarin.

\medterm{International Unit} A unit that measures the biological effect of a substance, such as insulin, a vitamin, or a vaccine. The amount represented by one unit differs between substances, so international units cannot be directly converted to milligrams without knowing the substance.

\medterm{Interneuron} An interneuron connects neurons within the central nervous system and helps integrate, inhibit, or coordinate neural signals.

\medterm{Interneurons} Neurons within the central nervous system that connect other neurons and help process or coordinate neural signals.

\medterm{Internist} An internist is a physician specializing in adult internal medicine, including diagnosis and management of complex medical conditions.

\medterm{Internuclear Ophthalmoplegia} A conjugate horizontal gaze disorder caused by a lesion of the medial longitudinal
fasciculus (MLF), a heavily myelinated fiber tract that connects the abducens nucleus
(CN VI) on one side to the oculomotor nucleus (CN III) on the contralateral side,
coordinating horizontal gaze.

\medterm{Interosseous} Interosseous means located between bones, such as the membrane or muscles between the forearm or lower-leg bones.

\medterm{Interpersonal Psychotherapy} A time-limited, evidence-based psychotherapy for depression that focuses on the
interpersonal context of mood symptoms. It identifies four problem areas: grief, role
disputes, role transitions, and interpersonal deficits. IPT helps patients improve
interpersonal functioning and social support.

\medterm{Interphase} Interphase is the portion of the cell cycle between divisions, including growth, DNA replication, and preparation for mitosis or meiosis.

\medterm{Interphase Cell} A cell between divisions that grows, performs normal functions, and copies its DNA before possible division.

\medterm{Interproximal Caries} Dental caries occurring on the contact surfaces between adjacent teeth. It may be
difficult to detect clinically and is commonly assessed with bitewing radiographs when
imaging is indicated.

\medterm{Intersexual} An older term for a person born with physical sex characteristics that do not fit typical definitions of male or female. Modern care uses “difference of sex development” and respects the person’s identity and choices.

\medterm{Interstitial} Relating to the spaces between cells or tissues, such as interstitial fluid or interstitial lung disease.

\medterm{Interstitial (Cornual) Pregnancy} An ectopic pregnancy implanted in the part of the fallopian tube that passes
through the muscular wall of the uterus. It can enlarge before rupture and cause severe internal bleeding.

\synonyms
Cornual

\medterm{Interstitial Cystitis} Long-lasting bladder or pelvic pain with frequent or urgent urination, often without infection. Symptoms may flare and improve; evaluation rules out other causes, and treatment may include bladder training, physical therapy, lifestyle changes, or selected medicines.

\synonyms
Bladder Inflammation
Cystitis, Interstitial
IC (Interstitial Cystitis)
Painful Bladder Syndrome

\medterm{Interstitial Fluid} Interstitial fluid is the fluid surrounding tissue cells, formed from blood plasma and returned to circulation through lymphatic vessels.

\medterm{Interstitial Granulomatous Dermatitis} A rare inflammatory dermatosis with interstitial or palisading histiocytes around collagen, often presenting with symmetric erythematous or violaceous plaques and associated systemic disease or medication exposure.

\medterm{Interstitial Keratitis} Inflammation of the corneal stroma without primary involvement of the corneal surface.
It may follow infection, autoimmune disease, or trauma and can cause pain, photophobia,
redness, and reduced vision.

\medterm{Interstitial Lung Disease} A heterogeneous group of disorders with inflammation and fibrosis of lung parenchyma,
causing restrictive physiology (reduced FVC and TLC with normal or elevated FEV1/FVC
ratio).

\synonyms
Diffuse Parenchymal Lung Diseases
Interstitial Lung Diseases
Lung Disease, Interstitial

\medterm{Interstitial Lung Disease (Interstitial Pneumonia)} Alternate index label for Interstitial Pneumonia.



\medterm{Interstitial Nephritis} Inflammation of the renal interstitium and tubules, which can be acute or chronic. Acute
interstitial nephritis is most commonly caused by medications, particularly NSAIDs,
antibiotics including penicillins and cephalosporins, and proton pump inhibitors.

\medterm{Interstitial Pneumonia} Inflammation affecting the supporting tissue around the lung air sacs. It may be caused by infection, autoimmune disease, medicines, or other lung disorders and can cause dry cough, fever, breathlessness, and low oxygen. Examination and imaging help determine the cause and treatment.

\synonyms
Interstitial Lung Disease (Interstitial Pneumonia)

\medterm{Interstitium} The supporting tissue between and around the lung air sacs, capillaries, and small airways; inflammation or fibrosis of this tissue can impair lung expansion and gas exchange.

\medterm{Intertrigo} Intertrigo is inflammation in opposing skin folds caused by moisture, heat, friction, and maceration; Candida or bacterial infection may complicate the tender erythematous rash.

\synonyms
Allergy Diaper (Intertrigo)

\medterm{Interval} An interval is the time, distance, or range between two points or events, such as a dosing or follow-up interval.

\medterm{Intervening Sequence} A nucleotide sequence within a gene transcript that lies between coding sequences; in many genes it is an intron removed during RNA processing.

\medterm{Intervention} An action intended to prevent, diagnose, treat, or change a health condition or outcome, such as a medicine, procedure, counseling, or public-health measure.

\medterm{Intervention Study} A study in which researchers assign or provide an intervention and compare its outcomes with a control.

\synonyms
Clinical trial

\medterm{Interventional} Relating to a procedure or active treatment used to diagnose or manage disease, often through minimally invasive access rather than open surgery.

\medterm{Interventional Cardiology} A subspecialty using catheter-based procedures to diagnose and treat coronary, structural, and selected vascular heart conditions, including angioplasty, stenting, and transcatheter valve procedures.

\medterm{Interventional Procedure} A targeted procedure used to diagnose or treat disease, often through a needle, tube, or small incision and guided by imaging. Examples include biopsy, drainage, vessel treatment, and tumor ablation; risks depend on the body area and technique.

\medterm{Interventional Procedures} Procedures that diagnose or treat disease through a needle, catheter, or small incision, often using imaging for guidance. Examples include biopsy, drainage, vessel opening, embolization, and ablation; less invasive does not mean risk-free.

\medterm{Interventional Radiologist} A physician trained to perform minimally invasive diagnostic and therapeutic procedures using medical imaging for guidance.

\medterm{Interventional Radiology} A specialty using imaging guidance to perform minimally invasive diagnostic and therapeutic procedures, such as biopsies, drainages, embolization, and vascular treatments.

\medterm{Interventricular Septum} The thick muscular wall separating the left and right ventricles. It has two portions:
the muscular (inferior, thicker) and membranous (superior, thinner) portions. The septum
contains parts of the cardiac conduction system, including the bundle of His.

\medterm{Intervertebral Disc} A fibrocartilaginous cushion between adjacent vertebral bodies, composed of a gelatinous
nucleus pulposus surrounded by a tough annulus fibrosus. It absorbs shock and permits
limited movement of the vertebral column.

\medterm{Intervertebral Disk} A strong, flexible cushion between neighboring spinal bones that absorbs shock and allows movement. A damaged or herniated disk can press on nearby nerves and cause back pain, numbness, tingling, or weakness.

\synonyms
Intervertebral disc

\medterm{Intervertebral Foramen} An opening between two neighboring bones of the spine through which a spinal nerve leaves the spinal canal. Narrowing of this opening can press on the nerve and cause pain, numbness, tingling, or weakness in its area of the body.

\synonyms
Neural foramen

\medterm{Intestinal Absorption} Intestinal absorption is the movement of digested nutrients, water, and electrolytes from the intestine into the blood or lymph.

\medterm{Intestinal Atresia} A congenital obstruction of the small intestine (jejunum or ileum) caused by an
intrauterine vascular accident (mesenteric ischemia) leading to aseptic necrosis and
resorption of a segment of bowel. In contrast to duodenal atresia, it is not associated
with Down syndrome.

\medterm{Intestinal Cancer} Intestinal cancer is a malignant growth arising in the small or large intestine; symptoms and treatment depend on its location and stage.

\medterm{Intestinal Duplication} Intestinal duplication is a congenital extra segment or cyst of bowel that may cause obstruction, bleeding, pain, or infection.

\medterm{Intestinal Fistula} An intestinal fistula is an abnormal passage connecting the intestine to another organ or the skin, often after inflammation, surgery, or injury.

\medterm{Intestinal Ganglioneuromatosis} Benign proliferation of ganglion cells and nerve fibers in the gastrointestinal tract
wall. May be diffuse involving the entire gastrointestinal tract or localized as
ganglioneuromas. Associated with multiple endocrine neoplasia type 2B where intestinal
ganglioneuromatosis may precede the development of medullary thyroid carcinoma.

\medterm{Intestinal Gas} Gas within the digestive tract that may cause belching, bloating, abdominal discomfort, passing wind, or symptoms from an underlying digestive disorder.

\medterm{Intestinal Gas (Belching, Bloating, Flatulence)} Alternate index label for Belching, Bloating, Flatulence.

\medterm{Intestinal Gas (Intestinal Gas (Belching, Bloating, Flatulence))} Alternate index label for Intestinal Gas (Belching, Bloating, Flatulence).

\medterm{Intestinal Ischemia} Too little blood flow reaching the intestines. It may develop suddenly from a clot or blockage, or gradually from narrowed arteries, and can cause severe pain, bloody stools, or bowel damage; sudden severe pain needs emergency care.

\synonyms
Ischemia, Intestinal

\medterm{Intestinal Leiomyoma} A usually benign smooth-muscle tumor arising in the wall of the intestine; symptoms depend on size and location and may include bleeding, obstruction, or abdominal pain.

\medterm{Intestinal Lipodystrophy} Alternate index label; see \textit{Whipple's disease}.

\medterm{Intestinal Lipomatosis} Benign overgrowth of adipose tissue within the intestinal wall, most commonly affecting
the ileum and colon. Often asymptomatic and discovered incidentally, but large lipomas
can cause intussusception, gastrointestinal bleeding, or bowel obstruction.

\medterm{Intestinal Lymphoma} Malignant proliferation of lymphoid cells arising in the small or large intestine.
Important types include enteropathy-associated T-cell lymphoma, monomorphic
epitheliotropic intestinal T-cell lymphoma, and diffuse large B-cell lymphoma.

\medterm{Intestinal Malrotation} Congenital failure of normal midgut rotation and fixation, predisposing to a narrow
mesenteric base, volvulus, and duodenal obstruction. Infants may present with bilious
vomiting, abdominal distension, or bowel ischemia; sudden symptoms require emergency
assessment.

\medterm{Intestinal Metaplasia} Replacement of the normal lining of an intestinal or gastric region by intestinal-type epithelium, usually in response to chronic injury or inflammation and sometimes associated with increased cancer risk.

\medterm{Intestinal Mucosa} The intestinal mucosa is the inner lining of the intestine, where digestion, nutrient absorption, mucus production, and immune defence occur.

\medterm{Intestinal Obstruction} A mechanical or functional blockage that prevents normal passage of intestinal contents.
Causes include adhesions, hernias, tumors, inflammation, volvulus, and disorders of intestinal motility.

\medterm{Intestinal Obstruction And Ileus} Impaired passage of intestinal contents caused by a mechanical blockage or by reduced bowel motility without a physical blockage; symptoms commonly include abdominal pain, distension, vomiting, and inability to pass stool or gas.

\medterm{Intestinal Obstruction Repair} Surgery or another intervention to relieve an intestinal blockage, remove a diseased segment when necessary, and restore bowel continuity or create a stoma when required.


\medterm{Intestinal Parasite} A protozoan or worm that lives in the gastrointestinal tract and may cause diarrhea, abdominal pain, anemia, malnutrition, or no symptoms.

\medterm{Intestinal Perforation} Full-thickness breach in the bowel wall allowing luminal contents to enter the
peritoneal cavity. Causes include diverticulitis, appendicitis, peptic ulcer disease,
and trauma. Presents with acute abdomen, peritonitis, and free air on imaging.

\medterm{Intestinal Polyp} An intestinal polyp is a growth projecting from the lining of the intestine; some types can become cancerous.

\medterm{Intestinal Polyposis} Intestinal polyposis is the presence of numerous polyps in the intestine, sometimes as part of an inherited cancer-predisposition syndrome.

\medterm{Intestinal Pseudo-Obstruction} Acute or chronic dilation of the bowel without mechanical obstruction due to
neuromuscular dysfunction. Acute colonic pseudo-obstruction also called Ogilvie syndrome
occurs in hospitalized or post-surgical patients.

\synonyms
Ogilvie syndrome.

\medterm{Intestinal Villi} Intestinal villi are tiny finger-like projections of the small-intestinal lining that increase surface area for nutrient absorption.

\medterm{Intestine} The part of the digestive tract extending from the stomach to the anus, consisting of the small intestine and large intestine.

\medterm{Intestine Disorder} An intestine disorder is a disease affecting the small or large intestine, such as infection, inflammation, obstruction, or cancer.

\medterm{Intestine Neoplasm} An intestine neoplasm is an abnormal growth in the small or large intestine; it may be benign or malignant.

\medterm{Intestine Obstruction} Intestinal obstruction is partial or complete blockage of the passage of intestinal contents, causing pain, swelling, vomiting, and failure to pass stool or gas.

\medterm{Intestine, Diseases Of} Diseases of the small or large intestine include infection, inflammation, poor absorption, blockage, polyps, and cancer. Symptoms may include pain, diarrhea, constipation, bleeding, vomiting, weight loss, or nutritional deficiency.

\medterm{Intimal Dissection} A tear in the inner lining of an artery that allows blood to enter and split the vessel wall, creating a false lumen. Aortic or other arterial dissection can compromise blood flow and is a medical emergency.

\medterm{Intimal Hyperplasia} Proliferation and migration of vascular smooth-muscle cells with extracellular-matrix
deposition within the intima after vascular injury or graft placement. It can produce
restenosis at anastomoses, stents, dialysis access, or bypass grafts.

\medterm{Intimate Partner Violence In Pregnancy} Physical, sexual, or psychological abuse by a current or former
partner during pregnancy. It can harm the pregnant person and fetus and is associated with complications such as injury,
preterm birth, and placental abruption.

\medterm{Intolerance} An adverse response to a substance that does not involve the immune mechanism of an allergy.

\synonyms
Nonallergic sensitivity
Nonceliac Gluten Sensitivity (Intolerance)

\medterm{Intolerance, Gluten (Nonceliac Gluten Sensitivity (Intolerance))} Alternate index label for Nonceliac Gluten Sensitivity (Intolerance).

\medterm{Intoplicine} Intoplicine is an investigational topoisomerase inhibitor studied for anticancer effects by disrupting DNA structure and replication.

\medterm{Intoxication} A state of physical or mental impairment caused by recent exposure to a substance, such as alcohol, a medicine, or another drug.

\medterm{Intra-Abdominal Abscess} Localized collection of pus within the abdomen commonly in the pelvis, between bowel
loops, or beneath the diaphragm. Results from intra-abdominal infection or surgery.
Presents with fever, abdominal pain, and leukocytosis. Diagnosed by CT with contrast
showing rim-enhancing fluid collection.

\medterm{Intra-Abdominal Hypertension} Sustained elevation of intra-abdominal pressure above twelve millimeters of mercury
measured through the urinary bladder. Graded as grade one twelve to fifteen, grade two
sixteen to twenty, grade three twenty-one to twenty-five, and grade four greater than
twenty-five millimeters of mercury.

\medterm{Intra-Aortic Balloon Pump} Mechanical circulatory support. Inflatable catheter in descending aorta. Inflates during
diastole (increases coronary perfusion), deflates during systole (reduces afterload).
Indications: cardiogenic shock, unstable angina, high-risk PCI. Complications: limb
ischemia, bleeding, infection.

\medterm{Intra-Arterial} Intra-arterial means within an artery, including a catheter, injection, infusion, or measurement performed directly in an arterial vessel.

\medterm{Intraabdominal Cancer} A broad term for malignancies arising within the abdominal cavity, including cancers of
the colon, stomach, pancreas, liver, ovaries, and peritoneum. Presentation varies widely
depending on the specific organ involved.

\medterm{Intraabdominal Hemorrhage} Bleeding inside the abdominal cavity, which may follow trauma, surgery, organ rupture, pregnancy complications, or internal disease.

\medterm{Intraabdominal Hernia} Herniation of abdominal viscera through a congenital or acquired defect or abnormal
mesenteric aperture within the peritoneal cavity, distinct from ventral or inguinal
hernias. Types include paraduodenal, foramen of Winslow, transmesocolic, and
transomental hernias.

\medterm{Intraabdominal Infection} An infection within the abdominal cavity, such as peritonitis or an abscess, often requiring antibiotics and sometimes drainage or surgery.

\medterm{Intraaortic Balloon Pump} A temporary circulatory-support device with a balloon in the descending aorta that inflates and deflates in coordination with the cardiac cycle.

\medterm{Intraarticular} Intraarticular means within a joint, such as an injection, fracture, loose body, or inflammatory process inside the joint space.

\medterm{Intrabronchial} Intrabronchial means within a bronchus, one of the air passages branching from the trachea into the lungs.

\medterm{Intracardiac Electrophysiology Study} An invasive test using electrode catheters inside the heart to record electrical activity and provoke or map arrhythmias, helping guide diagnosis and treatment.

\medterm{Intracartilaginous} Located within cartilage, such as a lesion, injection, or developing structure. The term identifies position and does not by itself indicate whether the finding is normal, injured, inflamed, or cancerous.

\medterm{Intracavernosal Injection} An injection of a prescription medicine into the erectile tissue of the penis to treat erectile dysfunction. The medicine relaxes muscle and widens blood vessels, producing an erection that does not require sexual stimulation. Pain, bruising, scar tissue, or an erection lasting too long can occur; a prolonged painful erection needs urgent treatment.

\synonyms
ICI; Penile Injection Therapy; Intracavernous Injection

\medterm{Intracavitary} Intracavitary means within a body cavity, such as a treatment, device, fluid, or lesion located in the thorax, abdomen, uterus, or another cavity.

\medterm{Intracavity Radiotherapy} Intracavity radiotherapy places a radiation source inside a body cavity near a tumor, commonly to treat selected cancers of the cervix, uterus, or vagina.

\medterm{Intracellular} Inside a cell or occurring within a cell. Intracellular fluid contains most of the body’s potassium and magnesium, whereas sodium is found mainly outside cells. The cell membrane controls movement of water and dissolved substances between the cell and its surroundings.

\medterm{Intracellular Fluid} The fluid contained within cells, comprising approximately two thirds of total body
water in a typical adult. Its major solutes include potassium, magnesium, phosphate, and
intracellular proteins, and its composition is maintained by selective membrane
transport.

\medterm{Intracellular Membrane} A membrane inside a cell that separates compartments and controls movement of substances between them.

\medterm{Intracellular Pathogens} Microorganisms that survive and replicate within host cells. Intracellular bacteria
include Mycobacterium, Listeria, and Chlamydia. Viruses are obligate intracellular
pathogens. Intracellular survival protects pathogens from antibiotics and immune
responses.

\medterm{Intracellular Transport} Movement of proteins and other materials within a cell, often in membrane-bound sacs along the cell's internal framework.

\medterm{Intracerebral} Intracerebral means within the brain tissue, as in an intracerebral hemorrhage, injection, electrode, or lesion.
\medterm{Intracerebral Hemorrhage} Bleeding directly into the brain parenchyma, accounting for approximately ten to fifteen
percent of all strokes and carrying a high mortality rate of approximately forty percent
at thirty days.

\medterm{Intracisternal} Intracisternal means within a cistern, especially a cerebrospinal-fluid space at the base of the brain.
\medterm{Intracoronal} Intracoronal means within the crown of a tooth or within a coronary structure, with the exact meaning determined by context.

\medterm{Intracoronary} Intracoronary means within a coronary artery or the coronary circulation, such as a catheter, stent, imaging study, or medicine delivery.

\medterm{Intracranial} Inside the skull. The term may describe the brain, blood vessels, fluid, bleeding, pressure, infection, or a growth within the skull. Extra fluid, blood, or tissue can raise pressure and damage the brain, causing headache, vomiting, confusion, weakness, seizures, or loss of consciousness.

\medterm{Intracranial Aneurysm} A weakened, bulging area in a blood vessel inside the skull, usually in an artery supplying the brain. Many cause no symptoms, but a leak or rupture can cause bleeding around the brain and a sudden, severe headache.

\medterm{Intracranial Hemorrhage} Bleeding within the cranial cavity. Epidural hematoma: arterial bleeding (middle
meningeal artery), typically from temporal bone fracture; lucid interval then rapid
deterioration. Subdural hematoma: venous bleeding from bridging veins; acute (<72h),
subacute, or chronic (elderly, alcoholics).

\medterm{Intracranial Pressure} The pressure within the cranial vault generated by brain tissue, cerebrospinal fluid,
and intracranial blood. Sustained elevation can reduce cerebral perfusion and cause
headache, vomiting, altered consciousness, pupillary abnormalities, and brain
herniation.

\medterm{Intracranial Pressure Monitoring} Continuous or intermittent measurement of pressure inside the skull, usually in critically ill patients with severe brain injury, hemorrhage, hydrocephalus, or other risk of intracranial hypertension.

\medterm{Intracranial Venous Malformations} Abnormal collections or connections of veins within the skull that may bleed or affect nearby brain tissue.

\medterm{Intractable} Difficult to control, relieve, or cure despite appropriate treatment, as in intractable pain or seizures.
\medterm{Intradermal} Located within the dermis, the layer of skin beneath the epidermis; an intradermal injection places medication or test material into that layer.

\medterm{Intraductal Carcinoma} Alternate index label; see \textit{Ductal carcinoma in situ}.

\medterm{Intraductal Papillary Mucinous Nneoplasm Pancreas (Pancreatic Cysts)} Alternate index label for Pancreatic Cysts.

\medterm{Intraductal Papilloma} A benign papillary tumor arising within a breast duct, often presenting with spontaneous
unilateral bloody or clear nipple discharge. It may be solitary near the nipple or
multiple in the peripheral breast.

\medterm{Intraduodenal} Intraduodenal means within the duodenum, the first part of the small intestine beyond the stomach.
\medterm{Intraepidermal} Intraepidermal means within the epidermis, the outermost layer of skin.

\medterm{Intraepithelial} Located within the layer of cells that covers a body surface or lines an organ. The word may describe infection, inflammation, or abnormal cell changes that have not yet grown through the supporting layer beneath them.

\medterm{Intraesophageal} Intraesophageal means within the esophagus, such as a probe, pressure sensor, lesion, or medication delivery.
\medterm{Intrahepatic} Within the liver. For example, intrahepatic bile-duct blockage or inflammation can cause jaundice and abnormal liver tests, while a tumor may begin in liver cells or in bile ducts inside the liver. The exact meaning depends on the condition being described.

\medterm{Intrahepatic Cholestasis} Reduced bile flow within the liver causing bile substances to build up in blood, often producing itching, jaundice, dark urine, or pale stools.

\medterm{Intrahepatic Cholestasis Of Pregnancy} A pregnancy-associated liver disorder characterized mainly by
itching without a primary rash and increased bile acids in the blood. It may increase fetal risk and requires clinical
evaluation and monitoring.

\medterm{Intralesional} Delivered directly into a lesion, such as an injection of corticosteroid into a keloid or a medicine into a skin, joint, or tumor lesion. The drug, dose, sterility, and lesion type determine safety.

\medterm{Intralymphatic} Intralymphatic means within a lymphatic vessel or channel, such as a medicine, tumor spread, infection, or imaging contrast.

\medterm{Intramedullary} Intramedullary means within the medulla of an organ or the central canal of a bone, especially an implant placed inside a long bone.

\medterm{Intramedullary Nail} A metal implant inserted into the medullary canal of a long bone to align and stabilize
a fracture. It shares load along the bone axis and is commonly used for femoral and
tibial shaft fractures.

\medterm{Intramural Hematoma} Collection of blood within the wall of the gastrointestinal tract most commonly the
duodenum. Caused by trauma, anticoagulation, or coagulopathy. Presents with abdominal
pain, vomiting, and sometimes obstructive symptoms. Diagnosis is by CT showing
circumferential wall thickening with high-attenuation hematoma.

\medterm{Intramural Plaque} An intramural plaque is an atherosclerotic deposit within an artery wall that can narrow the lumen or destabilize and cause thrombosis.

\medterm{Intramuscular} Within a muscle. An intramuscular injection delivers a medicine into muscle tissue, commonly in the upper arm, thigh, or hip, where it is absorbed into the bloodstream.

\medterm{Intramuscular Injection} Delivery of a medicine into a muscle with a needle. The injection site and needle size depend on the medicine and the person’s age and body size; pain, bleeding, infection, or nerve injury are possible.

\medterm{Intramuscular Myxoma} An intramuscular myxoma is a rare benign tumor of gelatinous myxoid tissue located within skeletal muscle.

\medterm{Intraocular} Within the eyeball. The term may describe a medicine injection, lens, infection, bleeding, foreign body, pressure, or other finding inside the eye. Increased pressure inside the eye can damage the optic nerve and cause permanent loss of vision.

\medterm{Intraocular Infection} An infection inside the eye, such as endophthalmitis, that can rapidly threaten vision and requires urgent ophthalmic treatment.

\medterm{Intraocular Lens} An artificial lens placed inside the eye, usually during cataract surgery, to focus light after the cloudy natural lens is removed. Lens choices can emphasize distance, near vision, or astigmatism correction. Possible complications include glare, clouding behind the lens, infection, or retinal problems.

\medterm{Intraocular Lens Dislocation} Postoperative displacement of an implanted artificial intraocular lens from its intended
position within the capsular bag or ciliary sulcus into the vitreous cavity or anterior
chamber. It is most commonly caused by progressive zonular dehiscence and may require
surgical repositioning or exchange.

\medterm{Intraocular Lymphoma} Intraocular lymphoma is lymphoma involving the eye, often affecting the vitreous or retina and sometimes occurring with lymphoma in the brain or elsewhere.

\medterm{Intraocular Pressure} The pressure generated by aqueous humor within the eye, normally maintained by a balance
between aqueous production and outflow. Abnormal IOP is an important risk factor for
glaucomatous optic neuropathy, but glaucoma can occur at statistically normal pressures.

\medterm{Intraocular Tumor} Any neoplasm arising within the eye, with uveal melanoma being the most common primary
malignancy in adults and retinoblastoma the most common in children. Diagnosis requires
dilated fundus examination, ultrasound, and MRI; treatment depends on tumour type, size,
and location.

\medterm{Intraoperative} Occurring during a surgical operation, including monitoring, findings, complications, or treatment decisions made during the procedure.

\medterm{Intraoperative Blood Transfusion} Intraoperative blood transfusion replaces lost blood volume and maintains
oxygen-carrying capacity during major surgery. Indications include hemoglobin less than
7 g/dL in stable patients, less than 8 g/dL in cardiac disease, or active hemorrhage
with hemodynamic instability.

\medterm{Intraoperative Neuromonitoring} Electrical tests used during selected operations to watch the function of the brain, spinal cord, nerves, or muscles. Signals can warn the surgical team about a possible change in nerve function, allowing them to reassess the procedure and protect nearby structures.

\medterm{Intraoral} Located within the mouth or performed through the mouth. Located inside the mouth or performed through the mouth, such as an intraoral examination, scan, injection, or surgical procedure.
\medterm{Intraosseous Access} Placement of a vascular-access needle into the medullary cavity of a suitable bone to
permit rapid infusion of fluids, blood products, and medications when intravenous access
is delayed or impossible. It is a temporary emergency route and can cause extravasation,
infection, or compartment injury.

\medterm{Intraosseous Infusion} Rapid delivery of fluids, blood, or medicines through a needle placed into the inner space of a bone when a vein cannot be accessed quickly, usually during an emergency. It is temporary and must be monitored for pain, swelling, leakage, infection, or bone injury.

\medterm{Intrapartum} The period during labor and childbirth, from the onset of regular uterine contractions
leading to cervical dilation until delivery of the baby and placenta.

\medterm{Intrapartum Fever} Fever during labor, which may result from infection, epidural analgesia, dehydration, or other causes and requires maternal and fetal assessment.

\medterm{Intrapartum Stillbirth} Intrapartum stillbirth is fetal death occurring during labor after a live fetus was documented before labor, requiring obstetric investigation and support.

\medterm{Intrapericardial} Located inside the pericardial sac, the thin membrane surrounding the heart. An intrapericardial mass or fluid collection may affect how the heart fills or beats.

\medterm{Intraperitoneal} Intraperitoneal means within the peritoneal cavity or administered into it, including selected medicines, dialysis fluid, or procedures.

\medterm{Intraperitonealroute} The intraperitoneal route delivers a medicine or fluid into the peritoneal cavity, commonly for peritoneal dialysis or selected treatments.

\medterm{Intraprostatic} Located within the prostate gland, describing a lesion, injection, infection, biopsy, procedure, or other finding centered inside prostate tissue.

\medterm{Intrapulmonary} Located within lung tissue or occurring inside the lung’s air spaces, blood vessels, or other structures. Intrapulmonary findings may represent normal anatomy, infection, fluid, bleeding, inflammation, or tumors.

\medterm{Intrapulmonary Shunt} Perfusion of alveoli without ventilation (V/Q = 0). Blood passes through the pulmonary
circulation without gas exchange, mixing with oxygenated blood and causing hypoxemia. It
is seen in ARDS, pneumonia, and atelectasis. It is refractory to supplemental oxygen.

\medterm{Intrasynovial} Intrasynovial means within the synovial lining or fluid-filled space of a joint or tendon sheath.
\medterm{Intratendinous} Located within a tendon, describing a tear, calcium deposit, injection, degenerative change, or other abnormality affecting the tendon's substance.

\medterm{Intratesticular} Located within testicular tissue. An intratesticular mass may be benign or malignant and is usually evaluated with physical examination and scrotal ultrasound, interpreted in the clinical context.

\medterm{Intrathecal} Intrathecal means within the fluid-filled space surrounding the brain and spinal cord; medicines may be given there by lumbar puncture.

\medterm{Intrathecal Chemotherapy} Cancer medicine placed directly into the fluid surrounding the brain and spinal cord, usually through a lumbar puncture or implanted reservoir. It is used for selected leukemias, lymphomas, or cancer spread to this fluid; infection, bleeding, headache, and nervous-system toxicity are possible.

\medterm{Intrathoracic} Intrathoracic means within the chest cavity, including the lungs, heart, mediastinum, vessels, and thoracic lymph nodes.

\medterm{Intrathoracic Localized Obstruction} A focal blockage within the thoracic cavity affecting the airways, esophagus, or great
vessels, leading to localized symptoms depending on the structure involved. Causes
include endobronchial tumors, foreign bodies, esophageal strictures, or vascular
anomalies. Diagnosis is by imaging and endoscopy.

\medterm{Intratracheal} Intratracheal means within the trachea, such as a tube, medicine, foreign body, lesion, or route of administration.

\medterm{Intratympanic} Intratympanic means delivered through the eardrum into the middle ear, often for selected inner-ear disorders or local treatment.

\medterm{Intrauterine} Located or occurring within the uterus, including a pregnancy, device, medicine, lesion, or procedure. The term describes location and requires additional information to explain the finding’s purpose, cause, or risk.

\medterm{Intrauterine Contraceptive Device} A small device placed in the uterus to prevent pregnancy; copper IUDs release copper, while hormonal IUDs release a progestin, and both are highly effective reversible contraception.
\medterm{Intrauterine Device} A small device placed inside the uterus to prevent pregnancy for several years while remaining reversible. Copper devices work without hormones; hormonal devices release a progestogen and often reduce menstrual bleeding. Rare risks include expulsion, perforation during insertion, and infection soon after placement.

\medterm{Intrauterine System} A small T-shaped contraceptive device placed inside the uterus that slowly releases a progestogen hormone to prevent pregnancy. It is also called IUS.

\medterm{Intrauterine Devices} Small devices placed in the uterus to prevent pregnancy; copper and hormonal IUDs are highly effective long-acting reversible contraceptives.

\medterm{Intrauterine Fetal Demise / Stillbirth} Death of a fetus before complete delivery. The gestational-age
definition varies by jurisdiction. Causes may include placental disease, umbilical-cord complications, infection,
genetic abnormalities, maternal disease, or unexplained factors.

\medterm{Intrauterine Growth Restriction} IUGR: a fetus with estimated weight below the 10th percentile (or growth velocity
decline), usually with placental insufficiency or maternal/fetal factors.

\medterm{Intrauterine Pressure Catheter} A fluid-filled catheter inserted into the uterine cavity alongside the fetus to measure
the intensity and frequency of contractions in mm Hg (Montevideo units, MVU).

\medterm{Intrauterine Transfusion} Ultrasonically guided intravascular or intraperitoneal transfusion of fetal red blood
cells to treat severe fetal anemia, most commonly from Rh isoimmunization. Performed
through the umbilical vein using a 22-25 gauge needle. Packed red blood cells,
CMV-negative and crossmatched, are transfused slowly.

\medterm{Intravascular} Intravascular means within a blood vessel, including a catheter, clot, injection, or process occurring in the circulation.

\medterm{Intravascular Catheter Infections} Infections caused by microorganisms entering or growing on a catheter placed inside a blood vessel. They may produce fever, chills, low blood pressure, or bloodstream infection. Diagnosis uses examination and blood cultures, and treatment may require antibiotics and catheter removal.

\medterm{Intravascular Filter} A device placed in a large vein to trap venous clots and reduce pulmonary embolism risk in selected patients who cannot receive
or have not responded adequately to anticoagulation.

\medterm{Intravascular Hemolysis} Destruction of red blood cells within the bloodstream, releasing hemoglobin into plasma; causes include immune hemolysis, mechanical trauma, infections, toxins, and transfusion reactions.

\medterm{Intravascular Lymphoma} A rare, fast-growing blood cancer in which abnormal white blood cells collect inside small blood vessels. It may affect the skin, brain, or other organs, causing fever, skin changes, or neurologic problems.

\medterm{Intravascular Ultrasound} Catheter-based imaging using a miniature ultrasound probe inside coronary arteries to
visualize vessel wall layers, plaque composition, and stent deployment. Provides more
detailed information than angiography alone, including vessel remodeling, plaque burden,
and true vessel diameter.

\medterm{Intravenous} Into or within a vein; intravenous medicines, fluids, or blood products are delivered directly into the bloodstream.

\medterm{Intravenous Anesthesia} Anesthetic induction or maintenance using drugs delivered into a vein. It permits rapid
titration but requires attention to dose, context-sensitive recovery, cardiovascular and
respiratory effects, and the availability of a secure airway.

\medterm{Intravenous Catheters} Flexible devices inserted into a vein to deliver fluids, medicines, blood products, nutrition, or contrast, or to obtain blood samples; complications include infection, thrombosis, infiltration, and extravasation.

\medterm{Intravenous Immunoglobulin} A purified preparation of pooled antibodies given through a vein (IVIg) for selected immune deficiencies, autoimmune diseases, inflammatory disorders, and some neurologic conditions.

\medterm{Intravenous Leiomyomatosis} Extension of benign uterine smooth-muscle tumor into venous channels, sometimes reaching the inferior vena cava or right heart. Imaging is important because extensive disease can cause serious cardiovascular complications.

\medterm{Intravenous Lipid Emulsion Therapy} Administration of a lipid emulsion as a rescue treatment for selected severe lipophilic
drug toxicities, especially local-anesthetic systemic toxicity. It may act as a lipid
sink and through metabolic and cardiovascular effects and requires critical-care
monitoring.

\medterm{Intravenous Pyelogram} An x-ray examination of the urinary tract after intravenous iodinated contrast, historically used to visualize the kidneys, ureters, and bladder; CT urography or ultrasound often replaces it.
of iodinated contrast material. The contrast is filtered by the glomeruli and excreted
through the collecting system, allowing sequential imaging of the renal parenchyma,
calyces, renal pelvis, ureters, and bladder.

\medterm{Intravenous Pyelography} An X-ray examination of the urinary tract after intravenous contrast administration; it is now often replaced by CT urography or ultrasound.

\medterm{Intravenous Sedation} Administration of sedative medicines through a vein to reduce anxiety, discomfort, or awareness during a procedure while maintaining appropriate monitoring and readiness to support breathing.

\medterm{Intraventricular Hemorrhage} Bleeding into the germinal matrix and ventricular system, primarily in premature infants
less than 32 weeks gestation. Graded I-IV by severity on cranial ultrasound. Grade I
(germinal matrix only) and Grade II (intraventricular without dilation) have good
prognosis.

\medterm{Intraventricular Hemorrhage Of The Newborn} Intraventricular hemorrhage in a newborn is bleeding into the brain’s ventricular system, most often in premature infants, and can cause hydrocephalus or neurologic injury.

\medterm{Intravital Microscopy} Intravital microscopy uses a microscope to observe cells or tissues in a living organism, usually in research.

\medterm{Intravitreal Injection} An ophthalmological procedure in which medication is injected directly into the vitreous
cavity of the eye through the pars plana, a safe entry zone located approximately 3.5 to
4 millimeters posterior to the limbus, to deliver therapeutic agents directly to the
retina and vitreous for the treatment of various retinal diseases.

\medterm{Intrinsic} Belonging naturally to a structure or process and arising from within it, rather than being externally imposed.

\medterm{Intrinsic Factor} A protein secreted by stomach parietal cells that is required for vitamin B12 absorption in the terminal ileum.

\medterm{Intrinsic Renal AKI} Acute kidney injury caused by damage to renal tissue. Causes include acute tubular injury,
acute interstitial nephritis, glomerulonephritis, and vascular disorders. Laboratory findings vary with the cause and
should not be interpreted in isolation.

\medterm{Intrinsic Sphincter Deficiency} Weakness of the urinary sphincter that causes stress urinary incontinence.

\synonyms
ISD

\medterm{Introitus} An entrance or opening, especially the entrance to the vagina.

\medterm{Intron} A non-coding segment of a gene that is transcribed into RNA but spliced out before the
messenger RNA is translated into a protein.

\medterm{Intrusive Thought} An unwanted, distressing thought, image, or impulse that enters awareness repeatedly; it does not necessarily reflect intent or desire.

\medterm{Intubate} To place a tube into a body passage, especially an endotracheal tube through the mouth or nose into the trachea to maintain an airway or support ventilation.

\medterm{Intubation} Placement of a tube into the airway or another body passage to maintain it or deliver treatment.

\medterm{Intussusception} Condition in which a segment of intestine telescopes into the adjacent distal segment
causing bowel obstruction and compromise of the mesenteric blood supply. It is the most
common cause of intestinal obstruction in infants and children between six months and
three years of age.

\medterm{Intussusception – Children} Intussusception occurs when one segment of intestine telescopes into another, causing intermittent pain, vomiting, obstruction, and possible bowel ischemia in children.

\medterm{Inulin} Inulin is a plant fructan used as a dietary fiber and in tests of glomerular filtration because it is freely filtered and minimally handled by the kidney.

\medterm{Inuresis} A rare or obsolete spelling related to enuresis, meaning involuntary urination, especially during sleep.

\medterm{Inutero} A variant or misspelling of in utero, meaning within the uterus before birth.

\medterm{Invasive} Able to enter and spread through tissues; the term may describe an infection, tumor, procedure, or disease that penetrates nearby structures.

\medterm{Invasive Candidiasis} A serious Candida infection that enters the bloodstream or deep organs, most often in people who
are critically ill or immunocompromised. It may cause persistent fever or organ dysfunction and
requires prompt medical diagnosis and treatment.

\medterm{Invasive Ductal Carcinoma} The most common type of breast cancer, characterized by malignant cells that have broken
through the duct wall and invaded surrounding breast tissue. It can metastasize via
lymphatic or hematogenous spread.

\medterm{Invasive Lobular Carcinoma} A breast cancer characterized by discobesive cells infiltrating in single-file lines.
More frequently bilateral and multifocal than ductal carcinoma and harder to detect on
imaging; usually hormone receptor positive.

\medterm{Invasive Mole} A hydatidiform mole that invades the uterine myometrium or vascular spaces; it is a form of gestational trophoblastic neoplasia and may persist or metastasize.
\medterm{Invasive Procedure} A medical procedure that enters the body through a needle, tube, incision, or other instrument. It may diagnose or treat disease and can cause pain, bleeding, infection, injury, or complications related to anesthesia.

\medterm{Inverse Agonist} A type of drug that binds to a cell receptor and actively produces the opposite
biological effect of an agonist, effectively turning down the receptor's baseline
activity.

\medterm{Inversion} Chromosomal segment reversed end-to-end. Paracentric: does not include centromere.
Pericentric: includes centromere. Usually benign unless breakpoint disrupts gene.

\medterm{Invert} To turn upside down, reverse, or turn inward; in anatomy, inversion can describe movement of the sole of the foot toward the midline.

\medterm{Invert Sugar} A mixture of glucose and fructose produced by splitting sucrose, used as a sweetener in foods and medicines. It has nutritional effects similar to other added sugars and contributes to total carbohydrate and energy intake.

\medterm{Invertase} An enzyme that splits sucrose, table sugar, into glucose and fructose. It is produced by some plants, fungi, and microorganisms and is also used in food processing.

\medterm{Invertebrata} Invertebrata is a broad nonformal grouping of animals without a vertebral column, including worms, mollusks, arthropods, and many other groups.

\medterm{Invertebrate} An animal without a backbone, including insects, worms, mollusks, spiders, and many other groups with diverse body structures and life cycles.

\medterm{Invertebrate Hormones} Chemical signals regulating growth, development, reproduction, molting, behavior, or metabolism in animals without backbones. Studying them helps explain animal biology and can support development of selective pest-control methods.

\medterm{Inverted Papilloma} A usually benign but locally destructive tumor arising most often in the nasal cavity or sinuses; it can recur and rarely become cancerous.

\medterm{Investigational} Being studied in clinical research and not yet established as safe and effective for a particular use. Investigational medicines and devices require appropriate oversight, informed consent, and clear communication of uncertainty.

\medterm{Inveterate} Long-standing, deeply established, or difficult to change, often describing a habit, condition, or behavior. The word indicates persistence but does not identify the cause, severity, treatment, or expected outcome.

\medterm{Involucrum} An involucrum is a sheath of new bone that forms around dead bone, classically in chronic osteomyelitis.

\medterm{Involuntary} Occurring without conscious control or deliberate choice, such as a reflex, autonomic function, or involuntary muscle movement.

\medterm{Involuntary Movement} A movement that occurs without conscious control, such as a tremor, tic, chorea, dystonia, or myoclonus.

\medterm{Involuntary Muscle} Muscle that contracts without conscious control, including smooth muscle in internal organs and cardiac muscle in the heart.

\medterm{Involuntary Psychiatric Treatment} Assessment, hospitalization, or treatment without the patient’s consent under applicable
law when serious mental disorder creates a substantial risk of harm or severe inability
to meet basic needs. The least restrictive alternative, due process, and ongoing
reassessment are required.

\medterm{Involute} To shrink, return toward a former size, or undergo involution; in anatomy, an involute structure may be tightly coiled or rolled inward.

\medterm{Involution} The process by which the uterus returns to its prepregnancy size after delivery. The
uterus weighs approximately 1000 grams immediately postpartum and returns to
approximately 60 grams by 6 weeks. Fundal height decreases approximately 1 cm per day.

\medterm{Iodamoeba Butschlii} A usually non-disease-causing intestinal amoeba spread by swallowing cysts in contaminated food or water. Finding it in stool suggests fecal exposure but usually does not require treatment unless another infection is present.

\medterm{Iodide} The negatively charged form of iodine, used by the thyroid to make hormones and present in medicines, supplements, and contrast compounds.

\medterm{Iodide-Associated Sialadenitis} Inflammation and swelling of the salivary glands caused by excessive iodide intake,
typically from medications such as amiodarone or iodine-containing contrast, resolving
with dose reduction.

\medterm{Iodides} Salts containing iodine, sometimes used in carefully selected thyroid disorders.

\medterm{Iodinated Contrast Media} Contrast agents containing iodine used primarily for radiography, CT, angiography, and
some fluoroscopic procedures. Important risks include immediate hypersensitivity-like
reactions, extravasation, thyroid effects, and kidney injury in susceptible patients.

\medterm{Iodine} An essential mineral used to make thyroid hormones, which regulate growth, brain development, and metabolism. Too little iodine can cause thyroid enlargement and hormone deficiency, especially harmful during pregnancy and childhood; too much can also disturb thyroid function. Main sources include iodized salt, seafood, dairy, and some supplements.

\medterm{Iodine Deficiency} Inadequate iodine intake or availability, which can impair thyroid-hormone production and cause goiter, hypothyroidism, or developmental problems in a fetus or child.


\medterm{Iodine Poisoning} Iodine poisoning follows excessive iodine exposure and can irritate the gastrointestinal tract, disrupt thyroid function, and cause systemic toxicity depending on dose and formulation.

\medterm{Iodine-123} Pure gamma emitter, half-life 13 hours. Used for thyroid imaging. No beta emission,
lower radiation dose than I-131.

\medterm{Iodine-131} A radioactive iodine isotope used in selected diagnostic and therapeutic procedures, especially for thyroid disease.
\medterm{Iodipamide} An iodinated contrast medium formerly used for imaging the biliary system and urinary tract; it can cause allergic or kidney-related reactions.

\medterm{Iodoacetamide} A chemical reagent that irreversibly modifies sulfhydryl groups, especially cysteine residues, and is used in biochemical research.

\medterm{Iododerma} A papulopustular or acneiform skin eruption caused by iodine excess, presenting on the
face and extensor surfaces, resolving with iodine restriction and symptomatic treatment.

\medterm{Iodoform} A yellow iodine-containing compound formerly used as an antiseptic and in some wound or dental preparations. It can irritate tissue and cause toxicity if absorbed or swallowed and is now used only in selected contexts.

\medterm{Iohexol} A water-soluble iodine-containing contrast medicine used during CT scans, angiography, and other imaging studies to make vessels and organs easier to see.

\medterm{Ion} An atom or molecule carrying a positive or negative electric charge. Physiologically important ions include sodium, potassium, calcium, chloride, and bicarbonate, whose concentrations influence nerves, muscles, fluid balance, and acid-base status.

\medterm{Ion Transport} Movement of charged particles across a cell membrane through channels, carriers, or pumps. It maintains membrane potentials, nerve signaling, muscle contraction, fluid balance, and epithelial transport.

\medterm{Ionic Strength} A measure of the total concentration and charge of ions in a solution. It affects protein interactions, membrane behavior, chemical equilibria, and the performance of some laboratory assays and medicines.

\medterm{Ionization} The formation of charged particles by gaining or losing electrons or protons. A medicine's degree of ionization affects its solubility, membrane passage, absorption, distribution, and urinary excretion.

\medterm{Ionizing Radiation} Radiation with enough energy to remove electrons from atoms and create ions, including
X-rays, gamma rays, and certain particulate emissions. It can cause deterministic tissue
effects at sufficient doses and stochastic cancer risk at lower doses.

\medterm{Ionomycin} A laboratory compound that carries calcium across cell membranes and raises calcium inside cells. It is used mainly in research to study cell signaling, secretion, and immune-cell activation, not as a routine medicine.

\medterm{Ionophore} A molecule that transports ions across a lipid membrane. Ionophores have research, antimicrobial, and veterinary uses; some can disrupt human cellular ion balance and are toxic.

\medterm{Ions} Electrically charged atoms or molecules, including the electrolytes that regulate nerve conduction, muscle function, hydration, and acid-base balance. Abnormal ion concentrations can cause weakness, seizures, arrhythmias, or altered consciousness.

\medterm{Iontophoresis} Use of a mild electrical current to drive an ionized medicine through the skin or alter sweat-gland activity. It is used in selected dermatologic and rehabilitation settings; skin irritation and burns can occur.

\medterm{Iopamidol} A nonionic iodine-containing contrast medicine used to improve CT, angiography, and other medical images. It can cause warmth, nausea, allergy, kidney problems, or thyroid effects, so risks require assessment beforehand.

\medterm{Iothalamate Clearance} Iothalamate clearance estimates glomerular filtration by measuring how quickly injected iothalamate is cleared from blood or urine.

\medterm{Iothalamate Meglumine} Iothalamate meglumine is an iodinated contrast agent formerly used in radiographic imaging and associated with dose- and kidney-related risks.

\medterm{Iothalamic Acid} An older water-soluble, iodinated contrast agent formerly used for radiographic imaging. It has largely been replaced by newer agents; prior contrast reactions, kidney function, and thyroid disease remain clinically relevant when reviewing exposure history.

\medterm{Ioxaglic Acid} An older iodinated contrast medium used for radiographic imaging, including selected angiographic and urinary-tract studies. It has largely been superseded by newer agents, but its history is relevant when assessing prior contrast reactions.

\medterm{IPC} Alternate name; see \textit{Indwelling Pleural Catheter}.

\medterm{Ipecac} A vomiting medicine formerly used after poisoning, now not routinely recommended because vomiting can cause aspiration and delay more appropriate emergency treatment.

\medterm{Ipilimumab} An immunotherapy antibody that blocks CTLA-4, removing a brake on T cells so they can attack cancer. It can cause immune inflammation in the bowel, skin, liver, lungs, hormone glands, or other organs.

\medterm{IPMN Pancreas (Pancreatic Cysts)} Alternate index label for Pancreatic Cysts.

\medterm{Ipratropium} Ipratropium is an inhaled short-acting muscarinic antagonist that reduces vagally mediated bronchoconstriction and nasal secretion; it is used in COPD and acute asthma with bronchodilators.

\medterm{Ipratropium Bromide} An inhaled anticholinergic (antimuscarinic) bronchodilator that blocks acetylcholine at
airway smooth muscle muscarinic receptors. It reduces bronchospasm and mucus secretion.
It is used in COPD and acute asthma exacerbations. It is less effective than beta-2
agonists in asthma but more effective in COPD.

\medterm{Iproniazid} Iproniazid is an older monoamine oxidase inhibitor antidepressant that is no longer commonly used because of serious liver toxicity and drug interactions.

\medterm{Ipsilateral} On or affecting the same side of the body as a stated structure or injury. For example, weakness in the right arm and right leg is ipsilateral to a finding on the right side of the body; the opposite term is contralateral.

\medterm{IQ Testing} IQ testing uses standardized tasks to estimate intellectual functioning; results should be interpreted with education, language, development, and clinical context.

\medterm{Iratumumab} An investigational antibody studied for cancers involving CD30-positive cells. It was designed to help the immune system recognize these cells, but it is not an established routine treatment.

\medterm{Irbesartan} Irbesartan is an angiotensin-II receptor blocker used for hypertension and diabetic kidney disease; it can cause hyperkalaemia, hypotension, and renal-function changes and is contraindicated in pregnancy.

\medterm{Iridium} Iridium is a dense metal whose radioactive isotopes, especially iridium-192, are used in industrial testing and some medical brachytherapy.

\medterm{Iridocyclitis} Inflammation involving the iris and ciliary body, a form of anterior uveitis. It may
cause eye pain, photophobia, redness, blurred vision, and cells or flare in the anterior
chamber and may be associated with autoimmune disease, infection, trauma, or other
causes.

\medterm{Iridotomy} A procedure that creates a small opening in the iris to allow aqueous humor to pass between the posterior and anterior
chambers. Laser peripheral iridotomy is used to treat or prevent selected forms of
angle-closure glaucoma.

\medterm{Irinotecan} A chemotherapy medicine that inhibits topoisomerase I and is used for colorectal and several other cancers.

\medterm{Irinotecan Hydrochloride} A chemotherapy medicine that blocks topoisomerase I and damages cancer-cell DNA, used for selected cancers with diarrhea and low blood counts possible.

\medterm{Iris} The colored, muscular ring between the cornea and lens that controls the size of the pupil and therefore the amount of light entering the eye. Its muscles constrict the pupil in bright light and enlarge it in dim light.

\medterm{Iris Atrophy} Thinning or loss of iris tissue resulting from trauma, chronic inflammation, ischaemia,
or neurodegenerative disease, leading to a translucent or patched iris appearance. It
may be sectoral or diffuse and can be associated with decreased pupillary response,
photophobia, or secondary glaucoma.

\medterm{Iris Cyst} An iris cyst is a fluid-filled lesion on or within the iris, usually benign but sometimes affecting vision, pressure, or the appearance of the eye.

\medterm{Iris Disease} A disorder of the colored part of the eye, including inflammation, injury, infection, developmental abnormalities, or tumors.

\medterm{Iris Melanoma} Iris melanoma is a malignant tumor of pigment cells in the iris and may appear as a growing pigmented lesion, pupil distortion, or glaucoma.

\medterm{Iris Neoplasm} An iris neoplasm is an abnormal growth arising in the colored part of the eye; it may be benign or malignant and can affect vision or eye pressure.

\medterm{Iritis} Inflammation of the iris, the most common form of anterior uveitis. Presents with acute
onset eye pain, photophobia, blurred vision, and redness (ciliary flush). Signs: cells
and flare in the anterior chamber, keratic precipitates (KP), synechiae (adhesions
between iris and lens).

\medterm{Irofulven} Irofulven is an investigational anticancer compound derived from illudin that damages DNA and has been studied in solid tumors.

\medterm{Iron} Essential for hemoglobin, myoglobin, and electron transport. Heme iron from animal foods
is generally more bioavailable than non-heme iron from plant foods. Deficiency is the
most common nutritional deficiency worldwide and may result from blood loss, increased
requirements, poor intake, or malabsorption.

\medterm{Iron Deficiency} The most common nutritional deficiency worldwide, resulting from inadequate iron intake,
impaired absorption, or chronic blood loss, leading to reduced hemoglobin synthesis and
microcytic hypochromic anemia. Clinical features include fatigue, pallor, koilonychia,
pica, and angular cheilitis.

\medterm{Iron Deficiency Anemia} Anemia caused by insufficient iron for hemoglobin production. It is often microcytic and hypochromic; treatment requires identifying the cause and replacing iron when appropriate.


\medterm{Iron Overdose} Iron overdose can cause gastrointestinal injury, metabolic acidosis, shock, liver failure, and multiorgan toxicity, particularly after a large ingestion by a child.

\medterm{Iron Overload} Excess iron stored in the body, which can damage the liver, heart, pancreas, joints, or endocrine organs.

\synonyms
Overload Iron (Iron Overload)

\medterm{Iron-Deficiency Anaemia} Anaemia caused by insufficient iron for haemoglobin production, commonly due to blood loss, inadequate intake, or impaired absorption.

\medterm{Iron-Deficiency Anaemia In Adolescents} Iron-deficiency anaemia during adolescence, commonly related to rapid growth, inadequate dietary iron, or menstrual blood loss. Assessment should identify the cause and guide dietary advice and clinical treatment.

\medterm{Iron-Deficiency Anemia} Anemia caused by insufficient iron for normal hemoglobin production. It may cause fatigue, weakness,
shortness of breath, paleness, headache, or unusual cravings such as ice; the cause of iron
loss or inadequate intake should be identified.

\medterm{Irradiated Foods} Irradiated foods are exposed to controlled ionizing radiation to reduce pathogens, insects, or spoilage; the process does not make food radioactive.

\medterm{Irradiation} Exposure to or treatment with radiation; irradiation may be used diagnostically or therapeutically, or may cause tissue injury depending on dose and type.

\medterm{Irregular} Not following the usual pattern, rhythm, shape, timing, or order. In medicine, it may describe an uneven heartbeat, an unpredictable menstrual cycle, an unusual skin edge, or a result that does not match the expected pattern.

\medterm{Irregular Menstruation} Menstrual bleeding that is unusually early, late, frequent, infrequent, prolonged, heavy, or unpredictable for the person’s age and reproductive context. Causes include pregnancy, ovulatory disorders, endocrine disease, medicines, structural lesions, and perimenopause.

\medterm{Irregular Respiration} An abnormal breathing pattern in rate, depth, rhythm, or pauses. It may occur with sleep disorders, metabolic disturbance, lung disease, drug effects, or neurologic injury and may require urgent assessment if accompanied by cyanosis or reduced consciousness.

\medterm{Irregular Sleep-Wake Syndrome} A circadian rhythm sleep-wake disorder with multiple irregular periods of sleep and
wakefulness across the 24-hour day rather than one consolidated nighttime sleep period.
It is associated with neurological disease, dementia, aging, and reduced exposure to
environmental time cues.

\medterm{Irreversible Cell Injury} Cell damage beyond the point of recovery, resulting in death by necrosis, apoptosis, or
another regulated cell-death pathway. Key features include inability to reverse
mitochondrial dysfunction and severe membrane or nuclear damage.

\medterm{Irrigate} To wash or flush a body part, wound, cavity, or medical device with a fluid to remove debris, contaminants, or secretions.

\medterm{Irrigation} Flushing a wound, body cavity, or organ with fluid to remove blood, debris, microorganisms, or unwanted material.

\medterm{Irrigation Solution} A sterile or appropriately prepared liquid used to wash a wound, body cavity, organ, or surgical field. The solution must be suitable for the tissue because contamination, pressure injury, electrolyte shifts, or chemical irritation can cause harm.

\medterm{Irritable Bowel Syndrome} A disorder of gut--brain interaction that causes repeated abdominal pain and changes in bowel habits, such as diarrhea, constipation, or both. No visible structural damage is usually found in the digestive tract.

\synonyms
IBS

\medterm{Irritable Bowel Syndrome Diet} An individualized eating plan used to reduce irritable bowel syndrome symptoms. A limited
trial of a low-FODMAP diet may help some people and should be followed by structured
food reintroduction.

\medterm{Irritable Mood} A mood marked by increased sensitivity, impatience, anger, or a tendency to become upset easily.
\medterm{Irritant Contact Dermatitis} Inflammation of the skin caused by direct chemical or physical injury rather than an immune-mediated allergy. It commonly produces redness, burning, scaling, fissures, or pain at the area of exposure.

\medterm{Irritants} Substances that inflame or damage tissues and cause symptoms without necessarily triggering a true immune allergy. Examples include smoke, chemicals, soaps, fumes, and pollutants affecting skin, eyes, airways, or digestion.

\medterm{Irritation} A local unpleasant or inflammatory response to mechanical, chemical, thermal, infectious, or allergic stimuli. It may cause redness, itching, burning, pain, cough, or tearing depending on the tissue involved.

\medterm{Isaacs Syndrome} A rare peripheral nerve hyperexcitability disorder causing continuous muscle-fibre
activity, stiffness, cramps, fasciculations, myokymia, and delayed muscle relaxation. It
may be autoimmune, paraneoplastic, or associated with inherited or acquired
peripheral-nerve disorders.

\medterm{Isavuconazole} A triazole antifungal agent used in the treatment of invasive aspergillosis and
mucormycosis. Isavuconazole inhibits fungal lanosterol 14-alpha demethylase, disrupting
ergosterol synthesis and fungal cell membrane integrity.

\medterm{Isazophos} Isazophos is an organophosphate pesticide that can inhibit acetylcholinesterase and cause cholinergic toxicity after exposure.
\medterm{Ischaemia} Inadequate blood flow and oxygen delivery to a tissue, usually because of narrowed or blocked blood vessels, causing dysfunction or injury.

\medterm{Ischaemic Heart Disease} Heart disease caused by reduced blood supply to the heart muscle, usually as a result of coronary artery disease; it may cause angina or myocardial infarction.

\medterm{Ischemia} Reduced blood flow causing too little oxygen for a tissue’s needs. It can cause pain, impaired function, permanent injury, or tissue death and may affect the heart, brain, limbs, or bowel.

\medterm{Ischemia, Intestinal} Alternate index label; see \textit{Intestinal ischemia}.

\medterm{Ischemia-Reperfusion Injury} The paradoxical tissue damage that occurs when blood flow is restored to previously
ischemic tissue, a phenomenon with major clinical significance in myocardial infarction
treatment, stroke management, organ transplantation, and shock resuscitation.

\medterm{Ischemic Bone Necrosis} Alternate index label; see \textit{Avascular necrosis (osteonecrosis)}.

\medterm{Ischemic Cardiomyopathy} Weakening of the heart muscle caused by coronary artery disease and previous or ongoing
myocardial ischemia. Scarred or poorly perfused myocardium can reduce contractile function and increase the risk of heart
failure and ventricular arrhythmias.

\medterm{Ischemic Chest Pain (Angina Symptoms)} Alternate index label for Angina Symptoms.

\medterm{Ischemic Colitis} Inflammation and injury of the colon caused by reduced blood flow. It often causes sudden abdominal pain, bloody diarrhea, or rectal bleeding and commonly affects areas such as the splenic flexure.

\synonyms
Colitis, Ischemic

\medterm{Ischemic Fasciitis} A benign fibroblastic and myofibroblastic proliferation associated with pressure-related ischemic injury, usually arising in deep soft tissue over bony prominences in debilitated or immobile patients.

\medterm{Ischemic Gangrene} Death of tissue due to prolonged severe arterial insufficiency, most commonly in the
extremities. Presents as dry, black, mummified tissue with a clear line of demarcation
from viable tissue. Requires revascularization when possible and surgical debridement or
amputation of necrotic tissue.

\medterm{Ischemic Heart Disease} Heart disease caused by inadequate blood flow through narrowed or blocked coronary arteries. It may cause chest pressure, breathlessness, heart attack, rhythm problems, heart failure, or sudden death.

\synonyms
Coronary artery disease; coronary heart disease

\medterm{Ischemic Preconditioning} A phenomenon in which brief, nonlethal episodes of ischemia make tissue more resistant to a later, prolonged ischemic insult. It is important in cardiovascular and research physiology, but clinical applications remain context-dependent.

\medterm{Ischemic Priapism} Alternate index label; see \textit{Priapism}.

\medterm{Ischemic Stroke} A neurologic deficit caused by inadequate blood flow to part of the brain, usually from
thrombotic or embolic arterial occlusion. Causes include atherosclerotic disease, cardiac
embolism, and small-vessel disease.


\medterm{Ischium} Inferior-posterior hip bone. Ischial tuberosity bears seated weight and serves as
hamstring attachment. Ischial spine is a pudendal nerve block landmark and determines
pelvic outlet level. Lesser sciatic notch lies between spine and tuberosity.

\medterm{Ishak Score} A pathology score describing the amount of liver inflammation and scarring seen in a biopsy, especially in chronic hepatitis. Higher scores indicate more advanced damage, but interpretation requires the full clinical picture.

\medterm{Islet Cell} A cell in a pancreatic islet that helps control blood glucose; beta cells make insulin and alpha cells make glucagon.

\medterm{Islet Cell Cancer} Alternate index label; see \textit{Pancreatic neuroendocrine tumors}.

\medterm{Islets Of Langerhans} Small clusters of hormone-producing cells in the pancreas that release insulin, glucagon, and other hormones regulating blood sugar and metabolism.

\medterm{Iso-Sulfan Blue} A blue dye used to map lymphatic drainage or identify sentinel lymph nodes during selected surgical procedures.

\medterm{Isoaminile} Isoaminile is an older sympathomimetic appetite suppressant that is no longer widely used because of cardiovascular and other adverse effects.

\medterm{Isoantibodies} Antibodies made against antigens from another individual of the same species, such as antibodies involved in blood-group incompatibility.

\medterm{Isoantigen} An antigen that differs between individuals of the same species and can trigger an immune response after exposure.

\medterm{Isocarboxazid} Isocarboxazid is a monoamine oxidase inhibitor used for depression; it has important interactions with other medicines and tyramine-rich foods.

\medterm{Isochromosomes} Abnormal chromosomes with two identical arms caused by division across the wrong axis of the centromere.

\medterm{Isocitrate Dehydrogenase} A regulated enzyme of the citric acid cycle that converts isocitrate to
alpha-ketoglutarate while producing NADH and releasing carbon dioxide. Mitochondrial
isoforms are stimulated by ADP and calcium and inhibited by ATP and NADH. Mutations in
IDH1 or IDH2 can produce the oncometabolite 2-hydroxyglutarate.

\medterm{Isocyanate} A reactive chemical group found in some industrial products; exposure can irritate the airways and cause occupational asthma or skin reactions.

\medterm{Isoelectric Point} The acidity level at which a molecule, especially a protein, carries no overall electrical charge and may behave differently in solution.

\medterm{Isoenzyme} One of several forms of an enzyme that catalyze the same chemical reaction but differ in structure, location, or regulation. Measuring selected isoenzymes can help identify the tissue affected by injury.

\medterm{Isoetharine} An older airway-opening medicine related to beta agonists, formerly used to relieve narrowed airways but largely replaced by newer treatments.

\medterm{Isofenphos} A poisonous organophosphate insecticide. Exposure can cause sweating, excess saliva, vomiting, diarrhea, small pupils, muscle twitching, breathing difficulty, seizures, or weakness because it disrupts nerve signals; suspected poisoning needs urgent medical help.

\medterm{Isoflavone} A plant compound, especially from soy, with weak estrogen-like activity that is studied in relation to bone, menopause, and cancer risk.

\medterm{Isoflurane} Isoflurane is an inhaled anesthetic used to maintain general anesthesia; it can lower blood pressure and depress breathing.

\medterm{Isoflurophate} Isoflurophate is a long-acting organophosphate cholinesterase inhibitor formerly used in eye treatments; it can cause prolonged cholinergic effects.

\medterm{Isogeneic Graft} A transplant between genetically identical individuals, such as identical twins, so immune rejection risk is relatively low.

\medterm{Isokinetic Exercise} Muscle contraction at constant speed throughout range of motion. Uses specialized
equipment. Assesses and improves strength.

\medterm{Isolate} To separate a person, organism, tissue, or substance from others; infection-control isolation separates people who may transmit a pathogen.

\medterm{Isolate Compound} To separate a single chemical compound from a mixture for identification, testing, or research. In clinical laboratories, isolation and characterization may support toxicology, pharmacology, or natural-product analysis.

\medterm{Isolated Systolic Hypertension} Elevated systolic blood pressure with normal diastolic pressure, commonly related to arterial stiffness in older adults and associated with increased cardiovascular risk.

\medterm{Isolation} Separating a person with a contagious infection from others to limit spread, using appropriate rooms, protective equipment, and infection-control procedures.

\medterm{Isolation (Social)} Lack of social contacts and interactions. Common in elderly: bereavement, mobility
limitation, sensory impairment, retirement. Consequences: depression, cognitive decline,
functional impairment, mortality. Interventions: social programs, transportation,
technology, volunteer visits.

\synonyms
Social

\medterm{Isolation Precautions} Measures used with standard precautions to reduce transmission of infection. Additional precautions may be contact, droplet, or airborne, depending on the mode of spread.

\medterm{Isoleucine} An essential amino acid obtained from food that helps build proteins, support energy production, maintain muscle, and repair body tissues.

\medterm{Isomerase} An enzyme that rearranges atoms within a molecule to form an isomer. Isomerases participate in metabolism, and inherited defects in related pathways can cause metabolic disease.

\medterm{Isometheptene} Isometheptene is a vasoconstrictor formerly used in some headache medicines; it may raise blood pressure and is unsuitable for some cardiovascular conditions.

\medterm{Isometric} Describing muscle contraction that produces force without a meaningful change in muscle length or joint movement.

\medterm{Isometropia} The condition in which both eyes have approximately equal refractive power, so neither eye has substantially more near-sightedness, far-sightedness, or astigmatism than the other.

\medterm{Isoniazid} Isoniazid is a first-line antituberculous drug that inhibits mycolic-acid synthesis in Mycobacterium tuberculosis; hepatotoxicity and peripheral neuropathy are important risks, and pyridoxine is often given.

\medterm{Isoniazid Toxicity} Harm from excessive or unusual sensitivity to isoniazid, especially liver injury, nerve damage, or seizures. Treatment may include vitamin B6, supportive care, activated charcoal, seizure control, and urgent monitoring.

\medterm{Isoprenaline} Isoprenaline, also called isoproterenol, is a beta-adrenergic medicine that can increase heart rate and is used in selected emergency cardiac situations.

\medterm{Isoprenylation} Addition of a lipid-derived isoprenyl group to a protein, often helping determine its membrane location and signaling activity. Abnormal isoprenylation contributes to some cancers and is a target of investigational therapies.

\medterm{Isoprocarb} Isoprocarb is a carbamate insecticide that can cause cholinergic poisoning if swallowed or absorbed in a significant amount.

\medterm{Isopropamide} Isopropamide is an anticholinergic medicine formerly used for gastrointestinal spasms; it may cause dry mouth, blurred vision, constipation, or urinary retention.

\medterm{Isopropanol Alcohol Poisoning} Poisoning from isopropanol, often found in rubbing alcohol, causing intoxication, vomiting, abdominal pain, low blood pressure, and sometimes coma.

\medterm{Isopropyl Alcohol} A volatile alcohol used as a disinfectant and solvent. Ingestion can cause rapid intoxication, vomiting, abdominal pain, hypotension, and coma; management is supportive and differs from treatment of methanol or ethylene-glycol poisoning.

\medterm{Isoprostane} A compound formed when free radicals oxidize membrane lipids; it is used as a marker of oxidative stress.

\medterm{Isoproterenol} Isoproterenol, also called isoprenaline, is a beta-adrenergic medicine that increases heart rate and contractility and is used in selected emergency settings.

\medterm{Isoquinolines} A group of nitrogen-containing ring compounds that includes chemicals studied in pharmacology, natural products, and organic chemistry.

\medterm{Isosorbide} Isosorbide is a nitrate medicine that relaxes blood vessels and is used mainly to prevent or relieve angina; headache and low blood pressure are common adverse effects.

\medterm{Isosorbide Dinitrate} An organic nitrate (and prodrug) that releases nitric oxide, relaxing vascular smooth
muscle—predominantly venous (preload reduction) with coronary dilation.

\medterm{Isospora} An older name for a group of intestinal coccidian parasites, now generally classified as Cystoisospora. Infection can cause watery diarrhea, especially in people with weakened immunity.

\medterm{Isospora Belli} The former name for Cystoisospora belli, a parasite spread by swallowing oocysts in contaminated food or water. It causes watery diarrhea, cramps, nausea, and weight loss and can be severe or prolonged in immunocompromised people.

\medterm{Isosporiasis} Intestinal infection caused by the protozoan Cystoisospora belli. It typically causes
watery diarrhea, abdominal cramps, nausea, fever, and weight loss and may be prolonged
or severe in people with impaired cellular immunity.

\medterm{Isothiocyanate} A reactive compound formed from glucosinolates in some plants, especially cruciferous vegetables, and studied for antimicrobial and cancer-related effects.

\medterm{Isothiocyanates} Reactive plant chemicals formed from glucosinolates, especially in cruciferous vegetables. They are studied for possible protective effects against cancer, but concentrated supplements may have different effects from normal food amounts.

\medterm{Isotonic} Describing muscle contraction in which muscle length changes to move a joint under a relatively constant external load.

\medterm{Isotonic Exercise} Exercise in which muscles change length while producing force against a relatively constant load, as in lifting or lowering a weight.

\medterm{Isotonic Solution} A solution with an effective osmotic concentration close to that of plasma, such as normal saline or balanced crystalloid. Isotonic fluids may expand extracellular volume, but choice and amount must be tailored because excess can cause fluid overload or electrolyte disturbance.

\medterm{Isotope} One of two or more atoms of the same element with the same number of protons but different numbers of neutrons.

\medterm{Isotretinoin} A systemic retinoid used for severe, recalcitrant nodular acne. It reduces sebum
production, normalizes follicular keratinization, decreases Propionibacterium acnes, and
has anti-inflammatory effects. Dose is 0.5 to 1 mg/kg/day for 15 to 20 weeks, targeting
cumulative dose of 120 to 150 mg/kg.

\medterm{Isovaleric Acidemia} A rare inherited disorder in which the body cannot properly break down the amino acid leucine. Harmful substances build up and may cause poor feeding, vomiting, unusual sleepiness, low muscle tone, or a sweaty-feet odor, especially during illness or fasting. Newborn screening can detect it; treatment uses a controlled diet and prevents prolonged fasting.

\synonyms
IVA; Isovaleryl-CoA Dehydrogenase Deficiency

\medterm{Isoxsuprine} Isoxsuprine is a vasodilating medicine formerly used for selected circulation disorders; it may cause dizziness, fast heartbeat, or low blood pressure.

\medterm{Isozyme} One of several molecular forms of an enzyme that catalyze the same reaction but differ in structure, tissue distribution, or regulation.

\medterm{Ispaghula Husk} Ispaghula husk, or psyllium, is a soluble fiber that absorbs water and helps regulate stool consistency; adequate fluid is important when using it.

\medterm{Ispinesib} Ispinesib is an investigational kinesin-spindle protein inhibitor studied for anticancer effects by disrupting chromosome separation during cell division.

\medterm{Isradipine} Isradipine is a calcium-channel blocker used to lower blood pressure; it may cause flushing, headache, ankle swelling, or dizziness.

\medterm{ISUP Grade Group} A five-level system that grades prostate cancer by how abnormal the cancer looks under a microscope. Group 1 is the lowest risk, usually with a Gleason score of 6 or less. Groups 2 and 3 represent intermediate risk, while groups 4 and 5 represent high or highest risk. The result helps guide monitoring, surgery, radiation, or other treatment.

\synonyms
Grade Group; ISUP Prostate Grading System; WHO/ISUP Grade Group

\medterm{Itch} Itch, or pruritus, is an unpleasant skin sensation that provokes scratching and may arise from dermatitis, allergy, infection, medicines, kidney or liver disease, blood disorders, nerve disease, or systemic illness. Persistent generalized itch requires evaluation.

\synonyms
Pruritus (Itch)

\medterm{Itching} Itching, or pruritus, is an unpleasant sensation that provokes scratching and may arise from skin disease, allergy, systemic illness, medications, neuropathy, or psychogenic factors.
be acute or chronic and may arise from skin disease, systemic illness, medication, or
neuropathic mechanisms.

\medterm{Itching Anal (Anal Itching)} Alternate index label for Anal Itching.

\medterm{Itching, Anal} Itching around the anus, commonly caused by irritant dermatitis, hemorrhoids, pinworms, infection, inflammatory skin disease, or hygiene-related irritation. Persistent, severe, bleeding, painful, or unexplained symptoms require examination.

\medterm{Itchy Ear} An itching sensation of the external ear or ear canal, which may result from dermatitis, allergy, infection, dryness, or irritation from objects placed in the ear.

\medterm{Iterative Reconstruction} A computer method that repeatedly improves medical images by comparing calculated results with measured scan data. It can reduce image noise and sometimes allow lower radiation exposure, but usually requires more processing time.

\medterm{ITP} Alternate index label; see \textit{Immune thrombocytopenia}.

\medterm{Itraconazole} Itraconazole is a triazole antifungal that inhibits fungal ergosterol synthesis and treats several superficial and systemic mycoses; absorption, extensive drug interactions, liver injury, and negative inotropic effects require attention.

\medterm{IUCD} An intrauterine contraceptive device, usually referring to an intrauterine device placed in the uterus to prevent pregnancy.

\medterm{IV Fluids} Intravenous fluid therapy for volume resuscitation and maintenance. Crystalloids: normal
saline (0.9\% NaCl, increases chloride, risk of hyperchloremic metabolic acidosis),
lactated Ringer (balanced, contains K+, Ca2+, lactate buffer, safer in large volumes),
Plasma-Lyte (balanced, closest to plasma).

\medterm{IV Treatment At Home} Intravenous treatment given outside a hospital, usually through a catheter with instructions for safe infusion, line care, monitoring, and follow-up.

\medterm{Ivabradine} Ivabradine slows the heart rate by acting on the sinus node and is used for selected patients with chronic heart failure or stable angina.

\medterm{IVC Filter} Device placed in inferior vena cava to prevent pulmonary embolism. Indications: DVT/PE
with contraindication to anticoagulation, recurrent PE despite anticoagulation.
Complications: filter thrombosis, migration, perforation, IVC occlusion.

\medterm{Ivermectin} An antiparasitic medicine used for selected worm and mite infections, including strongyloidiasis, onchocerciasis, and scabies, when prescribed appropriately.

\medterm{IVF} In vitro fertilization: the assisted reproductive technology in which oocytes are
retrieved, fertilized by sperm outside the body (in vitro), and the resulting embryos
transferred to the uterus. Indications: tubal factor, male factor (with ICSI),
endometriosis, unexplained, ovulation disorders, genetic testing, and same-sex
couples or individuals using donor sperm or eggs.

\medterm{Ixabepilone} Ixabepilone is a chemotherapy medicine used for some advanced breast cancers; important risks include nerve damage, low blood counts, and allergic reactions.

\medterm{Ixodes} A genus of hard ticks, including the ticks that transmit Lyme disease and several other infections. A bite may be painless; removing the tick promptly reduces, but does not eliminate, infection risk.
\medterm{Internal Carotid Artery} A branch of the common carotid artery that enters the skull and supplies the brain, eyes, and related structures. It contributes to the cerebral arterial circulation.

\medterm{Idiopathic Anaphylaxis} Repeated anaphylaxis, a severe allergic reaction, for which no specific trigger is found after appropriate evaluation. It is still an emergency: use prescribed epinephrine and seek urgent medical care, while specialist assessment looks for hidden triggers or mast-cell disorders.

\medterm{Idiopathic Guttate Hypomelanosis} Harmless, small, flat white spots caused by reduced skin pigment, usually on the shins or forearms of older adults. They are not contagious and usually do not need treatment.

\medterm{Iliotibial Band Friction Syndrome} An overuse injury causing pain on the outside of the knee, and sometimes the hip, when the iliotibial band becomes irritated. It is common in runners and cyclists.

\medterm{Immunoassays} Laboratory tests that use specific binding between antibodies and target substances to detect or measure them. Examples include ELISA, chemiluminescent tests, and some rapid tests; results must be interpreted with the clinical situation.

\medterm{Immunomodulatory} Describing a treatment or substance that changes immune activity. It may stimulate, suppress, or otherwise regulate the immune response, depending on the medicine and condition.

\medterm{Implants} Medical devices, materials, or tissues placed in the body to replace, support, monitor, or deliver treatment. Possible problems include infection, rejection, movement, breakage, or failure.

\medterm{Impotence} An older term usually referring to erectile dysfunction: persistent difficulty getting or keeping an erection firm enough for sexual activity. Erectile dysfunction can have vascular, neurologic, hormonal, medication-related, or psychological causes.

\medterm{In Vitro Fertilization (IVF)} An assisted-reproduction treatment in which eggs are collected, fertilized with sperm in a laboratory, and an embryo is transferred to the uterus. Risks include ovarian hyperstimulation, multiple pregnancy, ectopic pregnancy, and procedure-related complications.

\medterm{Infantile Bradycardia} An unusually slow heart rate in an infant. It can occur normally during sleep, but may also signal low oxygen, infection, medicines, or a heart or conduction problem; symptoms and the infant's age determine the evaluation.

\medterm{Infectious Diseases} Illnesses caused by bacteria, viruses, fungi, parasites, or prions. They may spread through contact, droplets or air, food or water, insects, blood, or sexual contact; prevention and treatment depend on the organism.

\medterm{Infrared Therapy} Treatment using infrared radiation, usually to provide heat or light to tissues. It is used in some physiotherapy and skin treatments, but benefits vary and excessive exposure can cause burns or eye injury.

\medterm{Ingrown Toenails} A condition in which the edge of a toenail grows into nearby skin, most often on the big toe. It causes pain, redness, and swelling; infection, diabetes, or severe symptoms require medical advice.

\medterm{International Classification of Diseases} A World Health Organization system for coding diseases, injuries, and other health conditions. Versions such as ICD-10 and ICD-11 support clinical records, statistics, and billing; a code does not replace a clinical diagnosis.

\medterm{Interrupted Aortic Arch} A rare congenital heart defect in which the aorta has a complete gap between its ascending and descending parts. Newborns can rapidly develop poor circulation and heart failure, so it requires emergency stabilization and surgery.

\medterm{Intrauterine Insemination} A fertility procedure in which prepared sperm is placed directly into the uterus around ovulation. It is used for selected infertility causes and generally requires at least one open fallopian tube; risks include cramps, infection, and multiple pregnancy.

\medterm{Intravesical Therapy} Treatment delivered directly into the bladder through a catheter. It is used for conditions such as some bladder cancers or interstitial cystitis and may cause bladder irritation or, less commonly, effects elsewhere in the body.

\medterm{Iodine Perfusion} The distribution of iodine-containing contrast through tissues or blood vessels, measured during imaging such as computed tomography or angiography. Contrast can rarely cause allergic reactions and may affect kidney or thyroid function.

\medterm{Ion Channel} A protein pore in a cell membrane that allows selected electrically charged particles, such as sodium, potassium, calcium, or chloride, to cross. Ion channels are important for nerve signals, muscle contraction, and secretion.

\medterm{Iridodonesis} Shaking or trembling of the iris when the eye moves, often because the lens is absent or its supporting fibers are weak. An eye examination can identify the cause and related lens or eye problems.

\medterm{Iridoschisis} A rare splitting of the iris tissue into layers, usually in older adults. It may be associated with glaucoma or corneal problems and should be assessed by an eye specialist.

\medterm{Ischemic Optic Neuropathy} Damage to the optic nerve caused by reduced blood flow, often producing sudden, usually painless vision loss. It may be arteritic or non-arteritic; urgent assessment is especially important in older adults or when headache, jaw pain, or scalp tenderness occurs.
